% !TEX program = xelatex
% Copyright 2026 Ma.XianLiang. All rights reserved.
% First released simultaneously on TSA.CN and GitHub, timestamp June 2026

\documentclass[12pt,leqno]{article}
% ========== 中文支持（XeLaTeX + ctex）==========
% 编译方式：XeLaTeX（TeXStudio 中“构建 -> 默认编译器”选 XeLaTeX 即可），LuaLaTeX 亦可
\usepackage[UTF8]{ctex}

% ---------- 字体细化（仅补中文支持，不改动原稿西文观感）----------
% 中文：沿用 ctex 的 mac 字库（宋体正文 / 黑体标题 / 仿宋等宽）
% 仿宋缺少粗体字面，补一个伪粗体，避免 \texttt{中文} 在 \textbf 中静默失去加粗
\setCJKmonofont[AutoFakeBold=4]{STFangsong}
% 首页 DOI 小框原用 Times（ptm）；TU 编码下 ptm 无字体声明文件，改用 Times 系字体
\IfFontExistsTF{Times New Roman}
  {\newfontfamily\ptmbox{Times New Roman}}
  {\newfontfamily\ptmbox{TeXGyreTermes-Regular}}

% ========== PACKAGES ==========
\usepackage{amsmath,amssymb,amsfonts}
\usepackage{mathtools} 
\usepackage{bm}
\usepackage{siunitx}

% ---------- Graphics and Plotting ----------
\usepackage[dvipsnames,table]{xcolor}
\usepackage{graphicx}
\usepackage{tikz,pgfplots}
\pgfplotsset{compat=1.18}

% ---------- Tables and Formatting ----------
\usepackage{array,booktabs}
\usepackage{float}
\usepackage{enumitem}
\usepackage{braket}
\usepackage[font=small]{caption}
\usepackage{url}

% ---------- Left Colume of First Page ------
\usepackage{eso-pic}
\usepackage{xcolor}


% ---------- Page Layout ----------
\setlength{\textheight}{22.0cm}
\setlength{\voffset}{-1.0cm}
\setlength{\textwidth}{16cm}
\setlength{\footskip}{1.4cm}
\setlength{\oddsidemargin}{0cm}
\setlength{\evensidemargin}{0cm}
\topmargin=0pt
\pagestyle{plain}
% ---------- 行距与段间距（本版调整：行距 +32%，段间距 +25%）----------
% 行距：基线拉伸 1.0 -> 1.32，即同一段落内行与行的距离增加 32%
\renewcommand{\baselinestretch}{1.32}
% 段间距：article 默认 \smallskipamount=3pt、\medskipamount=6pt、\bigskipamount=12pt
% 此处仅放大承担“段与段之间”分隔的 medskip / bigskip，各乘 1.25；
% smallskip 用于引出公式的行内小间隔，属段落内部，保持原值不动。
\setlength{\medskipamount}{7.5pt plus 2.5pt minus 2.5pt}
\setlength{\bigskipamount}{15pt plus 5pt minus 5pt}

% ========== PDF COMPRESSION (PdfLaTeX ONLY) ==========
% These commands are only valid in PdfLaTeX, so we check before using them
\ifx\pdfoutput\undefined
% Not PdfLaTeX (XeLaTeX or LuaLaTeX) - skip these commands
\else
% PdfLaTeX - safe to use
\pdfcompresslevel=9
\pdfobjcompresslevel=0
\fi

% ========== Bibliography ==========
\bibliographystyle{unsrt}


\def\a{\mathcal{A}}
\def\b{\mathcal{B}}
\def\c{\mathcal{C}}
\def\d{\mathcal{D}}
\def\e{\mathcal{E}}
\def\f{\mathcal{F}}
\def\g{\mathcal{G}}
\def\h{\mathcal{H}}
\def\i{\mathcal{I}}
\def\j{\mathcal{J}} 
\def\k{\mathcal{K}} 
\def\l{\mathcal{L}}
\def\m{\mathcal{M}}
\def\n{\mathcal{N}}
\def\p{\mathcal{P}}
\def\q{\mathcal{Q}}
\def\r{\mathcal{R}}
\def\s{\mathcal{S}}
\def\t{\mathcal{T}}
\def\u{\mathcal{U}}
\def\v{\mathcal{V}}
\def\w{\mathcal{W}}
\def\x{\mathcal{X}}
\def\y{\mathcal{Y}}
\def\z{\mathcal{Z}}

\def\sp{(s,\psi)}
\def\ss{(\rho,\psi)}
\def\bp{\bar{\psi}}
\def\pc{\chi\psi}
\def\tp{\frac{1}{2\pi i}}
\def\al{\alpha}
\def\ep{\varepsilon}
\def\vt{\vartheta}
\def\la{\lambda}
\def\bla{\bar{\lambda}}
\def\tg{\tilde{g}}
\def\vs{\varsigma}
\def\tz{\tilde{Z}}
\def\vk{\varkappa}
\def\bt{\bar{\theta}}
\def\bk{\bar{\varkappa}}
\def\be{\beta}
\def\tk{\tilde{\kappa}}
\def\hk{\hat{\kappa}}
\def\aa{\mathfrak{a}}
\def\bb{\mathfrak{b}}
\def\cc{\mathfrak{c}}
\def\ff{\mathfrak{f}}
\def\gh{\mathfrak{g}}
\def\xx{\mathfrak{x}}
\def\yy{\mathfrak{y}}

\newcommand{\tinysp}{\mspace{1.5mu}}


% ========== Embedded CSV Data for pgfplots Plotting, Compatible with arXiv Submission ==========

\usepackage{filecontents}

\begin{filecontents*}{2026-growing-ct-fine-model-1381.csv}
	X,Y
	-0.0000100000034500,1.0000000000000000,
	-0.0100134546067800,0.9993108172287610,
	-0.0200238093018500,0.9986218846653470,
	-0.0300410711049400,0.9979332020054410,
	-0.0400652470449900,0.9972447689031750,
	-0.0500963441635800,0.9965565850913760,
	-0.0601343695150200,0.9958686501475920,
	-0.0701793301663300,0.9951809637666520,
	-0.0802312331973400,0.9944935257101380,
	-0.0902900857006300,0.9938063355128980,
	-0.1003558947816700,0.9931193929550830,
	-0.1104286675587900,0.9924326976278250,
	-0.1205084111632200,0.9917462493153290,
	-0.1305951327391600,0.9910600476075540,
	-0.1406888394437800,0.9903740920877050,
	-0.1507895384472600,0.9896883826767870,
	-0.1608972369328200,0.9890029188153060,
	-0.1710119420968300,0.9883177002451410,
	-0.1811336611487500,0.9876327266020520,
	-0.1912624013111500,0.9869479976648600,
	-0.2013981698199000,0.9862635129367670,
	-0.2115409739240300,0.9855792723182100,
	-0.2216908208858900,0.9848952753496390,
	-0.2318477179810600,0.9842115215155820,
	-0.2420116724985300,0.9835280110018670,
	-0.2521826917407300,0.9828447428520080,
	-0.2623607830233600,0.9821617171450130,
	-0.2725459536756300,0.9814789335471560,
	-0.2827382110403000,0.9807963912420450,
	-0.2929375624736100,0.9801140907573310,
	-0.3031440153454100,0.9794320309203790,
	-0.3133575770391000,0.9787502119597160,
	-0.3235782549517200,0.9780686332759930,
	-0.3338060564940200,0.9773872947192020,
	-0.3440409890905100,0.9767061955989990,
	-0.3542830601794300,0.9760253360020740,
	-0.3645322772126800,0.9753447152835670,
	-0.3747886476562300,0.9746643333928620,
	-0.3850521789898200,0.9739841900988080,
	-0.3953228787070300,0.9733042846928250,
	-0.4056007543154500,0.9726246167163850,
	-0.4158858133366200,0.9719451864227630,
	-0.4261780633062300,0.9712659929674250,
	-0.4364775117738900,0.9705870363317990,
	-0.4467841663033200,0.9699083160072070,
	-0.4570980344725100,0.9692298315751930,
	-0.4674191238735300,0.9685515831231610,
	-0.4777474421128300,0.9678735702099730,
	-0.4880829968109500,0.9671957918558990,
	-0.4984257956027600,0.9665182483630900,
	-0.5087758461375700,0.9658409390516930,
	-0.5191331560789500,0.9651638642344890,
	-0.5294977331050500,0.9644870230277400,
	-0.5398695849083200,0.9638104149565880,
	-0.5502487191958300,0.9631340402074770,
	-0.5606351436891900,0.9624578980696040,
	-0.5710288661246400,0.9617819879897720,
	-0.5814298942529700,0.9611063104345390,
	-0.5918382358396100,0.9604308644687190,
	-0.6022538986647900,0.9597556498610280,
	-0.6126768905234700,0.9590806658655490,
	-0.6231072192255000,0.9584059133143940,
	-0.6335448925952900,0.9577313907461950,
	-0.6439899184725700,0.9570570979852140,
	-0.6544423047115700,0.9563830354104420,
	-0.6649020591817200,0.9557092017479680,
	-0.6753691897673800,0.9550355973014050,
	-0.6858437043680000,0.9543622215679810,
	-0.6963256108981600,0.9536890742717030,
	-0.7068149172874700,0.9530161551236630,
	-0.7173116314812600,0.9523434632522820,
	-0.7278157614389400,0.9516709992338400,
	-0.7383273151360500,0.9509987617969330,
	-0.7488463005632900,0.9503267511857110,
	-0.7593727257267600,0.9496549673035230,
	-0.7699065986479200,0.9489834090904410,
	-0.7804479273635700,0.9483120764442020,
	-0.7909967199261000,0.9476409692844550,
	-0.8015529844033800,0.9469700872434300,
	-0.8121167288790100,0.9462994298693640,
	-0.8226879614518500,0.9456289970776670,
	-0.8332666902368100,0.9449587870619170,
	-0.8438529233644900,0.9442888018505260,
	-0.8544466689808300,0.9436190398686290,
	-0.8650479352480700,0.9429495007925420,
	-0.8756567303441000,0.9422801842174900,
	-0.8862730624627400,0.9416110901958410,
	-0.8968969398138600,0.9409422177531320,
	-0.9075283706228600,0.9402735668461860,
	-0.9181673631315800,0.9396051374515730,
	-0.9288139255979600,0.9389369293001330,
	-0.9394680662958900,0.9382689417576380,
	-0.9501297935157100,0.9376011741014270,
	-0.9607991155636700,0.9369336267172590,
	-0.9714760407624800,0.9362662987068570,
	-0.9821605774508600,0.9355991903576290,
	-0.9928527339843600,0.9349323001810720,
	-1.0035525187343800,0.9342656290918130,
	-1.0142599400894400,0.9335991766684750,
	-1.0249750064543000,0.9329329413869750,
	-1.0356977262500500,0.9322669240204530,
	-1.0464281079145500,0.9316011240942000,
	-1.0571661599023800,0.9309355405224530,
	-1.0679118906845200,0.9302701737196790,
	-1.0786653087492400,0.9296050233480080,
	-1.0894264226010500,0.9289400888405510,
	-1.1001952407616900,0.9282753698965240,
	-1.1109717717695900,0.9276108660320070,
	-1.1217560241803200,0.9269465769311940,
	-1.1325480065661500,0.9262825020538740,
	-1.1433477275169800,0.9256186419693460,
	-1.1541551956390600,0.9249549950975620,
	-1.1649704195560800,0.9242915620942330,
	-1.1757934079095200,0.9236283425024050,
	-1.1866241693573100,0.9229653346470770,
	-1.1974627125747800,0.9223025400400910,
	-1.2083090462547100,0.9216399574311850,
	-1.2191631791073400,0.9209775868257800,
	-1.2300251198609500,0.9203154275088600,
	-1.2408948772603200,0.9196534791241570,
	-1.2517724600685500,0.9189917417922090,
	-1.2626578770659000,0.9183302145559450,
	-1.2735511370501500,0.9176688971391230,
	-1.2844522488373700,0.9170077897512690,
	-1.2953612212611100,0.9163468917864730,
	-1.3062780631728000,0.9156862028134260,
	-1.3172027834417200,0.9150257222964240,
	-1.3281353909548400,0.9143654508564560,
	-1.3390758946173800,0.9137053871182710,
	-1.3500243033523900,0.9130455302600020,
	-1.3609806261013800,0.9123858820198290,
	-1.3719448718238100,0.9117264395434010,
	-1.3829170494971100,0.9110672039169820,
	-1.3938971681174100,0.9104081746090690,
	-1.4048852366988900,0.9097493518644110,
	-1.4158812642740800,0.9090907340598190,
	-1.4268852598938500,0.9084323214525360,
	-1.4378972326276700,0.9077741140594290,
	-1.4489171915635700,0.9071161108893980,
	-1.4599451458081600,0.9064583119738110,
	-1.4709811044865500,0.9058007171249560,
	-1.4820250767426100,0.9051433254848770,
	-1.4930770717391800,0.9044861374628220,
	-1.5041370986575800,0.9038291522228210,
	-1.5152051666982600,0.9031723691849120,
	-1.5262812850806300,0.9025157884150430,
	-1.5373654630425600,0.9018594093795290,
	-1.5484577098414000,0.9012032318890130,
	-1.5595580347535300,0.9005472554036020,
	-1.5706664470745400,0.8998914793007140,
	-1.5817829561193800,0.8992359044287410,
	-1.5929075712222600,0.8985805282238550,
	-1.6040403017361700,0.8979253529515680,
	-1.6151811570339100,0.8972703769161450,
	-1.6263301465074900,0.8966155992223970,
	-1.6374872795690500,0.8959610201360290,
	-1.6486525656495000,0.8953066397331860,
	-1.6598260141998900,0.8946524568123190,
	-1.6710076346908500,0.8939984724600500,
	-1.6821974366126700,0.8933446844878330,
	-1.6933954294754400,0.8926910935446160,
	-1.7046016228090600,0.8920376993021810,
	-1.7158160261638900,0.8913845012587110,
	-1.7270386491094100,0.8907314987394270,
	-1.7382695012360300,0.8900786923239320,
	-1.7495085921538100,0.8894260802103070,
	-1.7607559314930800,0.8887736632916610,
	-1.7720115289038300,0.8881214410218080,
	-1.7832753940570400,0.8874694124017920,
	-1.7945475366442800,0.8868175786332450,
	-1.8058279663772300,0.8861659377893030,
	-1.8171166929876500,0.8855144885777030,
	-1.8284137262280500,0.8848632332927790,
	-1.8397190758724600,0.8842121695851720,
	-1.8510327517140200,0.8835612986260630,
	-1.8623547635680000,0.8829106195265910,
	-1.8736851212695300,0.8822601312813110,
	-1.8850238346758700,0.8816098335255320,
	-1.8963709136631700,0.8809597258479780,
	-1.9077263681307300,0.8803098096467880,
	-1.9190902079965800,0.8796600828030390,
	-1.9304624432024500,0.8790105452337120,
	-1.9418430837097800,0.8783611978347970,
	-1.9532321395013700,0.8777120380943370,
	-1.9646296205814700,0.8770630671317490,
	-1.9760355369760500,0.8764142841887480,
	-1.9874498987326600,0.8757656885854940,
	-1.9988727159197400,0.8751172807133410,
	-2.0103039986274800,0.8744690607236400,
	-2.0217437569679700,0.8738210262316990,
	-2.0331920010748800,0.8731731780081950,
	-2.0446487411045500,0.8725255166128770,
	-2.0561139872331100,0.8718780398427440,
	-2.0675877496605300,0.8712307488244900,
	-2.0790700386085700,0.8705836424686100,
	-2.0905608643205500,0.8699367204071540,
	-2.1020602370617200,0.8692899830629910,
	-2.1135681671202500,0.8686434294956340,
	-2.1250846648062700,0.8679970590550660,
	-2.1366097404522000,0.8673508708292730,
	-2.1481434044129400,0.8667048666975950,
	-2.1596856670658800,0.8660590453043860,
	-2.1712365388111200,0.8654134043049560,
	-2.1827960300715500,0.8647679454178740,
	-2.1943641512924800,0.8641226685870740,
	-2.2059409129417900,0.8634775732240440,
	-2.2175263255111400,0.8628326561619680,
	-2.2291203995134500,0.8621879213099220,
	-2.2407231454865400,0.8615433653683080,
	-2.2523345739903700,0.8608989883288350,
	-2.2639546956077100,0.8602547915444770,
	-2.2755835209456200,0.8596107736645690,
	-2.2872210606335400,0.8589669329352150,
	-2.2988673253247000,0.8583232706359180,
	-2.3105223256957100,0.8576797871734360,
	-2.3221860724467500,0.8570364803862950,
	-2.3338585763015600,0.8563933499568530,
	-2.3455398480072800,0.8557503961221350,
	-2.3572298983357400,0.8551076196543360,
	-2.3689287380814200,0.8544650181596380,
	-2.3806363780636100,0.8538225920607950,
	-2.3923528291239200,0.8531803418426680,
	-2.4040781021307100,0.8525382660523930,
	-2.4158122079750400,0.8518963644689780,
	-2.4275551575718300,0.8512546367997840,
	-2.4393069618601800,0.8506130825545380,
	-2.4510676318053100,0.8499717032241700,
	-2.4628371783950200,0.8493304953571370,
	-2.4746156126426100,0.8486894600702840,
	-2.4864029455852300,0.8480485977125390,
	-2.4981991882859200,0.8474079054858080,
	-2.5100043518313600,0.8467673869981300,
	-2.5218184473334900,0.8461270385291690,
	-2.5336414859291700,0.8454868609721330,
	-2.5454734787803100,0.8448468536451250,
	-2.5573144370736000,0.8442070167220510,
	-2.5691643720221900,0.8435673500032610,
	-2.5810232948613100,0.8429278512067090,
	-2.5928912168563500,0.8422885225517980,
	-2.6047681492947700,0.8416493621419120,
	-2.6166541034897900,0.8410103694649630,
	-2.6285490907811600,0.8403715443788640,
	-2.6404531225334600,0.8397328871326420,
	-2.6523662101381900,0.8390943974656690,
	-2.6642883650115700,0.8384560731565380,
	-2.6762195985956900,0.8378179164246300,
	-2.6881599223593400,0.8371799251512270,
	-2.7001093477981100,0.8365421000026670,
	-2.7120678864321900,0.8359044392719800,
	-2.7240355498085200,0.8352669432349170,
	-2.7360123495006900,0.8346296134164380,
	-2.7479982971086800,0.8339924463800360,
	-2.7599934042596100,0.8333554433223910,
	-2.7719976826063700,0.8327186031135920,
	-2.7840111438279700,0.8320819258164820,
	-2.7960337996314600,0.8314454122199670,
	-2.8080656617507400,0.8308090600144650,
	-2.8201067419467000,0.8301728700449700,
	-2.8321570520060900,0.8295368422615030,
	-2.8442166037440800,0.8289009730472470,
	-2.8562854090033200,0.8282652671724590,
	-2.8683634796517400,0.8276297205706590,
	-2.8804508275865500,0.8269943354264650,
	-2.8925474647320100,0.8263591083397780,
	-2.9046534030408300,0.8257240416058410,
	-2.9167686544922300,0.8250891320309910,
	-2.9288932310941100,0.8244543827962100,
	-2.9410271448807300,0.8238197904939670,
	-2.9531704079167200,0.8231853564431710,
	-2.9653230322921800,0.8225510788694910,
	-2.9774850301280700,0.8219169604698580,
	-2.9896564135710900,0.8212829968464100,
	-3.0018371947978300,0.8206491904063230,
	-3.0140273860126500,0.8200155375511220,
	-3.0262269994498000,0.8193820430052940,
	-3.0384360473711700,0.8187487031719950,
	-3.0506545420659600,0.8181155176550850,
	-3.0628824958543700,0.8174824846836370,
	-3.0751199210852300,0.8168496088176210,
	-3.0873668301355100,0.8162168839868400,
	-3.0996232354128000,0.8155843144827620,
	-3.1118891493521500,0.8149518966066520,
	-3.1241645844188800,0.8143196312187600,
	-3.1364495531075000,0.8136875181371140,
	-3.1487440679434100,0.8130555563390780,
	-3.1610481414794800,0.8124237459963930,
	-3.1733617862991400,0.8117920868821350,
	-3.1856850150159200,0.8111605784586840,
	-3.1980178402733500,0.8105292198781030,
	-3.2103602747452700,0.8098980108093440,
	-3.2227123311341100,0.8092669512421430,
	-3.2350740221752600,0.8086360405942460,
	-3.2474453606323200,0.8080052789107990,
	-3.2598263592996800,0.8073746651430830,
	-3.2722170310026800,0.8067441999913400,
	-3.2846173885973000,0.8061138810683740,
	-3.2970274449701400,0.8054837096766390,
	-3.3094472130384800,0.8048536854248610,
	-3.3218767057514700,0.8042238075098150,
	-3.3343159360888500,0.8035940756498090,
	-3.3467649170604600,0.8029644895503730,
	-3.3592236617095700,0.8023350471524540,
	-3.3716921831094000,0.8017057509701510,
	-3.3841704943649600,0.8010765987655020,
	-3.3966586086129500,0.8004475902782840,
	-3.4091565390219800,0.7998187254034240,
	-3.4216642987921500,0.7991900041012890,
	-3.4341819011567200,0.7985614252615680,
	-3.4467093593790000,0.7979329895452060,
	-3.4592466867556300,0.7973046948663950,
	-3.4717938966158800,0.7966765419938670,
	-3.4843510023213800,0.7960485309556140,
	-3.4969180172650200,0.7954206599744230,
	-3.5094949548742000,0.7947929303642890,
	-3.5220818286078300,0.7941653410542550,
	-3.5346786519578600,0.7935378910134740,
	-3.5472854384501100,0.7929105806022570,
	-3.5599022016416100,0.7922834085964380,
	-3.5725289551248900,0.7916563756322540,
	-3.5851657125244100,0.7910294781928610,
	-3.5978124874976200,0.7904027216203460,
	-3.6104692937365400,0.7897761016476510,
	-3.6231361449669500,0.7891496180334440,
	-3.6358130549484600,0.7885232724663720,
	-3.6485000374740200,0.7878970620311510,
	-3.6611971063701100,0.7872709865206730,
	-3.6739042754995800,0.7866450473523140,
	-3.6866215587567400,0.7860192428115490,
	-3.6993489700722300,0.7853935732724980,
	-3.7120865234109800,0.7847680376490700,
	-3.7248342327714600,0.7841426358381130,
	-3.7375921121877700,0.7835173679785160,
	-3.7503601757296500,0.7828922328716300,
	-3.7631384375006600,0.7822672286712090,
	-3.7759269116392700,0.7816423590265420,
	-3.7887256123212000,0.7810176198119570,
	-3.8015345537547000,0.7803930129898310,
	-3.8143537501849500,0.7797685361343820,
	-3.8271832158947100,0.7791441901947700,
	-3.8400229652003700,0.7785199746033910,
	-3.8528730124537500,0.7778958895188130,
	-3.8657333720435800,0.7772719318504880,
	-3.8786040583933300,0.7766481039653320,
	-3.8914850859663200,0.7760244053065820,
	-3.9043764692587000,0.7754008357403780,
	-3.9172782228071500,0.7747773919196800,
	-3.9301903611798700,0.7741540775149090,
	-3.9431128989846800,0.7735308876492000,
	-3.9560458508674200,0.7729078262154380,
	-3.9689892315090900,0.7722848913307440,
	-3.9819430556285800,0.7716620796609580,
	-3.9949073379813500,0.7710393940482970,
	-4.0078820933615800,0.7704168355858530,
	-4.0208673366007600,0.7697943999265060,
	-4.0338630825675600,0.7691720886738190,
	-4.0468693461693400,0.7685499012744850,
	-4.0598861423515800,0.7679278385946640,
	-4.0729134860969500,0.7673058951006850,
	-4.0859513924267400,0.7666840766921870,
	-4.0989998764011100,0.7660623811940260,
	-4.1120589531185200,0.7654408043721960,
	-4.1251286377164400,0.7648193524654080,
	-4.1382089453700300,0.7641980193326320,
	-4.1512998912959700,0.7635768056291260,
	-4.1644014907478100,0.7629557149597150,
	-4.1775137590190900,0.7623347420724800,
	-4.1906367114431100,0.7617138872816510,
	-4.2037703633932900,0.7610931532451220,
	-4.2169147302822100,0.7604725378709070,
	-4.2300698275605200,0.7598520398071860,
	-4.2432356707219000,0.7592316576376150,
	-4.2564122752974400,0.7586113951425260,
	-4.2695996568607500,0.7579912491748750,
	-4.2827978310248900,0.7573712170444930,
	-4.2960068134450000,0.7567513042101030,
	-4.3092266198162700,0.7561315039867780,
	-4.3224572658711600,0.7555118187656640,
	-4.3356987673889800,0.7548922501460970,
	-4.3489511401849100,0.7542727955281310,
	-4.3622144001217900,0.7536534531848860,
	-4.3754885630998700,0.7530342253830590,
	-4.3887736450601900,0.7524151125546880,
	-4.4020696619884200,0.7517961089529520,
	-4.4153766299100700,0.7511772180310110,
	-4.4286945648931800,0.7505584407457010,
	-4.4420234830494400,0.7499397728793300,
	-4.4553634005332000,0.7493212157478550,
	-4.4687143335386700,0.7487027707458270,
	-4.4820762983061700,0.7480844346967990,
	-4.4954493111160800,0.7474662104162710,
	-4.5088333882934500,0.7468480927578080,
	-4.5222285462074600,0.7462300836903460,
	-4.5356348012686500,0.7456121848785050,
	-4.5490521699330400,0.7449943929538440,
	-4.5624806686998400,0.7443767077407340,
	-4.5759203141102100,0.7437591313368970,
	-4.5893711227517000,0.7431416607397950,
	-4.6028331112567800,0.7425242962100110,
	-4.6163062963007400,0.7419070367684530,
	-4.6297906946041100,0.7412898821316310,
	-4.6432863229308400,0.7406728350443080,
	-4.6567931980919400,0.7400558923053280,
	-4.6703113369418700,0.7394390521106280,
	-4.6838407563807200,0.7388223161446660,
	-4.6973814733545700,0.7382056823667620,
	-4.7109335048560700,0.7375891521647960,
	-4.7244968679210200,0.7369727250020520,
	-4.7380715796329800,0.7363564005800410,
	-4.7516576571233400,0.7357401758126570,
	-4.7652551175654700,0.7351240531597390,
	-4.7788639781842400,0.7345080303658020,
	-4.7924842562452600,0.7338921057913630,
	-4.8061159690661200,0.7332762845854710,
	-4.8197591340093400,0.7326605629142910,
	-4.8334137684848200,0.7320449386382940,
	-4.8470798899501500,0.7314294122508760,
	-4.8607575159110700,0.7308139860428490,
	-4.8744466639197400,0.7301986542944610,
	-4.8881473515765300,0.7295834208689910,
	-4.9018595965315800,0.7289682853122470,
	-4.9155834164810600,0.7283532451688700,
	-4.9293188291716700,0.7277383007485060,
	-4.9430658523979600,0.7271234528447680,
	-4.9568245040034300,0.7265086994894050,
	-4.9705948018798300,0.7258940391415340,
	-4.9843767639691400,0.7252794749371180,
	-4.9981704082629500,0.7246650033226490,
	-5.0119757528015600,0.7240506239994900,
	-5.0257928156764900,0.7234363391410790,
	-5.0396216150275200,0.7228221462785330,
	-5.0534621690466900,0.7222080448183420,
	-5.0673149959752400,0.7215940345535760,
	-5.0811786141051700,0.7209801151873210,
	-5.0950545417783800,0.7203662874203860,
	-5.1089422973902900,0.7197525495759080,
	-5.1228418993860200,0.7191389011852330,
	-5.1367533662612300,0.7185253415888850,
	-5.1506767165648800,0.7179118704479220,
	-5.1646119688977600,0.7172984885310990,
	-5.1785591419111000,0.7166851948219500,
	-5.1925182543100600,0.7160719870996680,
	-5.2064893248510000,0.7154588666628020,
	-5.2204723723430500,0.7148458338078340,
	-5.2344674156491100,0.7142328869443440,
	-5.2484744736853100,0.7136200253261590,
	-5.2624935654191500,0.7130072488644540,
	-5.2765247098738100,0.7123945552693420,
	-5.2905679261254800,0.7117819492427050,
	-5.3046232333037400,0.7111694245997480,
	-5.3186906505930700,0.7105569842704180,
	-5.3327701972319000,0.7099446258648350,
	-5.3468618925117600,0.7093323510172310,
	-5.3609657557825700,0.7087201594062050,
	-5.3750818064468100,0.7081080471737450,
	-5.3892100639628800,0.7074960157875090,
	-5.4033505478441800,0.7068840649147180,
	-5.4175032776584500,0.7062721963952300,
	-5.4316682730332900,0.7056604058249060,
	-5.4458455536478000,0.7050486946407020,
	-5.4600351392412000,0.7044370628780220,
	-5.4742370496076500,0.7038255082291630,
	-5.4884513045978200,0.7032140322750740,
	-5.5026779241195500,0.7026026351535990,
	-5.5169169281392800,0.7019913132120190,
	-5.5311683366780600,0.7013800671975390,
	-5.5454321698181400,0.7007688977110630,
	-5.5597084476974700,0.7001578050856850,
	-5.5739971905131300,0.6995467868724670,
	-5.5882984185196400,0.6989358416043390,
	-5.6026121520311100,0.6983249705491890,
	-5.6169384114210300,0.6977141768518480,
	-5.6312772171217800,0.6971034535620630,
	-5.6456285896253100,0.6964928019883720,
	-5.6599925494815700,0.6958822246733360,
	-5.6743691173014400,0.6952717189498780,
	-5.6887583137546900,0.6946612826619100,
	-5.7031601595759100,0.6940509193586620,
	-5.7175746755545800,0.6934406263976720,
	-5.7320018825459500,0.6928304018715160,
	-5.7464418014635500,0.6922202484268340,
	-5.7608944532832300,0.6916101632951180,
	-5.7753598590411400,0.6910001454488440,
	-5.7898380398366200,0.6903901967227100,
	-5.8043290168297800,0.6897803146156940,
	-5.8188328112456200,0.6891705004134100,
	-5.8333494443699100,0.6885607531218190,
	-5.8478789375530500,0.6879510717390420,
	-5.8624213122055200,0.6873414569311700,
	-5.8769765898039500,0.6867319044254270,
	-5.8915447918878000,0.6861224206317080,
	-5.9061259400601800,0.6855129976049170,
	-5.9207200559909000,0.6849036403272120,
	-5.9353271614092900,0.6842943438157090,
	-5.9499472781125800,0.6836851116155080,
	-5.9645804279656900,0.6830759409436450,
	-5.9792266328942900,0.6824668331114800,
	-5.9938859148912900,0.6818577855756840,
	-6.0085582960148000,0.6812487999025130,
	-6.0232437983913700,0.6806398758174770,
	-6.0379424442117800,0.6800310080173560,
	-6.0526542557347700,0.6794222004734110,
	-6.0673792552848100,0.6788134541688890,
	-6.0821174652546100,0.6782047643001760,
	-6.0968689081034300,0.6775961321150570,
	-6.1116336063600500,0.6769875593555870,
	-6.1264115826187100,0.6763790404031840,
	-6.1412028595456700,0.6757705799214360,
	-6.1560074598728400,0.6751621746172310,
	-6.1708254064031700,0.6745538251983190,
	-6.1856567220066600,0.6739455309804910,
	-6.2005014296246700,0.6733372899414580,
	-6.2153595522670500,0.6727291035320730,
	-6.2302311130176400,0.6721209701539310,
	-6.2451161350252400,0.6715128911162030,
	-6.2600146415120000,0.6709048627350700,
	-6.2749266557726800,0.6702968876289160,
	-6.2898522011719700,0.6696889632238200,
	-6.3047913011451300,0.6690810896650990,
	-6.3197439792030400,0.6684732662824490,
	-6.3347102589231800,0.6678654957907650,
	-6.3496901639613900,0.6672577718838000,
	-6.3646837180408800,0.6666500963526510,
	-6.3796909449635300,0.6660424699451550,
	-6.3947118686002200,0.6654348918308350,
	-6.4097465129030300,0.6648273609864680,
	-6.4247949018886400,0.6642198760266940,
	-6.4398570596536000,0.6636124391070400,
	-6.4549330103688800,0.6630050462066860,
	-6.4700227782793800,0.6623977012866240,
	-6.4851263877079000,0.6617903977499210,
	-6.5002438630491300,0.6611831431973370,
	-6.5153752287786900,0.6605759279875340,
	-6.5305205094430100,0.6599687599354400,
	-6.5456797296721400,0.6593616307888880,
	-6.5608529141701800,0.6587545451478090,
	-6.5760400877166100,0.6581475028600110,
	-6.5912412751703000,0.6575405010075520,
	-6.6064565014681800,0.6569335404031960,
	-6.6216857916262800,0.6563266184102150,
	-6.6369291707395000,0.6557197380080110,
	-6.6521866639829800,0.6551128964823600,
	-6.6674582966086100,0.6545060916429360,
	-6.6827440939488300,0.6538993271232290,
	-6.6980440814189400,0.6532925998927500,
	-6.7133582845139400,0.6526859099477640,
	-6.7286867288073900,0.6520792544811470,
	-6.7440294399550500,0.6514726354101080,
	-6.7593864436966500,0.6508660545908890,
	-6.7747577658533200,0.6502595081366570,
	-6.7901434323267000,0.6496529948573160,
	-6.8055434691024900,0.6490465165676830,
	-6.8209579022520700,0.6484400703584340,
	-6.8363867579234800,0.6478336581973340,
	-6.8518300623536400,0.6472272771632830,
	-6.8672878418662200,0.6466209286491820,
	-6.8827601228630300,0.6460146125095320,
	-6.8982469318360600,0.6454083258687210,
	-6.9137482953605300,0.6448020698522740,
	-6.9292642400985600,0.6441958417143020,
	-6.9447947927988400,0.6435896451481810,
	-6.9603399802937300,0.6429834758596690,
	-6.9758998295046600,0.6423773365336940,
	-6.9914743674400800,0.6417712233658180,
	-7.0070636211979700,0.6411651374171230,
	-7.0226676179616500,0.6405590773445480,
	-7.0382863850063300,0.6399530414356490,
	-7.0539199496925900,0.6393470339067640,
	-7.0695683394715100,0.6387410512421800,
	-7.0852315818859000,0.6381350904509390,
	-7.1009097045660100,0.6375291569166300,
	-7.1166027352334800,0.6369232436549970,
	-7.1323107017034000,0.6363173531391700,
	-7.1480336318792500,0.6357114845933670,
	-7.1637715537599600,0.6351056388159630,
	-7.1795244954312000,0.6344998149997280,
	-7.1952924850739500,0.6338940096200800,
	-7.2110755509656800,0.6332882261411760,
	-7.2268737214720900,0.6326824597746760,
	-7.2426870250569100,0.6320767132219100,
	-7.2585154902753700,0.6314709836942060,
	-7.2743591457801600,0.6308652733027150,
	-7.2902180203164100,0.6302595795074700,
	-7.3060921427253300,0.6296539015881470,
	-7.3219815419474300,0.6290482405572390,
	-7.3378862470197000,0.6284425940915210,
	-7.3538062870724600,0.6278369624204450,
	-7.3697416913360200,0.6272313455853680,
	-7.3856924891396700,0.6266257421936910,
	-7.4016587099084200,0.6260201543799800,
	-7.4176403831691900,0.6254145764865930,
	-7.4336375385419900,0.6248090097268780,
	-7.4496502057555300,0.6242034577539630,
	-7.4656784146318700,0.6235979138368970,
	-7.4817221950984500,0.6229923805463020,
	-7.4977815771794100,0.6223868583166850,
	-7.5138565910054800,0.6217813458554980,
	-7.5299472668064200,0.6211758392360190,
	-7.5460536349162100,0.6205703424558870,
	-7.5621757257701500,0.6199648515669590,
	-7.5783135699079800,0.6193593693738010,
	-7.5944671979759900,0.6187538926025570,
	-7.6106366407235600,0.6181484221487870,
	-7.6268219290027000,0.6175429591353200,
	-7.6430230937729400,0.6169374953707050,
	-7.6592401661006300,0.6163320385743910,
	-7.6754731771589500,0.6157265854334060,
	-7.6917221582271300,0.6151211353386020,
	-7.7079871406924800,0.6145156851959060,
	-7.7242681560491100,0.6139102396764480,
	-7.7405652359029200,0.6133047932072080,
	-7.7568784119659700,0.6126993456296850,
	-7.7732077160621700,0.6120938998547030,
	-7.7895531801265300,0.6114884550742720,
	-7.8059148362002600,0.6108830065747340,
	-7.8222927164415000,0.6102775580182190,
	-7.8386868531163400,0.6096721062323810,
	-7.8550972786033600,0.6090666480469240,
	-7.8715240253962400,0.6084611890908700,
	-7.8879671261029200,0.6078557274466840,
	-7.9044266134416500,0.6072502591138130,
	-7.9209025202474900,0.6066447840738480,
	-7.9373948794722100,0.6060393031278920,
	-7.9539037241816100,0.6054338167627270,
	-7.9704290875573300,0.6048283235977280,
	-7.9869710028970300,0.6042228204319240,
	-8.0035295036186700,0.6036173106918160,
	-8.0201046232561800,0.6030117926666970,
	-8.0366963954635700,0.6024062622759450,
	-8.0533048540112600,0.6018007220387330,
	-8.0699300327917000,0.6011951712162680,
	-8.0865719658178700,0.6005896094495100,
	-8.1032306872235200,0.5999840336669430,
	-8.1199062312638900,0.5993784474121580,
	-8.1365986323154100,0.5987728471652300,
	-8.1533079248755300,0.5981672299268110,
	-8.1700341435719200,0.5975616006257260,
	-8.1867773231468200,0.5969559558587730,
	-8.2035374984735600,0.5963502935967180,
	-8.2203147045493700,0.5957446171881550,
	-8.2371089764989600,0.5951389214046710,
	-8.2539203495660100,0.5945332095887720,
	-8.2707488591285700,0.5939274802752220,
	-8.2875945406939300,0.5933217300175180,
	-8.3044574298903300,0.5927159595379010,
	-8.3213375624812500,0.5921101703135600,
	-8.3382349743587000,0.5915043586265950,
	-8.3551497015415000,0.5908985255910790,
	-8.3720817801830100,0.5902926713646210,
	-8.3890312465668800,0.5896867922606930,
	-8.4059981371141800,0.5890808903927200,
	-8.4229824883621100,0.5884749667489490,
	-8.4399843369974100,0.5878690158786610,
	-8.4570037198400700,0.5872630392179870,
	-8.4740406738403100,0.5866570375957380,
	-8.4910952360861900,0.5860510077145170,
	-8.5081674438010500,0.5854449527279740,
	-8.5252573343450900,0.5848388694898950,
	-8.5423649452196800,0.5842327551977630,
	-8.5594903140629200,0.5836266131966920,
	-8.5766334786442200,0.5830204415484600,
	-8.5937944768859700,0.5824142407512050,
	-8.6109733468419600,0.5818080043474830,
	-8.6281701267114900,0.5812017386752360,
	-8.6453848548317400,0.5805954408998310,
	-8.6626175696938600,0.5799891098808420,
	-8.6798683099148700,0.5793827423948780,
	-8.6971371142742900,0.5787763420186480,
	-8.7144240216714400,0.5781699083790020,
	-8.7317290711891900,0.5775634359506800,
	-8.7490523020216000,0.5769569281873500,
	-8.7663937535285900,0.5763503828924000,
	-8.7837534652131200,0.5757437976759200,
	-8.8011314767249300,0.5751371762137100,
	-8.8185278278652300,0.5745305139590420,
	-8.8359425585898700,0.5739238136261340,
	-8.8533757089995100,0.5733170704305490,
	-8.8708273193491400,0.5727102869959480,
	-8.8882974300468300,0.5721034626080840,
	-8.9057860816529400,0.5714965944535330,
	-8.9232933148875000,0.5708896825955600,
	-8.9408191706156600,0.5702827269957420,
	-8.9583636898663100,0.5696757277224720,
	-8.9759269138245500,0.5690686816709720,
	-8.9935088838287700,0.5684615906834170,
	-9.0111096413805200,0.5678544500067110,
	-9.0287292281370100,0.5672472650264960,
	-9.0463676859193500,0.5666400299126980,
	-9.0640250567072700,0.5660327463110700,
	-9.0817013826408400,0.5654254147149300,
	-9.0993967060274200,0.5648180304177560,
	-9.1171110693370400,0.5642105962899440,
	-9.1348445151995800,0.5636031111626490,
	-9.1525970864137900,0.5629955718298400,
	-9.1703688259434200,0.5623879806879540,
	-9.1881597769219300,0.5617803345055480,
	-9.2059699826470200,0.5611726361746470,
	-9.2237994865872900,0.5605648789292550,
	-9.2416483323815500,0.5599570704117990,
	-9.2595165638391800,0.5593492017497190,
	-9.2774042249418800,0.5587412759615870,
	-9.2953113598418600,0.5581332929961460,
	-9.3132380128667400,0.5575252511216550,
	-9.3311842285168500,0.5569171495481960,
	-9.3491500514711900,0.5563089876837570,
	-9.3671355265796300,0.5557007642851920,
	-9.3851406988763600,0.5550924789601800,
	-9.4031656135707600,0.5544841322282190,
	-9.4212103160522500,0.5538757215006410,
	-9.4392748518837300,0.5532672459141130,
	-9.4573592668198000,0.5526587077849660,
	-9.4754636067940600,0.5520501030049040,
	-9.4935879179195200,0.5514414305635110,
	-9.5117322464959600,0.5508326950237180,
	-9.5298966390102200,0.5502238890276130,
	-9.5480811421342900,0.5496150155452910,
	-9.5662858027248500,0.5490060729828710,
	-9.5845106678324700,0.5483970605028650,
	-9.6027557846910600,0.5477879763104780,
	-9.6210212007295100,0.5471788210687320,
	-9.6393069635676600,0.5465695934788230,
	-9.6576131210117300,0.5459602945297220,
	-9.6759397210749900,0.5453509216839650,
	-9.6942868119525400,0.5447414727911940,
	-9.7126544420424000,0.5441319501499540,
	-9.7310426599387600,0.5435223504597060,
	-9.7494515144296200,0.5429126754200860,
	-9.7678810545107500,0.5423029226405050,
	-9.7863313293668300,0.5416930911922080,
	-9.8048023883938200,0.5410831800739500,
	-9.8232942811857700,0.5404731894158950,
	-9.8418070575413900,0.5398631196092250,
	-9.8603407674641800,0.5392529672707600,
	-9.8788954611648900,0.5386427312608100,
	-9.8974711890606200,0.5380324140726780,
	-9.9160680017764600,0.5374220144722330,
	-9.9346859501485800,0.5368115262039760,
	-9.9533250852187500,0.5362009582805210,
	-9.9719854582467800,0.5355903013390670,
	-9.9906671207061200,0.5349795568896250,
	-10.0093701242882000,0.5343687250816960,
	-10.0280945208835000,0.5337578045688110,
	-10.0468403626171000,0.5331467974574120,
	-10.0656077018256000,0.5325356980314240,
	-10.0843965910643000,0.5319245072748740,
	-10.1032070831088000,0.5313132237281790,
	-10.1220392309583000,0.5307018516439080,
	-10.1408930878394000,0.5300903828172360,
	-10.1597687071961000,0.5294788214431990,
	-10.1786661427009000,0.5288671652185430,
	-10.1975854482525000,0.5282554151732620,
	-10.2165266779815000,0.5276435648811200,
	-10.2354898862508000,0.5270316179300570,
	-10.2544751276445000,0.5264195742794530,
	-10.2734824569879000,0.5258074294480950,
	-10.2925119293397000,0.5251951858484750,
	-10.3115635999865000,0.5245828440891920,
	-10.3306375244607000,0.5239703973320370,
	-10.3497337585284000,0.5233578490164430,
	-10.3688523581979000,0.5227451995801470,
	-10.3879933797130000,0.5221324428369030,
	-10.4071568795654000,0.5215195825534860,
	-10.4263429144908000,0.5209066182921770,
	-10.4455515414668000,0.5202935443498280,
	-10.4647828177177000,0.5196803667036840,
	-10.4840368007195000,0.5190670761455550,
	-10.5033135481975000,0.5184536807637010,
	-10.5226131181252000,0.5178401727211610,
	-10.5419355687297000,0.5172265543073410,
	-10.5612809584943000,0.5166128252642780,
	-10.5806493461572000,0.5159989815768070,
	-10.6000407907123000,0.5153850272625800,
	-10.6194553514178000,0.5147709555681560,
	-10.6388930877843000,0.5141567706618270,
	-10.6583540595882000,0.5135424684966580,
	-10.6778383268738000,0.5129280513933790,
	-10.6973459499453000,0.5123135127455850,
	-10.7168769893756000,0.5116988578400060,
	-10.7364315060056000,0.5110840850904830,
	-10.7560095609422000,0.5104691872881950,
	-10.7756112155739000,0.5098541720288270,
	-10.7952365315542000,0.5092390333031230,
	-10.8148855708148000,0.5086237707801960,
	-10.8345583955636000,0.5080083862363120,
	-10.8542550682907000,0.5073928749962520,
	-10.8739756517586000,0.5067772359013480,
	-10.8937202090186000,0.5061614723313860,
	-10.9134888034082000,0.5055455821622380,
	-10.9332814985382000,0.5049295602671330,
	-10.9530983583198000,0.5043134127471160,
	-10.9729394469433000,0.5036971293516430,
	-10.9928048288961000,0.5030807202169310,
	-11.0126945689596000,0.5024641760664770,
	-11.0326087322012000,0.5018474991510820,
	-11.0525473839940000,0.5012306881247400,
	-11.0725105900006000,0.5006137397033980,
	-11.0924984161933000,0.4999966607389560,
	-11.1125109288383000,0.4993794398281250,
	-11.1325481945111000,0.4987620813958280,
	-11.1526102800878000,0.4981445861546190,
	-11.1726972527579000,0.4975269506123540,
	-11.1928091800186000,0.4969091725272040,
	-11.2129461296732000,0.4962912536414170,
	-11.2331081698513000,0.4956731934485550,
	-11.2532953689853000,0.4950549867420870,
	-11.2735077958296000,0.4944366392880020,
	-11.2937455194618000,0.4938181442448940,
	-11.3140086092780000,0.4931995025656380,
	-11.3342971349971000,0.4925807143680820,
	-11.3546111666666000,0.4919617752340230,
	-11.3749507746616000,0.4913426869413390,
	-11.3953160296826000,0.4907234506433430,
	-11.4157070027741000,0.4901040627117330,
	-11.4361237652999000,0.4894845223318980,
	-11.4565663889714000,0.4888648258439520,
	-11.4770349458316000,0.4882449769801980,
	-11.4975295082731000,0.4876249727143540,
	-11.5180501490189000,0.4870048098931580,
	-11.5385969411512000,0.4863844927563160,
	-11.5591699580907000,0.4857640151131930,
	-11.5797692736111000,0.4851433790634320,
	-11.6003949618429000,0.4845225812176760,
	-11.6210470972612000,0.4839016224615590,
	-11.6417257547063000,0.4832805012209820,
	-11.6624310093743000,0.4826592188551690,
	-11.6831629368229000,0.4820377685218540,
	-11.7039216129746000,0.4814161563147650,
	-11.7247071141175000,0.4807943729601680,
	-11.7455195169121000,0.4801724256805450,
	-11.7663588983832000,0.4795503093554070,
	-11.7872253359381000,0.4789280196548860,
	-11.8081189073547000,0.4783055627317970,
	-11.8290396907908000,0.4776829322200380,
	-11.8499877647834000,0.4770601310136370,
	-11.8709632082608000,0.4764371531156170,
	-11.8919661005328000,0.4758140021448920,
	-11.9129965212941000,0.4751906729304260,
	-11.9340545506390000,0.4745671680030080,
	-11.9551402690524000,0.4739434861037460,
	-11.9762537574167000,0.4733196218793200,
	-11.9973950970129000,0.4726955788076990,
	-12.0185643695283000,0.4720713556248310,
	-12.0397616570537000,0.4714469479008360,
	-12.0609870420847000,0.4708223569828020,
	-12.0822406075334000,0.4701975815945690,
	-12.1035224367220000,0.4695726186601860,
	-12.1248326133844000,0.4689474716211520,
	-12.1461712216891000,0.4683221326915940,
	-12.1675383462093000,0.4676966072809420,
	-12.1889340719591000,0.4670708929021540,
	-12.2103584843692000,0.4664449850574170,
	-12.2318116693049000,0.4658188845552030,
	-12.2532937130652000,0.4651925904830210,
	-12.2748047023808000,0.4645661004477140,
	-12.2963447244444000,0.4639394151635140,
	-12.3179138668627000,0.4633125339181670,
	-12.3395122177075000,0.4626854537276150,
	-12.3611398654929000,0.4620581745631920,
	-12.3827968991858000,0.4614306934770770,
	-12.4044834082144000,0.4608030108719790,
	-12.4261994824582000,0.4601751275136710,
	-12.4479452122707000,0.4595470380951920,
	-12.4697206884543000,0.4589187430976040,
	-12.4915260022938000,0.4582902424018870,
	-12.5133612455432000,0.4576615336431160,
	-12.5352265104292000,0.4570326165361810,
	-12.5571218896566000,0.4564034889885020,
	-12.5790474764177000,0.4557741505585460,
	-12.6010033643880000,0.4551445991216510,
	-12.6229896477426000,0.4545148348253940,
	-12.6450064211137000,0.4538848558700650,
	-12.6670537796785000,0.4532546617602760,
	-12.6891318190812000,0.4526242477863980,
	-12.7112406354742000,0.4519936187504800,
	-12.7333803255244000,0.4513627686504570,
	-12.7555509864157000,0.4507316972040420,
	-12.7777527158328000,0.4501004021666760,
	-12.7999856119896000,0.4494688856307810,
	-12.8222497736143000,0.4488371420492700,
	-12.8445452999661000,0.4482051752832720,
	-12.8668722908251000,0.4475729812104390,
	-12.8892308465223000,0.4469405581486670,
	-12.9116210679070000,0.4463079044286190,
	-12.9340430563852000,0.4456750215668100,
	-12.9564969139026000,0.4450419056727860,
	-12.9789827429498000,0.4444085560655550,
	-13.0015006465750000,0.4437749715439620,
	-13.0240507283941000,0.4431411519327430,
	-13.0466330925696000,0.4425070930986980,
	-13.0692478438460000,0.4418727948672650,
	-13.0918950875215000,0.4412382586036160,
	-13.1145749294797000,0.4406034831334140,
	-13.1372874761826000,0.4399684633304600,
	-13.1600328346702000,0.4393331971577100,
	-13.1828111125808000,0.4386976887829900,
	-13.2056224181284000,0.4380619325551870,
	-13.2284668601473000,0.4374259266941490,
	-13.2513445480446000,0.4367896734420790,
	-13.2742555918577000,0.4361531693981570,
	-13.2972001022227000,0.4355164144979530,
	-13.3201781904003000,0.4348794035924810,
	-13.3431899682535000,0.4342421426543970,
	-13.3662355482869000,0.4336046207091890,
	-13.3893150436240000,0.4329668434355880,
	-13.4124285680354000,0.4323288069738620,
	-13.4355762359161000,0.4316905105542290,
	-13.4587581623141000,0.4310519519540150,
	-13.4819744629286000,0.4304131296553540,
	-13.5052252541024000,0.4297740451160420,
	-13.5285106528439000,0.4291346931024880,
	-13.5518307768286000,0.4284950751702800,
	-13.5751857444023000,0.4278551870983740,
	-13.5985756745714000,0.4272150303518390,
	-13.6220006870513000,0.4265745996803580,
	-13.6454609022139000,0.4259338983970140,
	-13.6689564411312000,0.4252929205684010,
	-13.6924874255771000,0.4246516684192720,
	-13.7160539780212000,0.4240101370803690,
	-13.7396562216488000,0.4233683275528260,
	-13.7632942803501000,0.4227262387098870,
	-13.7869682787382000,0.4220838671346160,
	-13.8106783421475000,0.4214412155626120,
	-13.8344245966476000,0.4207982768855690,
	-13.8582071690383000,0.4201550508512660,
	-13.8820261868663000,0.4195115379020520,
	-13.9058817784222000,0.4188677348894540,
	-13.9297740727580000,0.4182236415993570,
	-13.9537031996775000,0.4175792554257210,
	-13.9776692897565000,0.4169345776040850,
	-14.0016724743382000,0.4162896019137730,
	-14.0257128855486000,0.4156443298199650,
	-14.0497906562995000,0.4149987607779050,
	-14.0739059202977000,0.4143528899659300,
	-14.0980588120370000,0.4137067167450210,
	-14.1222494668272000,0.4130602433545610,
	-14.1464780207811000,0.4124134623803100,
	-14.1707446108367000,0.4117663763169170,
	-14.1950493747517000,0.4111189815439590,
	-14.2193924511209000,0.4104712763411190,
	-14.2437739794021000,0.4098232631745290,
	-14.2681940998039000,0.4091749337530560,
	-14.2926529534950000,0.4085262924087870,
	-14.3171506824592000,0.4078773316896280,
	-14.3416874295537000,0.4072280556168470,
	-14.3662633385002000,0.4065784618593560,
	-14.3908785539140000,0.4059285432197170,
	-14.4155332212924000,0.4052783049635850,
	-14.4402274870350000,0.4046277412380860,
	-14.4649614984410000,0.4039768514040920,
	-14.4897354037210000,0.4033256317659140,
	-14.5145493520136000,0.4026740865941520,
	-14.5394034933724000,0.4020222061514190,
	-14.5642979787934000,0.4013699961657670,
	-14.5892329602200000,0.4007174497413970,
	-14.6142085905349000,0.4000645689448420,
	-14.6392250235910000,0.3994113466250520,
	-14.6642824142088000,0.3987577865318580,
	-14.6893809181796000,0.3981038848900510,
	-14.7145206922817000,0.3974496426042900,
	-14.7397018942911000,0.3967950515490540,
	-14.7649246829852000,0.3961401148080950,
	-14.7901892181463000,0.3954848290314380,
	-14.8154956605838000,0.3948291937440910,
	-14.8408441721315000,0.3941732052726800,
	-14.8662349156605000,0.3935168615733650,
	-14.8916680550837000,0.3928601625959890,
	-14.9171437553896000,0.3922031081225850,
	-14.9426621826109000,0.3915456913106060,
	-14.9682235038718000,0.3908879153537470,
	-14.9938278873659000,0.3902297733059400,
	-15.0194755023930000,0.3895712680102250,
	-15.0451665193541000,0.3889123972142000,
	-15.0709011097604000,0.3882531543670300,
	-15.0966794462484000,0.3875935429860920,
	-15.1225017025919000,0.3869335569641420,
	-15.1483680537141000,0.3862731981521820,
	-15.1742786756787000,0.3856124609690720,
	-15.2002337457238000,0.3849513448916410,
	-15.2262334422655000,0.3842898494057050,
	-15.2522779449030000,0.3836279715878940,
	-15.2783674344345000,0.3829657090941680,
	-15.3045020928645000,0.3823030607987090,
	-15.3306821034190000,0.3816400237037480,
	-15.3569076505504000,0.3809765955520080,
	-15.3831789199689000,0.3803127776465910,
	-15.4094960986222000,0.3796485610200980,
	-15.4358593747227000,0.3789839521200810,
	-15.4622689377780000,0.3783189421566210,
	-15.4887249785672000,0.3776535315080560,
	-15.5152276891732000,0.3769877195328410,
	-15.5417772630039000,0.3763215030106820,
	-15.5683738947875000,0.3756548775461990,
	-15.5950177805826000,0.3749878461710610,
	-15.6217091178060000,0.3743204010913340,
	-15.6484481052393000,0.3736525433721010,
	-15.6752349430366000,0.3729842701818130,
	-15.7020698327463000,0.3723155829917620,
	-15.7289529773344000,0.3716464739674350,
	-15.7558845811338000,0.3709769417949460,
	-15.7828648499531000,0.3703069865933530,
	-15.8098939910513000,0.3696366071757750,
	-15.8369722131114000,0.3689657987527560,
	-15.8640997263160000,0.3682945578208050,
	-15.8912767423186000,0.3676228836992280,
	-15.9185034742866000,0.3669507775120300,
	-15.9457801368842000,0.3662782333143650,
	-15.9731069463183000,0.3656052483732980,
	-16.0004841203436000,0.3649318225063900,
	-16.0279118782701000,0.3642579512519170,
	-16.0553904409877000,0.3635836362621590,
	-16.0829200309733000,0.3629088713750650,
	-16.1105008723167000,0.3622336555535570,
	-16.1381331907303000,0.3615579843952270,
	-16.1658172135620000,0.3608818599104900,
	-16.1935531698247000,0.3602052762106030,
	-16.2213412901928000,0.3595282335078790,
	-16.2491818070375000,0.3588507266253730,
	-16.2770749544350000,0.3581727535173370,
	-16.3050209681844000,0.3574943149255910,
	-16.3330200858300000,0.3568154053673340,
	-16.3610725466672000,0.3561360208089320,
	-16.3891785917802000,0.3554561643537810,
	-16.4173384640369000,0.3547758276074010,
	-16.4455524081238000,0.3540950134349650,
	-16.4738206705622000,0.3534137160255440,
	-16.5021434997176000,0.3527319328635780,
	-16.5305211458259000,0.3520496622175370,
	-16.5589538610308000,0.3513669002704710,
	-16.5874418993613000,0.3506836468119540,
	-16.6159855167907000,0.3499998967936820,
	-16.6445849712358000,0.3493156492426660,
	-16.6732405225767000,0.3486309011386080,
	-16.7019524326930000,0.3479456518696850,
	-16.7307209654810000,0.3472598947977300,
	-16.7595463868613000,0.3465736271499500,
	-16.7884289648102000,0.3458868507098410,
	-16.8173689693702000,0.3451995588589460,
	-16.8463666727106000,0.3445117511745700,
	-16.8754223490952000,0.3438234240485750,
	-16.9045362749438000,0.3431345742397010,
	-16.9337087288304000,0.3424452015270310,
	-16.9629399915380000,0.3417552971003380,
	-16.9922303460432000,0.3410648647429340,
	-17.0215800775770000,0.3403738984115670,
	-17.0509894736327000,0.3396823973833150,
	-17.0804588239705000,0.3389903556260040,
	-17.1099884206825000,0.3382977712899380,
	-17.1395785581944000,0.3376046406705420,
	-17.1692295332893000,0.3369109655904710,
	-17.1989416451425000,0.3362167380213790,
	-17.2287151953490000,0.3355219572146900,
	-17.2585504879432000,0.3348266189663590,
	-17.2884478294324000,0.3341307205570320,
	-17.3184075288075000,0.3334342609487670,
	-17.3484298976034000,0.3327372348545570,
	-17.3785152499030000,0.3320396387003280,
	-17.4086639023616000,0.3313414732235200,
	-17.4388761742641000,0.3306427311914170,
	-17.4691523875247000,0.3299434110085390,
	-17.4994928667352000,0.3292435082582470,
	-17.5298979391949000,0.3285430234544830,
	-17.5603679349150000,0.3278419509200450,
	-17.5909031866972000,0.3271402860268640,
	-17.6215040301277000,0.3264380278596850,
	-17.6521708036221000,0.3257351706617670,
	-17.6829038484623000,0.3250317167559000,
	-17.7137035088179000,0.3243276565510500,
	-17.7445701317926000,0.3236229892024720,
	-17.7755040674606000,0.3229177104624190,
	-17.8065056688841000,0.3222118186092780,
	-17.8375752921610000,0.3215053117018890,
	-17.8687132964631000,0.3207981820054280,
	-17.8999200440676000,0.3200904276956420,
	-17.9311959003937000,0.3193820464494880,
	-17.9625412340437000,0.3186730335680680,
	-17.9939564168301000,0.3179633863429570,
	-18.0254418238294000,0.3172531009813920,
	-18.0569978334178000,0.3165421759434560,
	-18.0886248273117000,0.3158306042323880,
	-18.1203231905876000,0.3151183834502200,
	-18.1520933117641000,0.3144055108254400,
	-18.1839355828033000,0.3136919812727680,
	-18.2158503991742000,0.3129777910966960,
	-18.2478381598976000,0.3122629396654170,
	-18.2798992675734000,0.3115474178889710,
	-18.3120341284534000,0.3108312282593240,
	-18.3442431524542000,0.3101143626729310,
	-18.3765267532175000,0.3093968187002400,
	-18.4088853481735000,0.3086785894357910,
	-18.4413193585457000,0.3079596775767830,
	-18.4738292094533000,0.3072400731787340,
	-18.5064153299129000,0.3065197736201190,
	-18.5390781529140000,0.3057987788324770,
	-18.5718181154617000,0.3050770812439280,
	-18.6046356586416000,0.3043546747606390,
	-18.6375312276363000,0.3036315593240250,
	-18.6705052718308000,0.3029077286600120,
	-18.7035582448075000,0.3021831805265690,
	-18.7366906044568000,0.3014579081889960,
	-18.7699028129912000,0.3007319098182180,
	-18.8031953370351000,0.3000051819959430,
	-18.8365686476374000,0.2992777155577220,
	-18.8700232203858000,0.2985495122748690,
	-18.9035595354213000,0.2978205624303210,
	-18.9371780775280000,0.2970908663346600,
	-18.9708793361709000,0.2963604142536440,
	-19.0046638055883000,0.2956292090437150,
	-19.0385319848151000,0.2948972387545110,
	-19.0724843777931000,0.2941645043789890,
	-19.1065214933886000,0.2934309954579170,
	-19.1406438455025000,0.2926967157206210,
	-19.1748519531137000,0.2919616535487200,
	-19.2091463403653000,0.2912258065027540,
	-19.2435275366078000,0.2904891705738360,
	-19.2779960765008000,0.2897517408865990,
	-19.3125525000735000,0.2890135117128070,
	-19.3471973527962000,0.2882744793402580,
	-19.3819311856544000,0.2875346389509020,
	-19.4167545552425000,0.2867939846184080,
	-19.4516680238215000,0.2860525117149890,
	-19.4866721594042000,0.2853102149669890,
	-19.5217675358508000,0.2845670919876240,
	-19.5569547329385000,0.2838231345019170,
	-19.5922343364290000,0.2830783371256110,
	-19.6276069382019000,0.2823326966945630,
	-19.6630731362853000,0.2815862081030630,
	-19.6986335349909000,0.2808388649205950,
	-19.7342887449900000,0.2800906632846960,
	-19.7700393833758000,0.2793415957076940,
	-19.8058860738138000,0.2785916584266770,
	-19.8418294465793000,0.2778408444953800,
	-19.8778701386995000,0.2770891489582300,
	-19.9140087940349000,0.2763365694012310,
	-19.9502460633678000,0.2755830970127050,
	-19.9865826045482000,0.2748287264869910,
	-20.0230190825254000,0.2740734536075640,
	-20.0595561695369000,0.2733172686496210,
	-20.0961945451632000,0.2725601710096990,
	-20.1329348964532000,0.2718021513364140,
	-20.1697779180405000,0.2710432052705260,
	-20.2067243122565000,0.2702833271304930,
	-20.2437747892662000,0.2695225108738990,
	-20.2809300671468000,0.2687607484407880,
	-20.3181908720541000,0.2679980370335040,
	-20.3555579383077000,0.2672343669934680,
	-20.3930320085644000,0.2664697336099010,
	-20.4306138339007000,0.2657041326563250,
	-20.4683041739704000,0.2649375523883180,
	-20.5061037971420000,0.2641699927700960,
	-20.5440134806263000,0.2634014432328520,
	-20.5820340106265000,0.2626318986276820,
	-20.6201661824622000,0.2618613523783280,
	-20.6584108007334000,0.2610897962000300,
	-20.6967686794763000,0.2603172266462440,
	-20.7352406423013000,0.2595436348359920,
	-20.7738275225452000,0.2587690125188980,
	-20.8125301634365000,0.2579933553209910,
	-20.8513494182677000,0.2572166551649510,
	-20.8902861505509000,0.2564389040758750,
	-20.9293412341755000,0.2556600971935260,
	-20.9685155536039000,0.2548802253948300,
	-21.0078100040172000,0.2540992827598540,
	-21.0472254915334000,0.2533172609550140,
	-21.0867629333682000,0.2525341509529680,
	-21.1264232580301000,0.2517499472618080,
	-21.1662074054801000,0.2509646430950410,
	-21.2061163273933000,0.2501782305095570,
	-21.2461509873037000,0.2493906987613240,
	-21.2863123608289000,0.2486020433545260,
	-21.3266014358656000,0.2478122554146810,
	-21.3670192128268000,0.2470213245329810,
	-21.4075667048554000,0.2462292465944590,
	-21.4482449380277000,0.2454360109302500,
	-21.4890549516028000,0.2446416104972340,
	-21.5299977982338000,0.2438460361694900,
	-21.5710745442339000,0.2430492794365920,
	-21.6122862698056000,0.2422513322405250,
	-21.6536340693015000,0.2414521849669760,
	-21.6951190514428000,0.2406518295768430,
	-21.7367423396380000,0.2398502562284030,
	-21.7785050722013000,0.2390474586641570,
	-21.8204084026518000,0.2382434258651450,
	-21.8624535000138000,0.2374381474484910,
	-21.9046415490428000,0.2366316178069950,
	-21.9469737505729000,0.2358238242850930,
	-21.9894513217964000,0.2350147582280650,
	-22.0320754965805000,0.2342044096857340,
	-22.0748475257641000,0.2333927721749480,
	-22.1177686775040000,0.2325798302271210,
	-22.1608402375941000,0.2317655792041130,
	-22.2040635098079000,0.2309500060800750,
	-22.2474398162371000,0.2301331006777330,
	-22.2909704976571000,0.2293148546784710,
	-22.3346569138969000,0.2284952575315180,
	-22.3785004441844000,0.2276742963221460,
	-22.4225024875678000,0.2268519620745500,
	-22.4666644632887000,0.2260282433951920,
	-22.5109878111783000,0.2252031287974460,
	-22.5554739920695000,0.2243766091263770,
	-22.6001244882573000,0.2235486721032510,
	-22.6449408038832000,0.2227193081843920,
	-22.6899244654092000,0.2218885009159310,
	-22.7350770220781000,0.2210562450646030,
	-22.7804000463603000,0.2202225233107620,
	-22.8258951344600000,0.2193873258887940,
	-22.8715639068006000,0.2185506430050440,
	-22.9174080085281000,0.2177124580860860,
	-22.9634291100430000,0.2168727623366480,
	-23.0096289075109000,0.2160315414646800,
	-23.0560091234612000,0.2151887827676820,
	-23.1025715072940000,0.2143444728653110,
	-23.1493178358997000,0.2134986004270870,
	-23.1962499142530000,0.2126511510107470,
	-23.2433695759918000,0.2118021115692910,
	-23.2906786841236000,0.2109514689810880,
	-23.3381791315848000,0.2100992063502050,
	-23.3858728419501000,0.2092453127307100,
	-23.4337617701221000,0.2083897745791730,
	-23.4818479030459000,0.2075325757110310,
	-23.5301332603892000,0.2066737005443940,
	-23.5786198953390000,0.2058131370146500,
	-23.6273098953284000,0.2049508686793330,
	-23.6762053828660000,0.2040868801144370,
	-23.7253085163376000,0.2032211563982180,
	-23.7746214908372000,0.2023536805538040,
	-23.8241465390008000,0.2014844398929880,
	-23.8738859319764000,0.2006134139856720,
	-23.9238419802608000,0.1997405881090430,
	-23.9740170346841000,0.1988659464467860,
	-24.0244134873910000,0.1979894726445320,
	-24.0750337728044000,0.1971111486300550,
	-24.1258803687144000,0.1962309546132080,
	-24.1769557973030000,0.1953488783193840,
	-24.2282626262562000,0.1944648966648920,
	-24.2798034699053000,0.1935789929452660,
	-24.3315809904040000,0.1926911509506980,
	-24.3835978989377000,0.1918013487516010,
	-24.4358569569467000,0.1909095692149760,
	-24.4883609774679000,0.1900157920161700,
	-24.5411128264192000,0.1891199970164250,
	-24.5941154240099000,0.1882221655996430,
	-24.6473717461327000,0.1873222768441420,
	-24.7008848258759000,0.1864203072391560,
	-24.7546577550032000,0.1855162410478630,
	-24.8086936855551000,0.1846100533262040,
	-24.8629958314289000,0.1837017216418930,
	-24.9175674701566000,0.1827912257153890,
	-24.9724119445623000,0.1818785422593670,
	-25.0275326645801000,0.1809636483252020,
	-25.0829331091326000,0.1800465219732020,
	-25.1386168280682000,0.1791271367066870,
	-25.1945874441638000,0.1782054696880850,
	-25.2508486551762000,0.1772814976297280,
	-25.3074042360294000,0.1763551925925140,
	-25.3642580410241000,0.1754265314684700,
	-25.4214140061474000,0.1744954882691980,
	-25.4788761514792000,0.1735620355568450,
	-25.5366485836812000,0.1726261460738050,
	-25.5947354985967000,0.1716877942527530,
	-25.6531411838747000,0.1707469494039890,
	-25.7118700218426000,0.1698035846502860,
	-25.7709264923582000,0.1688576713859560,
	-25.8303151758105000,0.1679091789232420,
	-25.8900407562572000,0.1669580771470890,
	-25.9501080246825000,0.1660043361562450,
	-26.0105218823719000,0.1650479234569800,
	-26.0712873444400000,0.1640888088146800,
	-26.1324095435069000,0.1631269564574950,
	-26.1938937335074000,0.1621623367322260,
	-26.2557452936488000,0.1611949122615670,
	-26.3179697325956000,0.1602246490179430,
	-26.3805726927751000,0.1592515142515970,
	-26.4435599548751000,0.1582754674253290,
	-26.5069374425148000,0.1572964759242430,
	-26.5707112272465000,0.1563144978004780,
	-26.6348875335292000,0.1553294968110060,
	-26.6994727442741000,0.1543414342172510,
	-26.7644734062341000,0.1533502668955370,
	-26.8298962360144000,0.1523559553542710,
	-26.8957481260687000,0.1513584575554070,
	-26.9620361511932000,0.1503577282773580,
	-27.0287675751184000,0.1493537265753080,
	-27.0959498575940000,0.1483464043484970,
	-27.1635906617467000,0.1473357163971110,
	-27.2316978616767000,0.1463216154150070,
	-27.3002795506625000,0.1453040526519230,
	-27.3693440495629000,0.1442829791107730,
	-27.4388999157127000,0.1432583423859250,
	-27.5089559522734000,0.1422300913836830,
	-27.5795212180627000,0.1411981704095000,
	-27.6506050378370000,0.1401625275107170,
	-27.7222170131776000,0.1391231047619380,
	-27.7943670338978000,0.1380798440814890,
	-27.8670652901353000,0.1370326869256420,
	-27.9403222849320000,0.1359815723839600,
	-28.0141488476734000,0.1349264369962150,
	-28.0885561481704000,0.1338672158169970,
	-28.1635557115931000,0.1328038451074700,
	-28.2391594341192000,0.1317362559988560,
	-28.3153795996077000,0.1306643772127280,
	-28.3922288972197000,0.1295881392270640,
	-28.4697204399811000,0.1285074664852110,
	-28.5478677845210000,0.1274222819883180,
	-28.6266849519500000,0.1263325101407940,
	-28.7061864500721000,0.1252380676682160,
	-28.7863872968439000,0.1241388726958520,
	-28.8673030454835000,0.1230348386036660,
	-28.9489498109841000,0.1219258772310960,
	-29.0313442984041000,0.1208118970472720,
	-29.1145038329746000,0.1196928035424030,
	-29.1984463922998000,0.1185685005773110,
	-29.2831906404561000,0.1174388850765190,
	-29.3687559647482000,0.1163038556199300,
	-29.4551625146711000,0.1151633034601830,
	-29.5424312438187000,0.1140171166560060,
	-29.6305839546101000,0.1128651812842310,
	-29.7196433463745000,0.1117073767504980,
	-29.8096330668406000,0.1105435800514960,
	-29.9005777674880000,0.1093736624398220,
	-29.9925031631342000,0.1081974898414290,
	-30.0854360959882000,0.1070149252630520,
	-30.1794046047932000,0.1058258259042810,
	-30.2744379994409000,0.1046300403406600,
	-30.3705669415161000,0.1034274145821830,
	-30.4678235316031000,0.1022177869276260,
	-30.5662414038120000,0.1010009888005820,
	-30.6658558283384000,0.0997768441040718,
	-30.7667038229798000,0.0985451697574103,
	-30.8688242744131000,0.0973057755481686,
	-30.9722580705706000,0.0960584613592770,
	-31.0770482449228000,0.0948030181482763,
	-31.1832401345222000,0.0935392264883869,
	-31.2908815530400000,0.0922668582278138,
	-31.4000229807129000,0.0909856731693040,
	-31.5107177732559000,0.0896954182871169,
	-31.6230223919759000,0.0883958289597899,
	-31.7369966579873000,0.0870866272924426,
	-31.8527040335136000,0.0857675188066164,
	-31.9702119336807000,0.0844381940071175,
	-32.0895920732898000,0.0830983263214789,
	-32.2109208533394000,0.0817475707887869,
	-32.3342797926167000,0.0803855621178905,
	-32.4597560114275000,0.0790119117047352,
	-32.5874427749811000,0.0776262087670696,
	-32.7174401057798000,0.0762280165641789,
	-32.8498554757278000,0.0748168687261570,
	-32.9848045912102000,0.0733922694957749,
	-33.1224122865776000,0.0719536866433936,
	-33.2628135448195000,0.0705005519197755,
	-33.4061546681914000,0.0690322553486653,
	-33.5525946260951000,0.0675481408387304,
	-33.7023066144623000,0.0660474999145740,
	-33.8554798678444000,0.0645295682513428,
	-34.0123217760936000,0.0629935176238923,
	-34.1730603697522000,0.0614384467671624,
	-34.3379472557309000,0.0598633745616101,
	-34.5072611059616000,0.0582672268435299,
	-34.6813118307827000,0.0566488249481765,
	-34.8604456083187000,0.0550068698584664,
	-35.0450509915024000,0.0533399234368277,
	-35.2355663882737000,0.0516463887739232,
	-35.4324893081340000,0.0499244787885793,
	-35.6363879090464000,0.0481721876548684,
	-35.8479155791052000,0.0463872466945184,
	-36.0678295804308000,0.0445670726525953,
	-36.2970152168076000,0.0427087042223450,
	-36.5365176500522000,0.0408087136283790,
	-36.7875845198325000,0.0388630969098111,
	-37.0517241768172000,0.0368671229285639,
	-37.3307870747662000,0.0348151298712040,
	-37.6270825637244000,0.0327002385759480,
	-37.9435517353031000,0.0305139398067884,
	-38.2840328003670000,0.0282454796242779,
	-38.6536871075325000,0.0258809030226199,
	-39.0597218908024000,0.0234014895917324,
	-39.5127059439955000,0.0207810181518199,
	-40.0292119846333000,0.0179805544484925,
	-40.6377530663221000,0.0149372301486279,
	-41.3955013165958000,0.0115352072795492,
	-42.4543433878785000,0.0075013897265686,
	-44.9381790145330000,0.0013364414409767,
	-45.2261985164242000,0.0009224484378228,
\end{filecontents*}



\title{
	\Large XianLiang猜想：\\[0.6em]
	\Large 宇宙学红移就是光速膨胀！ \\[0.6em]
	\Large $G(t)\tinysp c(t)^2 = \boldsymbol{P}$、修正后的爱因斯坦场方程 \textbf{$\boldsymbol G_{\mu\nu}$}，\\[0.6em]
	\Large 与零光速原初宇宙
}


\author{\\XianLiang Ma\\[8pt]
	\small{\texttt{Beijing DeepQuant.AI Technologies Co., Ltd}}\\
	\small{\texttt{中国 北京市朝阳区 100029}}}
\date{}


\setcounter{tocdepth}{1}

\begin{document}
	
	% First Page Zenodo.org DOI
	\AddToShipoutPictureFG*{%
		\AtPageLowerLeft{%
			\put(30,180){%
				\rotatebox{90}{%
					\textcolor[RGB]{90,105,168}{%
						\rmfamily\ptmbox\selectfont\LARGE
						{https://doi.org/10.5281/zenodo.23279668}
					}%
				}%
			}%
		}%
	}
	
	\maketitle
	
	\setlength{\leftskip}{1.1cm}
	\setlength{\rightskip}{1.1cm}
	
	\section *\small{ \centerline{\textbf{摘要}} }
	\medskip\noindent
	% 两侧各缩进 1cm（16-2=14cm），仅对摘要生效；正文自动复位
	
	\small{爱因斯坦 1905 年的狭义相对论（SR）与 1915 年的广义相对论（GR），只要求光速满足严格的洛伦兹参考系不变性；它们从未要求光速 $c$ 在百亿年的宇宙时标上保持恒定。这一关键区别，一个世纪以来一直被一个未经检验的宇宙学假设默默误读。在 1983 年的国际单位制（SI）定义中，$c$ 被定义为 299792458 m/s，其标称测量不确定度为 $\pm 1.1$ m/s，对应相对精度约为 $3.67\times 10^{-9}$。关键在于，若 $c(t)$ 的每日相对增长率处于 $10^{-13}\ \mathrm{day}^{-1}$ 量级——正如本研究的模型所预言——如此微小的变化便远低于 1983 年 SI 标准的精度极限。因此，这一细微漂移在 1983 至 2026 年整整 43 年间躲过了一切实验室测量，是完全可以理解的。狄拉克（1937）的大数假说曾设想引力参数 $G(t) \propto 1/t$ 与光速 $c(t)$ 随时间变化。费曼（1960）与戴维斯（1995）都指出，宇宙学红移的标准解释在其物理机制上存在一个未被回答的因果缺口。更晚近地，JWST（2023）与 DESI（2025）的数据揭示出一系列与当前 $\Lambda$CDM 模型的理论冲突，尤其在高红移区域以及暗能量的本质方面。秉承狄拉克、费曼与戴维斯所开创的怀疑与追问精神，本研究断言，这个未被回答的因果缺口恰恰暗示：宇宙学红移就是光速膨胀本身！正是在这一洞见的指引下，我们探索一个新的宇宙演化理论框架，它建立在两条核心公设之上：$(1)$ 光速随时间变化，而宇宙耦合常数 $P$ 保持不变，满足 $G(t)\,c(t)^2 = P$；$(2)$ 精细结构常数 $\alpha$ 永恒不变。在这两条假设之下，随时间增长的光速 $c(t)$ 成为宇宙演化的唯一驱动力，引力常数 $G$ 不再是基本常数，而是依赖于 $c(t)$ 的导出物理量。由 SI 米所定义的长度这一物理量纲，在本框架下被重新审视。真空磁导率 $\mu_0$、真空介电常数 $\varepsilon_0$、氢原子的玻尔半径 $a_{\scriptscriptstyle\mathrm{B}}$、史瓦西半径 $r_s$、约化普朗克常数 $\hbar$、普朗克质量 $m_\text{p}$、普朗克长度 $\ell_\text{p}$、普朗克密度 $\rho_\text{p}$，以及玻尔兹曼电离公式中的热力学温度 $T_\mathrm{ion}$，都随 \(c(t)\) 一同演化。含时函数 $c(t)$ 被推断为与宇宙标度因子 $a(t)$ 以及标准 $\Lambda$CDM 模型中定义的红移因子 $1+z$ 同步标度。在与 $1+z=1100$ 对应的 CMB 退耦时期，光速 $c(t)$ 为当前光速 $c_0$ 的 $1/1100$，玻尔半径 $a_{\scriptscriptstyle\mathrm{B}}$ 为其当前值的 $1/1100$，电离温度 $T_\mathrm{ion}$ 仅为 $82.64\,$K，这一结果与标准 $\Lambda$CDM 模型惯常的高温假设显著不同。在该增长 $c(t)$ 模型中，宇宙形成时标被延展至 $[-59.1,\,-45.2]\,$$\mathrm{Gyr}$，远超标准模型所假定的 -13.8$\,\mathrm{Gyr}$，这暗示着一个缓慢演化的宇宙，而非剧烈的大爆炸图景。本文重新审视了爱因斯坦曲率场方程 $G_{\mu\nu}$ 的解读，揭示出在早期星系形成时期（$1+z=25$）与 CMB 时期（$1+z=1100$）存在极高的时空曲率。本文还探讨了宇宙耦合常数 $P$ 的内在含义，其关系式为 $P = \dfrac{c(t)^6 G_{\mu\nu}}{8\pi T_{\mu\nu}} \approx \dfrac{c(t)^4 G_{\mu\nu}}{8\pi M_{\mu\nu}}$。此处 $P$ 刻画了宇宙质能分布、时空曲率与含时光速 $c(t)$ 之间的基本耦合。在宏观尺度上，动态的质量分布同时支配着时空曲率与光传播速度的演化。光速的演化编码了宇宙从早期完全均匀的强子化后等离子体状态，向物质与能量在时空中局域成团的转变。其当前的相对日变化率被计算为 $[0.86,\,1.89]\times 10^{-13}\ \mathrm{day}^{-1}$，并借助高能光子源（HEPS）与拟议中的 Mission-M 研究计划，探讨了检验这一变化的可行性。DESI DR2 对动态暗能量的偏好、宇宙光子预算危机、早期星系恒星质量张力、星族 III 恒星时标冲突、早期太阳锂丰度之谜，以及高级智慧生命为何出现得如此之晚等悬而未决的问题，都可以在该增长-\(c(t)\) 模型下重新审视。我们将宇宙的终极起源设定为 \textit{ZLZT}（零光速与零温度）原初态，其中光速与温度都趋于零。同时，本文定义了 \textit{LEM-SCDF}（光—电—磁同步宇宙暗场），用以实现电磁参数的全局同步。更重要的是，这一抛弃了常数 \(G\)、并配备动态 $c(t)$ 的修正版爱因斯坦场方程，为探索人们长期追寻的统一场论提供了一个突破性的视角与可行的路径。 \\}
	
	\setlength{\leftskip}{0cm}
	\setlength{\rightskip}{0cm}
	
	
	\section*{
		\begin{center}
			目\quad 录
		\end{center}
	}
	\medskip
	\noindent
	
	\noindent
	1. 经典基本方程
	
	\medskip\noindent
	2. 基本公设
	
	\medskip\noindent
	3. 符号定义与约定
	
	\medskip\noindent
	4. 被忽视的光速时变性
	
	\medskip\noindent
	5. S 型增长的光速 \(c(t)\) 模型
	
	\medskip\noindent
	6. 超精细增长 \(c(t)\) 模型
	
	\medskip\noindent
	7. 增长 \(c(t)\) 宇宙的不同物理图景
	
	\medskip\noindent
	8. 修正后的 ${G_{\mu\nu}}$、常数 $P$ 与 LEM‑SCDF 场
	
	\medskip\noindent
	9. ZLZT 原初态、囚禁生成物质与 ${c(t)}$ 的演化
	
	\medskip\noindent
	10. 宇宙学红移：一个未被回答的因果缺口及其物理闭合
	
	\medskip\noindent
	11. 关于微小演化 ${c(t)}$ 实验可探测性的讨论
	
	\medskip\noindent
	12. 以增长 \(c(t)\) 框架重新审视若干悬而未决的问题
	
	\medskip\noindent
	后记
	
	\medskip\noindent
	附录：倘若该猜想成立
	
	\medskip\noindent
	参考文献
	


	\section*{
		\begin{center}
			1. 经典基本方程
		\end{center}
	}
	\medskip
	\noindent
	
	本论文建立在若干物理学基本方程之上，它们为后续的理论推导奠定基础。为清晰与完整起见，下面简要列出这些标准方程并附简短注释，作为我们修正后的增长-${c(t)}$ 框架的基准。
	
	\bigskip
	真空中的光速与真空磁导率、真空介电常数相关，可由麦克斯韦方程组导出：
	$$
	c = \frac{1}{\sqrt{\mu_0\tinysp \varepsilon_0}} \enspace \quad \Rightarrow \quad \enspace
	\mu_0\tinysp \varepsilon_0\tinysp c^2 = 1 \eqno(1.1)
	$$
	
	\smallskip
	精细结构常数是一个无量纲基本常数，它描述电磁耦合强度，构成一切微观物质质量背后最基本的内禀量：
	$$
	\alpha = \frac{e^2}{4\pi\tinysp \varepsilon_0\tinysp \hbar\tinysp c}\eqno(1.2)
	$$
	
	\smallskip
	普朗克质量，即自然单位制中的基本质量单位：
	$$
	m_{\scriptscriptstyle\mathrm{P}} = \sqrt{\frac{\hbar\tinysp c}{G}}\eqno(1.3)
	$$
	
	\smallskip
	普朗克长度，即自然单位制中的基本长度单位：
	$$
	\ell_{\scriptscriptstyle\mathrm{P}} = \sqrt{\frac{\hbar\tinysp G}{c^3}}\eqno(1.4)
	$$
		
	\smallskip
	普朗克密度与普朗克力，即量子引力中的特征密度与特征力量级，以比例形式表示：
	$$
	\rho_{\scriptscriptstyle\mathrm{P}} \propto \frac{m_{\scriptscriptstyle\mathrm{P}}}{{\ell_{\scriptscriptstyle\mathrm{P}}}^3} = \frac{c^5}{\hbar\tinysp G^2}, \qquad
	F_{\scriptscriptstyle\mathrm{P}} = \frac{c^4}{G}
	\eqno(1.5)
	$$
	
	
	\smallskip
	一个联系微观普朗克量与宏观常数 $Gc^2$ 的唯象关系式。该式仅由量纲分析导出，尚无任何具有物理意义的内在解释。
	$$
	\frac{\hbar^4} {m_{\scriptscriptstyle\mathrm{P}}^5\tinysp \ell_{\scriptscriptstyle\mathrm{P}}^3} =
	\frac{\hbar^4\tinysp \rho_{\scriptscriptstyle\mathrm{P}}} {m_{\scriptscriptstyle\mathrm{P}}^6} =
	G\tinysp c^2\eqno(1.6)
	$$
	
	\smallskip
	氢原子的玻尔半径，即基态氢原子中电子轨道的最概然半径，如下：
	$$
	a_{\scriptscriptstyle\mathrm{B}} = \frac{4\pi\tinysp \varepsilon_0\tinysp \hbar^2}{m_e e^2}\tinysp \eqno(1.7)
	$$
	
	\smallskip
	氢原子的玻尔兹曼电离判据，以两种等价形式表述：
	$$
	k_{\scriptscriptstyle\mathrm{B}}\tinysp T_\mathrm{ion} > \frac{e^2}{4\pi\tinysp \varepsilon_0\tinysp a_\text{B}} = \alpha^2\tinysp m_e\tinysp c^2\eqno(1.8)
	$$
	
	\smallskip
	不含宇宙学常数的广义相对论爱因斯坦曲率场方程：
	$$
	G_{\mu\nu} = \frac{8\pi G}{c^4} T_{\mu\nu}\eqno(1.9)
	$$
	
	\smallskip
	含宇宙学常数的爱因斯坦曲率场方程，描述带有暗能量贡献的时空曲率：
	$$
	G_{\mu\nu} + \Lambda g_{\mu\nu} = \frac{8\pi G}{c^4} T_{\mu\nu}\eqno(1.10)
	$$
	
	\smallskip
	史瓦西时空中的引力频移，描述光子频率如何随引力势变化；史瓦西半径则刻画了这一效应的尺度：	
	\[
	\frac{\nu}{\nu_0} = \sqrt{1 - \frac{2GM}{c^2r}}, \qquad r_s = \frac{2GM}{c^2}\eqno(1.11)
	\]
	
	\smallskip
	牛顿万有引力定律，即引力的经典描述：
	$$
	F_g = G \,\frac{m_1 m_2}{r^2}\eqno(1.12)
	$$
	
	\smallskip
	真空中两个点电荷之间静电相互作用的库仑定律，其中 $k_e$ 表示库仑常数，$\varepsilon_0$ 为真空介电常数：
	$$
	F_e = k_e \frac{q_1 q_2}{r^2}, \enspace \quad k_e=\frac{1}{4\pi\varepsilon_0}\eqno(1.13)
	$$
	
	\smallskip
	由爱因斯坦场方程导出的弗里德曼方程，描述宇宙的膨胀。在本研究中，我们把 ${\rho_{\scriptscriptstyle\mathrm{F}}(t)}$ 定义为宇宙平均密度，以区别于其他密度情形。第一个方程定义了哈勃参数 $H(t) = \dot{a}(t)/a(t)$：
	$$
	H^2(t) = \left(\frac{\dot{a}(t)}{a(t)}\right)^2 = \frac{8\pi G}{3}{\rho_{\scriptscriptstyle\mathrm{F}}(t)} - \frac{kc^2}{a(t)^2} + \frac{\Lambda c^2}{3}\eqno(1.14)
	$$
	
	\smallskip
	第二个弗里德曼方程支配宇宙膨胀的加速度，包含物质、辐射与暗能量的贡献：
	$$
	\frac{\ddot{a}(t)}{a(t)} = -\frac{4\pi G}{3}\left({\rho_{\scriptscriptstyle\mathrm{F}}(t)} + \frac{3p(t)}{c^2}\right) + \frac{\Lambda c^2}{3}\eqno(1.15)
	$$
		
	\smallskip
	狄拉克注意到两个大数巧合。其一是电子与质子之间电磁力 ${F_\text{em}}$ 与引力 ${F_\text{g}}$ 之比，其二是宇宙年龄 ${T_{0}}$ 与质子的光穿越时间之比。 
	$$
	N_{\text{Dirac,1}} = \frac{F_\text{em}}{F_\text{g}} 
	= \frac{e^2}{ 4\pi\,G\,\varepsilon_0\,m_e\,m_p} \approx 2.27 \times 10^{39} \eqno(1.16)
	$$
	
	$$
	N_{\text{Dirac,2}} = \frac{T_{0}}{d_p / c} = \frac{c \, T_{0}}{d_p} \approx 7.67 \times 10^{40} \eqno(1.17)
	$$
	
	
	
	\section*{
		\begin{center}
			2. 基本公设
		\end{center}
	}
	\medskip
	\noindent
	
	本研究建立在两条互补的核心公设之上，它们定义了描述时变宇宙的理论框架。

	\bigskip	
	{\bf 宏观公设}\,\,\,{存在一个普适的宇宙耦合常数 \(\boldsymbol{P}\)，它在宇宙的整个历史中始终保持不变。光速 \(c(t)\) 并非常数，而是宇宙演化的含时函数，满足如下基本守恒律：} 
	$$
	{G(t)}\tinysp {c(t)^2} = \boldsymbol{P}\eqno(2.1)
	$$
	{其中 \(G(t)\) 为含时引力系数。在这一公设之下，\(c(t)\) 成为宇宙演化的首要驱动力，而 \(G(t)\) 不再是基本常数，而是一个导出的物理量。}

	\medskip
	\textit{\textbf{注：}光速 \(c(t)\) 是一个全局宇宙参数。在任意给定时刻 \(t\)，真空中的所有光子共享同一速度 \(c(t)\)，与它们发射时刻的光速无关。这一共同速度由一个迄今尚未被探究的宇宙暗场所支配，后文将把它确认为 LEM-SCDF 场。}

	\bigskip
	{\bf 微观公设}\,\,\,{精细结构常数 \(\boldsymbol{\alpha}\) 在 CMB 之后的整个时期普适不变，从中性原子形成时期直至今日始终保持恒定。在本框架中，其定义式为：}
	$$
	\boldsymbol{\alpha} = \frac{e^2}{4\pi{\varepsilon}_0(t)\tinysp {\hbar}(t)\tinysp {c}(t)}\eqno(2.2)
	$$
	{即使 \(\varepsilon_0(t)\)、\(\hbar(t)\) 与 \(c(t)\) 本身随宇宙时间演化，它依然保持不变。在 CMB 退耦之前的原初等离子体阶段，这一不变性尚无定义。}
	
	
	
	
	


	\section*{
		\begin{center}
			3. 符号定义与约定
		\end{center}
	}
	\medskip
	\noindent
	
	本节定义并统一后文讨论中使用的关键符号、下标规则与参考时期，以确保记号一致、物理含义清晰。
	
	\medskip
	\noindent
	
	{\bf 宇宙参考时期（下标 $z$）}\,\,\,基准宇宙时期以红移定义；下标 $\boldsymbol*$ 表示 CMB 时期之前的原初时期：
	\begin{itemize}
		\item $\boldsymbol{z_{*}}$：QGP 囚禁阈值时期（强子化时期）
		\item $\boldsymbol{z_0}$：$1+z=1100$，CMB 退耦时期
		\item $\boldsymbol{z_1}$：$1+z=25$，再电离晚期与早期星系形成时期
		\item $\boldsymbol{z_2}$：$1+z=3.7$，恒星形成与黑洞活动的高峰时期，即宇宙正午
		\item $\boldsymbol{z_3}$：$1+z=1.43$，太阳系形成时期
	\end{itemize}
	
	\medskip
	\noindent
	
	{\bf 归一化光速与积分量}\,\,\,我们定义两个核心量，用以刻画时变的宇宙光速及其累积传播效应：
	\begin{itemize}
		\item $\boldsymbol{C(t)}$：无量纲宇宙光速，定义为演化中的光速 $c(t)$ 与其当前值 $c_0$ 之比：
		\[
		C(t) = \dfrac{c(t)}{c_0}
		\]
		
		\item $\boldsymbol{D(t_1,t_2)}$：积分有效光行时，由 $C(t)$ 在任意宇宙时间区间 $[t_1, t_2]$ 上作路径积分得到：
		\[
		D(t_1,t_2) = \int_{t_1}^{t_2} C(t)\,dt
		\]
		$D(t_1,t_2)$ 以 $\mathrm{Gyr}$ 为单位。当以当前光速 $c_0$ 折算时，这一时间量对应传统宇宙学所采用的标准光行距离（单位为十亿光年，Gly）。
		
	\end{itemize}
	
	\medskip
	\noindent
	
	{\bf 光速}\,\,\,在各基准时期取值的归一化光速：
	\begin{itemize}
		\item $\left.\boldsymbol{C(t)}\right|_{\boldsymbol{z_{*}}}$：夸克囚禁阈值时期（$z_{*}$）的光速
		\item $\left.\boldsymbol{C(t)}\right|_{\boldsymbol{z_0}}$：CMB 时期（$z_0$）的光速
		\item $\left.\boldsymbol{C(t)}\right|_{\boldsymbol{z_1}}$：早期星系形成时期（$z_1$）的光速
		\item $\left.\boldsymbol{C(t)}\right|_{\boldsymbol{z_2}}$：宇宙正午时期（$z_2$）的光速
		\item $\left.\boldsymbol{C(t)}\right|_{\boldsymbol{z_3}}$：太阳系形成时期（$z_3$）的光速
	\end{itemize}
	
	\medskip
	\noindent
	
	{\bf 光子本征频率}\,\,\,光子的固定本征参考频率，由原子能级与精细结构常数确定：
	\begin{itemize}
		\item $\boldsymbol{\nu_{n0}}=1.0\times10^{15}\ \text{Hz}$：主本征频率
	\end{itemize}
	
	\medskip
	\noindent
	
	{\bf 介质折射率}\,\,\,
	\begin{itemize}
		\item $\boldsymbol{N}_{\text{glass}}=1.5$：标准光学玻璃的折射率（固定值）
	\end{itemize}
	
	\medskip
	\noindent
	
	{\bf 等离子体电离温度}\,\,\,各基准时期氢原子的临界玻尔兹曼电离温度：
	\begin{itemize}
		\item $\left.\boldsymbol{T_\mathrm{ion}}\right|_{\boldsymbol{z_0}}$：CMB 时期（${z_0}$）的临界温度
		\item $\left.\boldsymbol{T_\mathrm{ion}}\right|_{\boldsymbol{z_1}}$：早期星系形成时期（${z_1}$）的临界温度
		\item $\left.\boldsymbol{T_\mathrm{ion}}\right|_{\boldsymbol{z_2}}$：宇宙正午时期（${z_2}$）的临界温度
		\item $\left.\boldsymbol{T_\mathrm{ion}}\right|_{\boldsymbol{z_3}}$：太阳系形成时期（${z_3}$）的临界温度
	\end{itemize}
	
	\medskip
	\noindent
	
	{\bf 爱因斯坦时空曲率}\,\,\,在各基准时期取值的爱因斯坦曲率张量：
	\begin{itemize}
		\item $\left.\boldsymbol{G_{\mu\nu}}\right|_{\boldsymbol{z_0}}$：CMB 时期（${z_0}$）的曲率
		\item $\left.\boldsymbol{G_{\mu\nu}}\right|_{\boldsymbol{z_1}}$：早期星系形成时期（${z_1}$）的曲率
		\item $\left.\boldsymbol{G_{\mu\nu}}\right|_{\boldsymbol{z_2}}$：宇宙正午时期（${z_2}$）的曲率
		\item $\left.\boldsymbol{G_{\mu\nu}}\right|_{\boldsymbol{z_3}}$：太阳系形成时期（${z_3}$）的曲率
	\end{itemize}
	
	\medskip
	\noindent
	
	{\bf 当前宇宙参考值（下标 $0$）}\,\,\,
	在{当前宇宙时期}取值的基本物理常数与宇宙学尺度，作为贯穿本工作的普适参考基准：
	\begin{itemize}
		\item $\boldsymbol{c_0}$：当前真空光速
		\item $\boldsymbol{G_0}$：当前引力常数
		\item $\boldsymbol{a_0}$：当前宇宙标度因子（弗里德曼度规）
	\end{itemize}
	
	\medskip
	\noindent
	
	{\bf 宇宙时间的三种定义}\,\,\,
	我们引入三个不同的宇宙时间变量，以区分本工作所发展的{\it 变光速宇宙学}与标准的{\it 恒光速宇宙学}：
	\begin{itemize}
		\item $\boldsymbol{t}$：{\it 变光速框架}中的宇宙时间。
		按约定，处处 $t \le 0$，并把{\it 当前时期}设为 $t = 0$。
		更古老的宇宙时期，对应{\it 更负的 $t$ 值}。
		$t$ 的单位为 $\mathrm{Gyr}$（十亿年）。
		
		\item $\boldsymbol{t^\star}$：{\it 标准恒光速宇宙学}中的宇宙时间。
		这是 $\Lambda$CDM 中通用的宇宙时间，同样定义为 $t^\star \le 0$，且当前时期为 $t^\star = 0$。
		{\it 宇宙正午}（宇宙恒星形成活动的高峰）在标准宇宙学中对应于
		$t^\star \approx -11.4\,\mathrm{Gyr}$。
		
		\item $\boldsymbol{t^{\Delta}}$：标准 $\Lambda$CDM 模型中定义的正向宇宙时间。它以宇宙大爆炸时刻为原点 $t^\Delta = 0$，故恒有 $t^\Delta \ge 0$。
		当前时期对应于 $t^\Delta = 13.8\,\mathrm{Gyr}$。
		
		\item[]
		特别地，{\it 宇宙正午}在这一正向时间系统中位于 $t^\Delta = 2.4\,\mathrm{Gyr}$。这两个标准宇宙学时间坐标满足：
		\[
		t^\Delta + |t^\star| = 13.8\,\mathrm{Gyr}.\eqno(3.1)
		\]
		
		\item[] 
		\textit{\textbf{注：}上标 $\star$ 表示以当前为零时的标准恒光速宇宙学下的量。上标 $\vartriangle$ 表示从大爆炸起算的正向 $\Lambda$CDM 时间，常用于红移分析。这一记号体系把我们的框架与传统宇宙学区分开来。}
		
	\end{itemize}
	
	
	
	
	
	
	
	\section*{
		\begin{center}
			4. 被忽视的光速时变性
		\end{center}
	}
	\medskip
	\noindent
	
	{\bf 公理化的陷阱}\,\,\,
	一个多世纪以来，光速 \(c\) 的恒定性已被抬升为不容置疑的教条。爱因斯坦的理论在局域范围内取得了非凡的成功，却无意间施加了一重认知上的囚禁：研究者习惯于把实验室尺度的测量结果投射到整条宇宙时间线上。这种公理化的惯性，使得少有学者敢于探索光速可变的前景。
	
	\medskip
	然而，爱因斯坦 1905 年的狭义相对论（SR）与 1915 年的广义相对论（GR），只要求 \(c\) 具有\textbf{\textit{严格的}}洛伦兹参考系不变性——它们\textbf{\textit{从未}}要求 \(c\) 在百亿年的宇宙时期中恒定不变。这一微妙的区别，一个世纪以来一直被默默误读为一条普适的宇宙学假设。自 R\o mer（1676）首次指出光以有限速度传播以来，人类对光速 \(c\) 的测量跨度尚不足四个世纪——对宇宙这位巨人而言，不过一次心跳。
	
	\medskip
	在 1983 年的 SI 修订中，光速被固定为 $299\,792\,458 \mathrm{m/s}$，NBS 报告~\cite{NBS:1987} 记载其相对精度约为 \(3.67\times 10^{-9}\)。关键在于，若 \(c(t)\) 的相对日增长率处于 \(10^{-13}\ \mathrm{day}^{-1}\) 量级——正如第 5 节与第 6 节后续模型所提示的那样——那么在 1983 年以来的 43 年间，没有任何实验能够探测到如此微小的漂移，因为它远低于 1983 年标准的精度极限。
		
	
	\bigskip				
	{\bf 观测的不对称性：根本性的盲区}\,\,\,
	现代宇宙学建立在一个显眼却长期被忽视的不对称之上。远古光子及其红移现象是可以观测的，但对于早期宇宙中真实的光速 \(c\) 与引力强度 \(G\)，我们掌握的直接证据为\textbf{\textit{零}}。标准模型把宇宙膨胀的一切后果都约化为频率红移这一唯一可观测现象，同时盲目地把光速 \(c\) 永久固定。这一判断缺乏实际检验，只不过是一条从未经受严格审查、缺乏依据的形而上学预设。
	\bigskip
	
	{\bf 拓扑的必然性：二维几何的强制性要求}\,\,\,
	标准宇宙学框架的致命几何缺陷正在于此。宇宙膨胀并非一条线的一维拉伸，而是弯曲的二维时空流形的均匀标度形变。黎曼几何与拓扑学严格规定：度规曲面的各向同性膨胀，必须把其形变分配到所有固有的自由度之上。光具有两种相互共轭的运动学属性：本征频率与传播速度。把全部度规形变都只分配给频率，是一种数学上的反常。标准假设实际上截去了时空一半的自由度，导致一种物理上不完整、拓扑上不自洽的描述。
	\bigskip
	
	{\bf 恢复对称性的公设}\,\,\,
	为修复这一被破坏的几何对称性，最简约、也在数学上最不可避免的结论是：
	\smallskip
	$$
	\boldsymbol{\frac{\nu_0}{\nu_e}} \propto \boldsymbol{\frac{a(t_e)}{a_0}} \enspace \quad \Rightarrow \quad \enspace  \boldsymbol{\frac{c(t)}{c_0}} \propto \boldsymbol{\frac{a(t)}{a_0} }
	\eqno(4.1)
	$$
	
	\smallskip
	光速必须与宇宙标度因子同步演化。这不仅仅是一个物理假设，更是恢复膨胀宇宙之拓扑完整性的必然要求。
	
	\smallskip
	\textit{\textbf{注：}在后文第 10 节中，我们通过质疑标准宇宙学红移理论内部的\textbf{因果缺口}，进一步论证宇宙学红移就是光速膨胀本身。}
	
	\bigskip
	在接下来的三节中，我们构建两个不同的增长-\(c(t)\) 模型：第 5 节中 \textit{S 型}的紧凑模型，以及第 6 节中经 \texttt{Planck18} 标定的超精细模型。随后在第 7 节，我们尝试分析在增长-\(c(t)\) 假设之下，不同宇宙时期所呈现出的截然不同的物理图景。
	
	
	
	
	
	
	

	\section*{
		\begin{center}
			5. S 型增长的光速 $\boldsymbol{c(t)}$ 模型
		\end{center}
	}
	\medskip
	\noindent
	
	本节引入一个具有物理动机的 \textit{S 型}模型，用来描述含时光速 $c(t)$，旨在刻画光速从早期宇宙到今日自然而渐进的抬升。
	\bigskip
	
	{\bf 标准 S 型函数}\,\,\,
	S 型函数是一个光滑、单调且有界的函数，广泛用于生物学、植物学、生命科学与数据科学，以描述基本的增长、成熟、饱和与转变现象。其标准形式为
	\[
	\sigma(x)=\frac{1}{1+e^{-x}}\,\,,\eqno(5.1)
	\]
	
	其中对一切实数 $x$ 有 $\sigma(x)\in(0,1)$。该函数关于点 $(0,\,0.5)$ 对称，并在整个实域上单调递增。为清晰起见，表 1 列出若干代表性节点处的关键数值。
	
	\begin{table}[H]
		\centering
		\captionsetup{skip=2pt, font=small}
		\caption{标准 S 型函数 $\sigma(x)$ 的取值}
		\footnotesize
		\renewcommand{\arraystretch}{1.20}
		\begin{tabular}{c|ccccccccc}
			\toprule
			$\sigma(x)$ & $0.047$ & $0.119$ & $0.269$ & $0.378$ & $\mathbf{0.500}$ & $0.622$ & $0.731$ & $0.881$ & $0.953$ \\
			\midrule
			$x$ & $-3.0$ & $-2.0$ & $-1.0$ & $-0.5$ & $0.0$ & $+0.5$ & $+1.0$ & $+2.0$ & $+3.0$ \\
			\bottomrule
		\end{tabular}
	\end{table}
	
	下面给出 S 型函数 $\sigma(x)$ 的标准数学曲率 $\kappa(x)$ 表达式，不再详加推导：
	\[
	\kappa(x) = \frac{\bm{\left|}\ddot{\sigma}(x)\bm{\right|}}{\big[1 + (\dot{\sigma}(x))^2\big]^{3/2}} = \frac{\bm{\left|}\sigma(x)\big[1 - \sigma(x)\big]\big[1 - 2\sigma(x)\big]\bm{\right|}}{\big\{1 + \big[\sigma(x)\big(1 - \sigma(x)\big)\big]^2\big\}^{3/2}}\,\,,\eqno(5.2)
	\]
	
	\medskip
	曲率函数 $\kappa(x)$ 有三个关键点：两个对称的极大值与一个极小值。左极大值出现在 $x = -\ln(2+\sqrt{3}) \approx -1.317$，右极大值出现在 $x = \ln(2+\sqrt{3}) \approx +1.317$。在 $x=0$ 处曲率为零，此即曲线改变凹凸性的拐点。当 $x \to \pm\infty$ 时，曲率渐近趋于零。
	\medskip
	
	值得注意的是，若宇宙红移函数 $z(t)$ 取指数形式 $z(t) \propto e^{-t}$，则量 $1/\big(1+z(t)\big)$ 与标准 S 型函数 $\sigma(x) = 1/\big(1+e^{-x}\big)$ 高度相似。若这一对应成立，便能在 S 型函数形式与宇宙红移演化之间架起一座自然的桥梁，为下一小节所述的分析奠定坚实基础。
	\bigskip
	\medskip
	
	
	
	{\bf 红移曲线 $\boldsymbol{z(t)}$：形态与标定}\,\,\,
	我们使用 Python 库 \texttt{astropy.cosmology}~\cite{Astropy:2022} 获取标准的 $\Lambda$CDM 红移—时间关系，具体采用 \texttt{Planck18} 宇宙学实现~\cite{Planck2018:cosmo}，它给出的红移值作为自大爆炸以来正向宇宙时间 $t^\Delta$ 的函数，在数值上保持一致。
	
	
	观测研究中一个惯常使用的标定基石，是所谓的{\textit{宇宙正午}}时期，它对应于 $t^\Delta = 2.4\, \mathrm{Gyr}$，匹配红移 $z = 2.7$。我们采用这一参考点来构建简化的指数模型：
	\[
	z(t^\Delta)
	= \underbrace{A^\Delta\,e^{-k^\Delta \left(t^\Delta - \beta^\Delta\right)}}_{\text{指数衰减}} \eqno(5.3)
	\]
	
	各参数取为
	\[
	A^\Delta = 2.7,\quad \beta^\Delta = 2.4\,\mathrm{Gyr},\quad k^\Delta = 1.0\,\,,\eqno(5.4)
	\]
	如此，指数宗量恰好在宇宙正午时期等于零：
	\[
	t^\Delta - \beta^\Delta = 0 \quad \Rightarrow \quad e^{0} = 1 \quad \Rightarrow \quad z(t^\Delta) = A^\Delta\,\,.\eqno(5.5)
	\]
	
	\bigskip
	尽管上述简化模型抓住了宇宙红移的总体趋势，我们仍构建了一个改进的复合模型，以更好地拟合 $\Lambda$CDM \texttt{Planck18} 的数值数据。我们在基本式的基础上再加入两项：一个双曲修正项与一个伪椭圆调整项。完整表达式为：
	\medskip
	\[
	z(t^\Delta)
	=
	\underbrace{A^\Delta\,e^{-k^\Delta \left(t^\Delta - \beta^\Delta\right)}}_{\text{指数衰减}}
	\cdot
	\underbrace{\left(\frac{t^\Delta}{\beta^\Delta}\right)^{k^\Delta - 1}}_{\text{双曲修正}}
	\cdot
	\underbrace{\left(\frac{13.8 - t^\Delta}{13.8 - \beta^\Delta}\right)^{1 - k^\Delta}}_{\text{伪椭圆调整}}\eqno(5.6)
	\]
	
	\medskip
	双曲修正项减缓衰减速率，使曲线形态更接近 $\Lambda$CDM。这两个附加项采用互补的指数，以避免数值奇点与分母为零的问题。伪椭圆项则强制今日边界条件 $z(t^\Delta) = 0$ 在 $t^\Delta = 13.8\,\mathrm{Gyr}$ 处成立，这是指数模型所无法达到的特性。
	
	\medskip
	尽管外观复杂，改进模型只包含一个独立的自由参数 $k^\Delta$。用于整体形态调优的标定参数集取为
	\medskip
	\[
	A^\Delta = 2.7,\quad \beta^\Delta = 2.4\,\mathrm{Gyr},\quad k^\Delta = 0.16\,\,,\eqno(5.7)
	\]
	
	\medskip
	下图比较了两个模型所给出的红移曲线。
	\medskip
	
	
	\begin{figure}[H]
		\centering
		% -------------------- 左图 (a) 简单指数拟合 --------------------
		\begin{minipage}[b]{0.48\textwidth}
			\centering
			\begin{tikzpicture}
				\begin{axis}[
					width=7cm, height=6cm,
					axis lines = box,
					xlabel = {(a) 宇宙年龄 $t^\Delta$ [Gyr]},
					ylabel = {红移 $z$},
					xmin=0.2, xmax=13.8,
					ymin=0, ymax=20,
					xtick={0.0,2.5,5.0,7.5,10.0,12.5},
					ytick={0,5,10,15,20},
					tick label style={font=\small},
					label style={font=\small},
					smooth,
					every tick/.style={thick, black},
					]
					
					% 红移曲线
					\addplot[color=red, line width=1.4pt, domain=0.2:13.8, samples=300]
					{2.7*exp(-1*(x-2.4))};
					
					% 参考线 t = 2.4 Gyr
					\draw[blue!80!black, line width=0.8pt] (axis cs:2.4, 0) -- (axis cs:2.4, 20);
					% 参考线 z = 2.7
					\draw[blue!80!black, line width=0.8pt] (axis cs:0.2, 2.7) -- (axis cs:13.8, 2.7);
					
					% 文字标注
					\node[blue!80!black, font=\small, anchor=south west, rotate=90] at (axis cs:2.65, 9) {$t^\Delta=2.4\ \mathrm{Gyr}$};
					\node[blue!80!black, font=\small, anchor=south west] at (axis cs:6.5, 2.9) {$z=2.7$};
					
				\end{axis}
			\end{tikzpicture}
		\end{minipage}
		\hfill
		% -------------------- 右图 (b) 精细复合拟合 --------------------
		\begin{minipage}[b]{0.48\textwidth}
			\centering
			\begin{tikzpicture}
				\begin{axis}[
					width=7cm, height=6cm,
					axis lines = box,
					xlabel = {(b) 宇宙年龄 $t^\Delta$ [Gyr]},
					ylabel = {红移 $z$},
					xmin=0.2, xmax=13.8,
					ymin=0, ymax=20,
					xtick={0.0,2.5,5.0,7.5,10.0,12.5},
					ytick={0,5,10,15,20},
					tick label style={font=\small},
					label style={font=\small},
					smooth,
					every tick/.style={thick, black},
					]
					
					% 精细复合拟合曲线
					\addplot[color=red, line width=1.4pt, domain=0.2:13.8, samples=300]
					{2.7*exp(-0.16*(x-2.4)) * ((x/2.4)^(-0.84)) * (((13.8-x)/(13.8-2.4))^0.86)};
					
					% 参考线 t = 2.4 Gyr
					\draw[blue!80!black, line width=0.8pt] (axis cs:2.4, 0) -- (axis cs:2.4, 20);
					% 参考线 z = 2.7
					\draw[blue!80!black, line width=0.8pt] (axis cs:0.2, 2.7) -- (axis cs:13.8, 2.7);
					
					% 文字标注
					\node[blue!80!black, font=\small, anchor=south west, rotate=90] at (axis cs:2.65, 9) {$t^\Delta=2.4\ \mathrm{Gyr}$};
					\node[blue!80!black, font=\small, anchor=south west] at (axis cs:6.5, 2.9) {$z=2.7$};
					
				\end{axis}
			\end{tikzpicture}
		\end{minipage}
		
		\vspace{0.1cm}
		
		% 关键修改：用 footnotesize 让图注略小
		\captionsetup{width=0.80\textwidth, font=footnotesize}
		\caption{两个模型给出的红移 $z(t^\Delta)$ 曲线：
			(a) 简单指数模型，式~(5.3)；\enspace \enspace
			(b) 由指数项、双曲项与伪椭圆项构成的改进复合模型，式~(5.6)。\enspace \enspace
			两模型均锚定于参考点 $t^\Delta=2.4\ \mathrm{Gyr}$ 与 $z=2.7$。}
		\label{fig:redshift_fit_comparison}
		
	\end{figure}
	
	
	
	\medskip
	简单指数模型式 (5.3) 与改进复合模型式 (5.6) 都能很好地重现观测红移曲线的指数衰减形态。同时，改进模型在所有宇宙时期对 $\Lambda$CDM \texttt{Planck18} 数据取得了更优的整体拟合。
	\medskip
	
	尽管如此，我们仍惯常采用简单的指数拟设 $z(t) \propto e^{-t}$ 来推进关于变光速的研究。$c(t)$、标度因子 $a(t)$ 与红移 $z(t)$ 之间的关系如下：
	\medskip
	\[
	\boldsymbol{\frac{c(t)}{c_0}} = \boldsymbol{\frac{a(t)}{a_0}} = \frac{\boldsymbol{1}} {\boldsymbol{1} + \boldsymbol{z(t)}} \,\,,\eqno(5.8)
	\]
	
	\medskip
	我们采用这一简单形式来构建 $c(t)$ 模型，理由如下：
	\begin{itemize}
		\item 它具有紧凑的数学形式，且抓住了宇宙红移的主要演化趋势。
		
		\item 由标准 $\Lambda$CDM \texttt{Planck18} 数据导出的宇宙时间与膨胀历史 $z(t^\Delta)$，本身就是在光速恒定的假设下从理论上推断出来的。
		
		\item 迄今为止，尚无直接观测证据揭示：若 $c$ 确实随时间变化，光速在整个宇宙历史中究竟如何演化。
		
		\item 采用变光速之后，我们的理论框架与恒光速的标准宇宙学存在根本区别。
		
		\item 光滑的 S 型函数是由红移 $z(t)$ 数据作出的合理外推；其良好的数学性质也将便于我们构建玩具模型并开展理论分析。
	\end{itemize}
	\bigskip
	\medskip
	
	
	
	{\bf 增长的光速 $\boldsymbol{c(t)}$ 模型}\,\,\,
	我们的模型假设，光速从早期宇宙到今日持续上升，并且在未来仍将缓慢增加。在这一框架中，我们把今日的宇宙时间设为 $t = 0$，而更早的宇宙时间取更大的负值。这一时间约定与前述红移模型所采用的变量 $t^\Delta$ 不同。
	\medskip
	
	据此，我们用一个广义 S 型函数来描述演化中的 $c(t)$，它具有关于其拐点的局部对称性：
	\[
	\frac{c(t)}{c_0} = \frac{1}{1 + e^{-k(t+\beta)}} \,\,.\eqno(5.9)
	\]
	
	\medskip
	此处，$k$ 与 $\beta$ 是待定的两个正参数。在极早期宇宙中，当 $t \to -\infty$ 时，指数项 $e^{-k(t+\beta)} \to +\infty$，这意味着当时的光速远低于其当前值 $c_0$。
	
	
	\bigskip
	为使函数在今日（$t=0$）归一化，即强制满足条件 $c(0)/c_0 = 1$，我们在分子上加入一个常数因子 $(1 + e^{-k\beta})$。由此得到最终的归一化公式：
	\medskip
	\[
	\boldsymbol{C(t)} = \frac{c(t)}{c_0} = \frac{1 + e^{-k\beta}}{1 + e^{-k(t+\beta)}} \,\,.\eqno(5.10)
	\]
	\bigskip
	\medskip
	
	
	{\bf $\boldsymbol{C(t)}$ 的积分与导数}\,\,\,
	为评估跨越任意宇宙时期的有效光行距离，我们首先给出归一化光速函数式 (5.10) 在一般时间区间 $[t_1, t_2]$ 上的定积分。该积分对于把我们的模型与观测宇宙时标相比较至关重要。
	
	\medskip
	引入一个常数以简化式 (5.10) 中的归一化曲线：
	\[
	J = 1 + e^{-k\beta}\,\,,\eqno(5.11)
	\]
	
	于是
	\[
	\frac{c(t)}{c_0} = \frac{J}{1 + e^{-k(t+\beta)}}\,\,.\eqno(5.12)
	\]
	\medskip
	
	在 $[t_1, t_2]$ 上的\textbf{定积分}为：
	\medskip
	\[
	D(t_1,t_2) = \int_{t_1}^{t_2} \frac{c(t)}{c_0}\,\mathrm{d}t
	= \frac{J}{k}
	\Big[ k(t+\beta) + \ln\big(1+e^{-k(t+\beta)}\big) \Big]
	\Bigg|_{t_1}^{t_2}\,.\eqno(5.13)
	\]
	\medskip
	
	
	我们进一步导出一阶与二阶导数，它们分别描述演化光速的变化率与加速度。
	\medskip
	
	\textbf{一阶导数}（光速变化率）的形式为：
	\medskip
	\[
	\dot{C}(t) = \frac{d}{dt}\left( \frac{c(t)}{c_0} \right)
	=
	\frac{k\,J\,e^{-k(t+\beta)}}
	{\Big(1 + e^{-k(t+\beta)}\Big)^2}\,\,.\eqno(5.14)
	\]
	\medskip
	
	
	至于\textbf{二阶导数}（光速变化的加速度），我们有：
	\medskip
	\[
	\ddot{C}(t) = \frac{d^2}{dt^2}\left( \frac{c(t)}{c_0} \right)
	=
	\frac{k^2\,J\,e^{-k(t+\beta)}\,\Big(e^{-k(t+\beta)} - 1\Big)}
	{\Big(1 + e^{-k(t+\beta)}\Big)^3}\,\,.\eqno(5.15)
	\]
	
	\medskip
	拐点记为 $t_{\text{inf}}$：
	\[
	t_{\text{inf}} = -\beta\,\,.\eqno(5.16)
	\]
	
	这些表达式对一切容许的正值 $k$、$\beta$ 均成立。
	
	
	
	\medskip
	\bigskip
	{\bf $\boldsymbol{C(t)}$ 的左极大曲率点}\,\,\,
	参照式 (5.14) 与式 (5.15) 所定义的 $C(t)$ 的一阶与二阶导数，我们现在导出归一化光速曲线的曲率函数。
	\medskip
	
	按照式 (5.2) 引入的曲率标准定义，曲率 $\kappa(t)$ 因分子含绝对值而恒为非负。因此，$\kappa(t)$ 的一阶导数等于零的点，对应一个局部极大曲率。
	\medskip
	
	把导数代入曲率公式，即得 $\kappa(t)$ 的显式表达式：
	\medskip
	\[
	\kappa(t) = \frac{\bm{\left|}\ddot{C}(t)\bm{\right|}}{\big[1 + (\dot{C}(t))^2\big]^{3/2}}
	= \frac{
		\bm{\left|}
		\displaystyle
		\frac{k^2 J e^{-k(t+\beta)} \left(e^{-k(t+\beta)} - 1\right)}
		{\left(1 + e^{-k(t+\beta)}\right)^3}
		\bm{\right|}
	}{
		\left[1 + \left(
		\displaystyle
		\frac{k J e^{-k(t+\beta)}}
		{\left(1 + e^{-k(t+\beta)}\right)^2}
		\right)^2\right]^{3/2}
	}\,\,,\eqno(5.17)
	\]
	\medskip
	
	
	我们不直接求解这一复杂表达式，而是采用一种近似方案。参数 $k$ 与 $J$ 会轻微破坏 S 型曲线严格的平移不变性，但在 $kJ\approx 1$ 的条件下，有效的近似不变性依然成立。于是左极大曲率点可近似为
	\[
	{t}_{\text{leftmax}} \approx -\beta - \frac{\ln(2+\sqrt{3})}{k}\,\,.\eqno(5.18)
	\]
	
	
	\medskip
	\bigskip
	{\bf 借助宇宙正午锚定求解参数}\,\,\,
	由于宇宙光速的真实演化曲线无法直接观测，我们选择最具物理意义的宇宙时期作为锚点，来求解函数式 (5.10) 中的参数 $\boldsymbol{k}$ 与 $\boldsymbol{\beta}$。
	
	\medskip
	宇宙正午对应宇宙恒星形成的高峰时期，位于 $t^\star \approx -11.4\ \mathrm{Gyr}$；它标志着宇宙演化节律中生长最旺盛的阶段，与模型演化曲线的曲率极大值相合，是最具代表性的时间锚点。
	
	\medskip
	以 $\boldsymbol{t_w}$ 记宇宙正午时刻，我们施加三条独立的物理约束：
	\begin{itemize}
		\item 由式 (5.13) 给出的有效光行时积分满足 $D(t_w,0)=11.4$$\mathrm{Gyr}$，即等于从宇宙正午到今日的总积分光传播距离。
		
		\item 宇宙正午处的归一化光速为 $C(t_w)=1/3.7$，这一约束来自实测红移 $z=2.7$ 与式 (5.10) 的模型定义。
		
		\item 基于物理上的自洽性，宇宙正午与式 (5.18) 给出的左曲率极大值点重合，即 $t_w = t_{\text{leftmax}}$，此处 $c(t)$ 的加速度变化率达到峰值。
	\end{itemize}
	
	\medskip
	这三条约束构成关于三个未知量 $\boldsymbol{k}$、$\boldsymbol{\beta}$ 与 $\boldsymbol{t_w}$ 的封闭代数方程组，显式写为
	\[
	\left\{
	\begin{aligned}
		\displaystyle \quad \frac{J}{k}\Big[ k(t+\beta) + \ln\big(1+e^{-k(t+\beta)}\big) \Big]\Bigg|_{t_w}^{0} \enspace &= \enspace 11.4\\[6pt]
		\displaystyle \quad \frac{1 + e^{-k\beta}}{1 + e^{-k(t_w+\beta)}} \enspace &= \enspace \frac{1}{3.7}\\[6pt]
		\displaystyle \quad -\beta - \frac{\ln(2+\sqrt{3})}{k} \enspace &= \enspace t_w 
	\end{aligned}
	\right.\eqno(5.19)
	\]
	
	
	
	解得的参数为
	\[
	\boldsymbol{k} = 0.1442, \quad \boldsymbol{\beta} = 8.854, \quad \boldsymbol{t_w} = -17.98\,\mathrm{Gyr}, \quad \boldsymbol{J} = 1.2789. \eqno(5.20)
	\]
	
	\medskip
	至此，这个紧凑的增长-$c(t)$ 模型已完全建立。下面给出其演化曲线。
					
	\begin{figure}[H]
		\centering
		\begin{tikzpicture}
			\begin{axis}[
				width=14cm, height=7cm,
				axis lines = box,
				xlabel = {宇宙时间 \(t\) [Gyr]},
				ylabel = {S 型增长-\(c(t)\)},
				xmin=-68, xmax=40,
				ymin=-0.06, ymax=1.4,
				xtick = {-70, -60, -50, -40, -30, -20, -10, 0, 10, 20, 30, 40},
				ytick = {0.0, 0.2, 0.4, 0.6, 0.8, 1.0, 1.2, 1.4},
				tick label style={font=\small},
				label style={font=\small},
				smooth,
				every tick/.style={thick, black},
				]
				
				\addplot[color=red, line width=1.6pt, solid, domain=-59.1:0, samples=600]
				{ (1 + exp(-0.1442 * 8.854)) / (1 + exp(-0.1442*(x + 8.854))) };
				
				\addplot[color=purple!80!black, line width=1.6pt, dashed, domain=0:35, samples=600]
				{ (1 + exp(-0.1442 * 8.854)) / (1 + exp(-0.1442*(x + 8.854))) };
				
								
				% 参考线 z_future = 未来极大值, 1.2789 
				\draw[purple!80!black, line width=0.5pt] (axis cs:-70, 1.2789) -- (axis cs:35, 1.2789);	
				% 文字标注
				\node[purple!80!black, font=\scriptsize] at (axis cs:-12, 1.2289) {$\boldsymbol{J = 1.2789}$};

				
				% 参考线 z_symmetric = 对称线, 0.63945 
				\draw[purple!80!black, line width=0.5pt] (axis cs:-14, 0.63945) -- (axis cs:35, 0.63945);
				\draw[purple!80!black, line width=1.2pt] (axis cs:-8.854, 0.43945) -- (axis cs:-8.854, 0.90);
				% 文字标注
				\node[purple!80!black, font=\scriptsize] at (axis cs:17, 0.58945) {$\boldsymbol{0.5*J = 0.63945}$};
				\node[purple!80!black, font=\scriptsize, anchor=south west, rotate=90] at (axis cs:-8.854, 0.70) {$\boldsymbol{-\beta}$};

				
				% 参考线 z_now = 今日, -0Gyr, 1.0 
				\draw[blue!80!black, line width=0.8pt] (axis cs:0, -0.08) -- (axis cs:0, 1.1);
				\draw[blue!80!black, line width=0.3pt] (axis cs:-70, 1.00) -- (axis cs:3, 1.00);
				% 文字标注
				\node[blue!80!black, font=\scriptsize, anchor=south west, rotate=90] at (axis cs:-0.05, 0.03) {今日};

				% 参考线 z_3 = 太阳系, -7.55Gyr, 0.70 
				\draw[blue!80!black, line width=0.8pt] (axis cs:-7.55, -0.08) -- (axis cs:-7.55, 0.9);
				\draw[blue!80!black, line width=0.3pt] (axis cs:-70, 0.70) -- (axis cs:-5, 0.70);
				% 文字标注
				\node[blue!80!black, font=\scriptsize, anchor=south west, rotate=90] at (axis cs:-7.60, 0.03) {${z_3}$\,\,\,太阳系};
				
				% 参考线 z_2 = 宇宙正午, -36.8Gyr, 0.04 
				\draw[blue!80!black, line width=0.8pt] (axis cs:-17.83, -0.08) -- (axis cs:-17.83, 0.80);
				\draw[blue!80!black, line width=0.3pt] (axis cs:-70, 0.27) -- (axis cs:-15, 0.27);
				% 文字标注
				\node[blue!80!black, font=\scriptsize, anchor=south west, rotate=90] at (axis cs:-17.90, 0.03) {${z_2}$\,\,\,宇宙正午};
				
				% 参考线 z_1 = 早期星系, -32.7Gyr, 0.04 
				\draw[blue!80!black, line width=0.8pt] (axis cs:-32.7, -0.08) -- (axis cs:-32.7, 0.80);
				\draw[blue!80!black, line width=0.3pt] (axis cs:-70, 0.04) -- (axis cs:-30, 0.04);
				% 文字标注
				\node[blue!80!black, font=\scriptsize, anchor=south west, rotate=90] at (axis cs:-32.8, 0.03) {${z_1}$\,\,\,早期星系};
				
				% 参考线 z_0 = CMB 退耦, -59.1Gyr, 0.092%
				\draw[blue!80!black, line width=0.8pt] (axis cs:-59.1, -0.08) -- (axis cs:-59.1, 0.80);
				\draw[blue!80!black, line width=0.3pt] (axis cs:-70, 0.0009) -- (axis cs:-59.1, 0.0009);
				% 文字标注
				\node[blue!80!black, font=\scriptsize, anchor=south west, rotate=90] at (axis cs:-59.2, 0.03) {${z_0}$\,\,\,CMB 退耦};
												
			\end{axis}
		\end{tikzpicture}
		
		
		\vspace{0.2cm}
		\captionsetup{width=0.85\textwidth, font=footnotesize}
		\caption{
			参数确定后的紧凑 S 型增长-$c(t)$ 模型的 $\mathrm{C(t)}$ 曲线，其中 $k=0.1442$、$\beta=8.854$、$t_w=-17.98\,\mathrm{Gyr}$、$J=1.2789$。四个基准时期 $z_0, z_1, z_2, z_3$ 所对应的宇宙时间将在第 7 节中说明。}
		\label{fig:growing_ct_model}
			
	\end{figure}
	
	
	\medskip
	\bigskip
	{\bf 关于 S 型模型曲线的讨论}\,\,\,
	由式 (5.20) 得到的模型使我们能够提炼出若干值得注意的特征，归纳如下。
	\begin{itemize}
		\item \textit{未来的渐近光速。} 当 \( t \to +35 \) Gyr 时，归一化光速 \( C(t) \) 趋于渐近值 \( J = 1.2789 \)，这意味着若这一 S 型增长-$c(t)$ 成立，光速将在未来 35 Gyr 内继续增加约 \( 27.89\% \)，达到 \( 1.2789\,c_0\)。
		
		\item \textit{大幅延展的宇宙时间线。} CMB 退耦时期 \( z_0 \) 对应于 \( t = -59.1 \) Gyr，远超标准 \(\Lambda\)CDM 的 \( -13.8 \) Gyr。
		
		\item \textit{曲线的对称中心。} S 型函数关于其拐点 \( \boldsymbol{t} = \boldsymbol{-\beta} = -8.854 \) Gyr 对称，此处 \( C(t) = 0.5J = 0.63945 \)。
				
		\item \textit{一个耐人寻味的对称。} 值得注意的是，宇宙正午（\( t_w = -17.98 \) Gyr）位于左极大曲率点附近，而今日（\( t = 0 \)）则靠近右极大曲率点，二者满足关系 \( 2\times(-8.854) = -17.708 \approx -17.98\) Gyr。
		
		\item \textit{一点哲学反思。} 这一近乎对称的现象，出自一个仅以宇宙正午为锚点的模型，颇值得玩味。左极大值标示恒星形成与黑洞活动的高峰时期，对应宇宙更新最剧烈的阶段；右极大值则对应宇宙成熟到足以孕育智慧生命的时代。两个极值竟如此自然地与宇宙史上这两个决定性时刻相合，似乎并非纯属偶然；它暗示着一种或可称之为"存在之美"的共振。
	\end{itemize}

		
	\medskip
	\textit{\textbf{注：}也可通过式 (5.6) 的改进复合模型来定义函数，从而构造增长-$c(t)$ 模型的一个变体；然而该形式在数学结构上过于繁复，故超出本工作的范围。}
	
	
	
	
	
	\section*{
		\begin{center}
			6. 超精细增长 $\boldsymbol{c(t)}$ 模型
		\end{center}
	}
	\medskip
	\noindent
	
	在标准 $\Lambda$CDM 宇宙学中，只有宇宙红移可以被直接观测。精确的宇宙膨胀历史及与之对应的宇宙时间，在实践中无法测量。
	
	\medskip
	存储在 \texttt{Planck18} 数据集中的所有宇宙时间条目，都是模型导出的量。这类计算依赖三条基本前提：{\textit{均匀且各向同性的宇宙}}、{\textit{总物质与总能量不变}}、{\textit{光速恒定}}。宇宙年龄则是在这些限制下由弗里德曼方程求解得到的。
	
	\medskip
	在我们 $G(t)\,c(t)^2 = P$ 的公设之下，$G(t)$ 随时间变化，这意味着 $\Lambda$CDM 中所用的恒星质量估计方法本身含有内在偏差。这一问题的详细讨论留待第 12 节。
	
	\medskip
	尽管受限于预设的理论假设以及模型推导中不可避免的不确定性，\texttt{Planck18} 数据仍是最好的可及参考。因此，我们以该数据集为基础构建超精细增长-$c(t)$ 模型。
	
	
	
	\medskip
	\bigskip
	{\bf 数据提取}\,\,\,
	\texttt{Planck18} 标准红移—宇宙时间数据集 $(z,\, {t}^\vartriangle)$ 以 float64 数值精度从 Python \texttt{astropy.cosmology} 库中提取。
	
	
	\medskip
	宇宙总年龄约为 $13.8\ \mathrm{Gyr}$。我们以 $10^4$ 年（10 kyr）的固定步长对整个宇宙时间线作均匀采样，由此得到约 $13.8\times10^{5}$ 个采样数据点。
	
	\medskip
	采样自当前宇宙时间 ${t}^\vartriangle=13.79069_{\boldsymbol\cdots}\mathrm{Gyr}$ 开始，此后每个采样点依次减少 $10^4$ 年：第二、第三个采样时刻分别为 ${t}^\vartriangle=13.79068_{\boldsymbol\cdots}$ 与 ${t}^\vartriangle=13.79067_{\boldsymbol\cdots}$。
	
	\medskip
	CMB 时期的红移为 $z\simeq 1090$，它定义了观测的最上边界。为此，采样在 ${t}^\vartriangle=0.00037_{\boldsymbol\cdots}$ 与 ${t}^\vartriangle=0.00036_{\boldsymbol\cdots}$ 处截断，分别对应 $z=1083.07$ 与 $z=1100.61$，即大爆炸之后 370 kyr 与 360 kyr，符合这一物理极限。
	
	\medskip
	最终共提取 $N=1\,379\,024$ 组等时间间隔的原始时间—红移数据对。
	\medskip
	
	
	\bigskip
	{\bf 区间距离 $\boldsymbol{yr_i}$ 的定义}\,\,\,
	按分段平均规则，各时间区间的分段时距定义为
	$$
	\mathit{yr}_i=1+\frac12\left(z_i+z_{i+1}\right) \,\,,\eqno(6.1)
	$$
	
	其中 $i\in (0,\,1379022)$，$z_i$ 表示第 $i$ 个采样点处的红移。所有 $yr_i$ 值在全部有效区间上依次计算。
	
	\medskip
	每个分段 $yr_i$ 代表在增长-$c(t)$ 模型下导出的时间标度因子。例如，$yr_i=1.5$ 对应由 $z_i$ 与 $z_{i+1}$ 得出的平均红移 $\bar z=0.5$；在增长-$c(t)$ 框架内，原本 $10^4$ 年的采样间隔被折算为 $1.5\times 10^4$ 年。
	\medskip
	

	\bigskip
	{\bf 累积求和 $\boldsymbol{\mathrm{YR}_i}$}\,\,\,
	从 $i=0$ 起推进至 $i=1379022$，我们按式 (6.2) 累加序列 $yr_i$，以计算增长-$c(t)$ 模型下的宇宙总时标。为抑制数百万次连续相加所带来的浮点舍入误差，采用 Kahan 求和算法作残差补偿。
	$$
	\mathrm{YR}_0 = yr_0 \,\,,\quad \mathrm{YR}_i = \mathrm{YR}_{i-1} + yr_i \,\,.\eqno(6.2)
	$$
	
	\medskip
	所得的 $\mathrm{YR}_i$ 数组以 $10^4$ 年为单位，量化了增长-$c(t)$ 时标。每个 $\mathrm{YR}_i$ 乘以 $-1$ 以表示负的宇宙时间。CMB 时期的形成点对应于 $-45.2\ \mathrm{Gyr}$。
	\medskip
	
	
	\bigskip
	{\bf $\boldsymbol{C(t)}$ 的归一化与 $\boldsymbol{\mathrm{YR}_i}$ 的降采样}\,\,\,
	归一化的 $C(t)$ 由采样红移序列按 $C(t) = 1/(1+z_t)$ 给出，由此产生含 $1\,379\,023$ 条记录的全分辨率 $C(t)$ 数据。
	
	\medskip
	为便于在本文中呈现，我们按每 1000 个原始采样保留一个数据的规则降采样。该过程给出 1381 组离散的增长-$c(t)$ 及其对应时间 $\mathrm{YR}_i$，相邻数据之间的有效间距为 $t^\Delta = 10^7\ \mathrm{years}$。
	
	\medskip
	该数据集即所提出的超精细增长-$c(t)$ 模型的核心数据库。其对应曲线如下图。
	\bigskip
	
	\begin{figure}[H]
		\centering
		\begin{tikzpicture}
			\begin{axis}[
				width=14cm, height=7cm,
				axis lines = box,
				xlabel = {宇宙时间 \(t\) [Gyr]},
				ylabel = {增长-\(c(t)\)},
				xmin=-50, xmax=5,
				ymin=-0.05, ymax=1.2,
				xtick = {-50, -45, -40, -35, -30, -25, -20, -15, -10, -5, 0, 5},
				ytick = {0.0, 0.2, 0.4, 0.6, 0.8, 1.0, 1.2},
				tick label style={font=\small},
				label style={font=\small},
				every tick/.style={thick, black},
				]
				
				% 读取 1381 组实测数据，红色小方块，不连线
				\addplot[
				color=red,
				only marks,
				mark=square,
				mark size=0.8pt,
				] table [col sep=comma, x=X, y=Y] {2026-growing-ct-fine-model-1381.csv};
				
				
				% 参考线 z_3 = 太阳系, -5.57Gyr, 0.70 
				\draw[blue!80!black, line width=0.8pt] (axis cs:-5.57, 0) -- (axis cs:-5.57, 1.2);
				\draw[blue!80!black, line width=0.4pt] (axis cs:-60, 0.70) -- (axis cs:-3, 0.70);
				% 文字标注
				\node[blue!80!black, font=\scriptsize, anchor=south west, rotate=90] at (axis cs:-5.68, 0.50) {${z_3}$\,\,\,太阳系};
				
				% 参考线 z_2 = 宇宙正午, -36.8Gyr, 0.04 
				\draw[blue!80!black, line width=0.8pt] (axis cs:-20, 0) -- (axis cs:-20, 1.2);
				\draw[blue!80!black, line width=0.4pt] (axis cs:-60, 0.27) -- (axis cs:-16, 0.27);
				% 文字标注
				\node[blue!80!black, font=\scriptsize, anchor=south west, rotate=90] at (axis cs:-20.1, 0.29) {${z_2}$\,\,\,宇宙正午};
				
				% 参考线 z_1 = 早期星系, -36.8Gyr, 0.04 
				\draw[blue!80!black, line width=0.8pt] (axis cs:-36.8, 0) -- (axis cs:-36.8, 1.2);
				\draw[blue!80!black, line width=0.4pt] (axis cs:-60, 0.04) -- (axis cs:-32, 0.04);
				% 文字标注
				\node[blue!80!black, font=\scriptsize, anchor=south west, rotate=90] at (axis cs:-36.9, 0.13) {${z_1}$\,\,\,早期星系};
				
				% 参考线 z_0 = CMB 退耦, -45.2Gyr, 0.092%
				\draw[blue!80!black, line width=0.8pt] (axis cs:-45.2, 0) -- (axis cs:-45.2, 1.2);
				% 文字标注
				\node[blue!80!black, font=\scriptsize, anchor=south west, rotate=90] at (axis cs:-45.3, 0.03) {${z_0}$\,\,\,CMB 退耦};
				
			\end{axis}
		\end{tikzpicture}
		
		\vspace{0.2cm}
		\captionsetup{width=0.85\textwidth, font=footnotesize}
		\caption{
			由 \texttt{Planck18} 数据导出的超精细增长-$c(t)$ 模型的 $\mathrm{C(t)}$ 曲线。
			四个基准时期 $z_0, z_1, z_2, z_3$ 所对应的宇宙时间将在下一节中说明。}
		\label{fig:growing_ct_fine_model}
		
	\end{figure}
	
	
	
	
	
	
	
	\section*{
		\begin{center}
			7. 增长-\(\boldsymbol{c(t)}\) 宇宙的 \\
			不同物理图景
		\end{center}
	}
	\medskip
	\noindent
	
	本节描绘了在不同宇宙时期下，宇宙演化所呈现的一系列截然不同的物理图景。
	
	\medskip
	宇宙起源时标被延展至 $[-59.1,\,-45.2]\,$$\mathrm{Gyr}$；而在 CMB 退耦时期，光速 \(c(t)\) 等于当前光速 \(c_0\) 的 \(1/1100\)，玻尔半径 \(a_\text{B}\) 仅为当前值的 \(1/1100\)，氢原子的电离温度 ${T_\mathrm{ion}(t)}$ 也仅达 82.64$\,\mathrm{K}$。
	
	\medskip
	经典基本方程中的许多物理量都随 \(c(t)\) 一同演化，以下各小节将依次阐述每个主题。
	
	\medskip
	本工作所用的全部基本物理常数当前值（例如玻尔半径 $a_{\scriptscriptstyle\mathrm{B},0}$、约化普朗克常数 $\hbar_0$、质子直径 $d_p$）均取自标准参考资料~\cite{Tipler:2008}。

	\bigskip
	{\bf 宇宙起源时标}\,\,\,
	CMB 退耦时期标志着可观测宇宙的起点，其对应红移满足 $1+z=1100$。
	
	\medskip
	在式 (5.10) 所定义的紧凑增长-$c(t)$ 模型下，$C(t)=1/1100$ 给出 $t=-59.1\,\mathrm{Gyr}$；对超精细增长-$c(t)$ 模型而言，在红移 $z=1083.07$ 的早期 CMB 阶段，宇宙时间为 $t=-45.2\,\mathrm{Gyr}$。这些结果暗示一个平缓演化的宇宙，而非剧烈的大爆炸图景~\cite{Weinberg:2008}。
	
	\medskip
	一项初步评估表明，QGP 阈值时期（$\boldsymbol{z_{\star}}$）与 CMB 退耦时期（$\boldsymbol{z_{1}}$）之间的时距，超过了标准 $\Lambda$CDM 模型给出的 360–370\,$\mathrm{kyr}$。
	
	\medskip
	宇宙的终极起源被设定为{\textit{零光速零温度}}（\textit{ZLZT}）原初态，其中光速与温度都趋于零。我们推断，若在 QGP 阈值时期（$\boldsymbol{z_{\star}}$）宇宙光速降到 42\,$\mathrm{km/s}$（21$\,\mathrm{km/s}$ 或 \allowbreak 14 $\mathrm{km/s}$）以下，而环境温度维持在 2.75$\,\mathrm{K}$ 附近，则一切强子结构都将完全解离为自由夸克与胶子，形成一个没有强子的纯夸克—胶子宇宙。详细讨论见第 9 节。
	
	
	\bigskip
	{\bf 宇宙正午与 $\boldsymbol{C(t)}$}\,\,\,
	依据式 (5.20) 的紧凑增长-$c(t)$ 模型，宇宙正午出现在 $t_w = -17.98\,\mathrm{Gyr}$。对于式 (6.2) 所述的超精细增长-$c(t)$ 模型，该时期落在 $t_w = -20.21\,\mathrm{Gyr}$。两个时间标记都与标准 $\Lambda$CDM 模型预言的 $t^\star = -11.4\,\mathrm{Gyr}$ 相去甚远。
		
	\medskip
	在宇宙正午，宇宙光速达到当前值 $c_0$ 的 $27.03\%$，相当于 $1/3.7\,c_0$。按式 (1.11) 导出的标度关系 $r_s \propto c^{-4}$，该时期的史瓦西半径比今日大 $3.7^4 = 187.4$ 倍，这解释了为何在这一宇宙时代黑洞能够高效形成。

		
	\bigskip
	{\bf 太阳系形成}\,\,\,
	对太阳系形成时期（$\boldsymbol{z_3}$），\texttt{Planck18} 数据集在红移 $z=0.43$ 处给出 $t^\star = -4.67\,\mathrm{Gyr}$。
	
	\medskip
	紧凑模型与超精细增长-$c(t)$ 模型分别给出 $-7.55\,\mathrm{Gyr}$ 与 $-5.57\,\mathrm{Gyr}$。我们认为后一结果更可靠，因为超精细模型是用完整的 \texttt{Planck18} 数据集标定的，而紧凑模型只依赖单个数据点。该时期的相对光速为 $69.93\%$。
	
	\medskip
	早期星系形成时期（$\boldsymbol{z_1}$）的数据汇总于下表，不再进一步讨论。
	
	\bigskip
	\begin{table}[H]
		\centering
		\captionsetup{skip=2pt, font=small}
		\caption{基准时期 $\boldsymbol{z_i}$ 的 $\boldsymbol{C(t)}$ 与 $\boldsymbol{t}$}
		\footnotesize
		\renewcommand{\arraystretch}{1.50}
		\begin{tabular}{>{\centering\arraybackslash}p{2.5cm}|>{\centering\arraybackslash}p{2cm}>{\centering\arraybackslash}p{2cm}|>{\centering\arraybackslash}p{2cm}>{\centering\arraybackslash}p{2cm}}
			\toprule
			\textbf{参数} & \textbf{$\boldsymbol{z_0}$} & \textbf{$\boldsymbol{z_1}$} & \textbf{$\boldsymbol{z_2}$} & \textbf{$\boldsymbol{z_3}$} \\
			\midrule
			$\boldsymbol{C(t)}$                     & $1/1100$    & $1/25$    & $1/3.7$   & $1/1.43$   \\
			$\boldsymbol{C(t)\,\,\%}$               & $0.09\%$    & $4.00\%$  & $27.03\%$ & $69.93\%$  \\
			$\boldsymbol{t}\mathsf{_{compact}}$     & $-59.1$     & $-32.7$   & $-17.83$  & $-7.55$    \\
			$\boldsymbol{t}\mathsf{_{fine}}$        & $-45.2$     & $-36.6$   & $-20.21$  & $-5.57$    \\
			$\boldsymbol{t^\star}$                  & $-13.8$     & $-13.6$   & $-11.40$  & $-4.67$    \\
			\bottomrule
		\end{tabular}
	\end{table}
	
	
	
	\medskip
	{\bf 量纲分析}\,\,\,
	物理学定义了七个基本量纲，简记为 \texttt{LMTI$\Theta$NJ}。它们分别代表长度（\texttt{L}）、质量（\texttt{M}）、时间（\texttt{T}）、电流（\texttt{I}）、热力学温度（\texttt{$\Theta$}）、物质的量（\texttt{N}）与发光强度（\texttt{J}）。
	
	\medskip
	本研究建立在第 2 节所列的两条核心公设之上。宏观公设为 ${G(t)}\tinysp {c(t)^2} = \boldsymbol{P}$，它表明 $G(t) \propto c(t)^{-2}$。微观公设则指出，精细结构常数 $\boldsymbol{\alpha}$ 在宇宙演化过程中普适不变。
	
	\medskip
	基于这些公设，第 1 节中的一部分物理量随 $c(t)$ 演化。我们导出它们的量纲以及与 $c(t)$ 的幂律关系，汇总于下表。
	
	\medskip
	\begin{table}[H]
		\centering
		\captionsetup{skip=2pt, font=small}
		\caption{量纲分析及与 $\boldsymbol{c(t)}$ 的幂律标度}
		\footnotesize
		\renewcommand{\arraystretch}{1.50}
		\begin{tabular}{>{\centering\arraybackslash}p{2.5cm}|>{\hspace{0.5cm}\raggedright\arraybackslash}p{4.2cm}|>{\hspace{0.5cm}\centering\arraybackslash}p{5.5cm}}
			\toprule
			\textbf{演化量} & \textbf{量纲形式} & \textbf{与 $\boldsymbol{c(t)}$ 的幂律标度} \\
			\midrule
			$G(t)$                                  & $[G]$ : $\mathrm{L}^{3}\, \mathrm{M}^{-1}\, \mathrm{T}^{-2}$                                & ${\propto c^{-2}}$ \\
			$\mu_0(t)$                              & $[\mu_0]$ : $\mathrm{L}^{1}\, \mathrm{M}\, \mathrm{T}^{-2}\, \mathrm{I}^{-2}$               & ${\propto c^{+1}}$ \\
			$\varepsilon_0(t)$                      & $[\varepsilon_0]$ : $\mathrm{L}^{-3}\, \mathrm{M}^{-1}\, \mathrm{T}^{4}\, \mathrm{I}^{2}$   & ${\propto c^{-3}}$ \\
			$\hbar(t)$                              & $[\hbar]$ : $\mathrm{L}^{2}\, \mathrm{M}\, \mathrm{T}^{-1}$                                 & ${\propto c^{+2}}$ \\
			$a_{\scriptscriptstyle\mathrm{B}}(t)$   & $[a_{\mathrm{B}}]$ : $\mathrm{L}^{1}$                                                       & ${\propto c^{+1}}$ \\
			$k_e(t)$                                & $[k_e]$ : $\mathrm{L}^{3}\mathrm{M}\, \mathrm{T}^{-4}\, \mathrm{I}^{-2}$                    & ${\propto c^{+3}}$ \\
			\bottomrule
		\end{tabular}
	\end{table}
	
	\medskip
	其余七个量由上述结果导出。它们的量纲形式及与 $c(t)$ 的幂律标度列于下表。
	
	\medskip
	\begin{table}[H]
		\centering
		\captionsetup{skip=2pt, font=small}
		\caption{其余四个量与 $\boldsymbol{c(t)}$ 的幂律标度}
		\footnotesize
		\renewcommand{\arraystretch}{1.50}
		\begin{tabular}{>{\centering\arraybackslash}p{2.5cm}|>{\centering\arraybackslash}p{1.8cm}|>{\hspace{0.5cm}\raggedright\arraybackslash}p{8.2cm}}
			\toprule
			\textbf{演化量} & \textbf{参考式} & \textbf{与 $\boldsymbol{c(t)}$ 的幂律标度} \\
			\midrule
			$m_{\scriptscriptstyle\mathrm{P}}(t)$    &$\mathrm{Eq.}(1.3)$      & ${\hbar \propto c^{+2}},\enspace {c^{+1}\propto c^{+1}},\enspace {G^{-1}\propto c^{+2}}  \quad \quad \Rightarrow \quad {\propto \boldsymbol{c^{\frac{5}{2}}}  }$  \\
			$\ell_{\scriptscriptstyle\mathrm{P}}(t)$ &$\mathrm{Eq.}(1.4)$      & ${\hbar \propto c^{+2}},\enspace {c^{-3}\propto c^{-3}},\enspace {G^{+1}\propto c^{-2}}  \quad \quad \Rightarrow \quad {\propto \boldsymbol{c^{-\frac{3}{2}}} }$  \\
			$\rho_{\scriptscriptstyle\mathrm{P}}(t)$ &$\mathrm{Eq.}(1.5)$      & ${\hbar^{-1} \propto c^{-2}},\enspace {c^{+5} \propto c^{+5}},\enspace {G^{-2} \propto c^{+4}} \quad \Rightarrow \quad {\propto \boldsymbol{c^{7}}      }$        \\
			$F_{\scriptscriptstyle\mathrm{P}}(t)$ &$\mathrm{Eq.}(1.5)$         & ${c^{+4} \propto c^{+4}},\enspace {G^{-1} \propto c^{+2}} \quad \Rightarrow \quad {\propto \boldsymbol{c^{6}}      }$        \\
			$T_{\mathrm{ion}}(t)$  &$\mathrm{Eq.}(1.8)$    & ${\quad \Rightarrow \quad {\propto \boldsymbol{c^{2}}}}     $  \\
			$r_{\mathrm{s}}(t)$ &$\mathrm{Eq.}(1.11)$      & ${\quad \Rightarrow \quad {\propto \boldsymbol{c^{-4}}}}    $  \\
			$N_{\mathrm{Dirac,1}}(t)$ &$\mathrm{Eq.}(1.16)$    & $ {G \propto c^{-2}},\enspace \varepsilon_0 \propto c^{-3} \quad \Rightarrow \quad {\propto \boldsymbol{c^{5}}      }$        \\
			\bottomrule
		\end{tabular}
	\end{table}
	
	
	\medskip
	\textbf{$\boldsymbol{\varepsilon_0(t)}$、$\boldsymbol{\mu_0(t)}$ 的演化}\,\,\,
	$\varepsilon_0(t)$ 与 $\mu_0(t)$ 是支配 $c(t)$ 演化的核心量。依 $\varepsilon_0(t) \propto c(t)^{-3}$ 与 $\mu_0(t) \propto c(t)$，随 $c(t)$ 上升，这两个参数分别减小与增大，并相互制约，与麦克斯韦方程组的核心形式相符，如式 (7.1) 所示。
	\[
	\mu_0(t)\,\varepsilon_0(t)\, c(t)^2 = 1\,\,. \eqno(7.1)
	\]
	
	\medskip
	$\mu_0(t)$、$\varepsilon_0(t)$ 与 $c(t)$ 这三个量由{\textit{光—电—磁同步宇宙暗场}}（\textit{LEM-SCDF}）所支配，后文将对此作完整阐述。
	
	
	\medskip
	\bigskip
	\textbf{玻尔半径 $\boldsymbol{a_{\scriptscriptstyle\mathrm{B}}(t)}$ 的演化}\,\,\,
	在我们的框架中，氢的玻尔半径 $a_{\scriptscriptstyle\mathrm{B}}(t)$ 与 $c(t)$ 成线性标度，这与惯常假设其为常数不同。在 $1+z=25$ 的时期，其值等于今日值的 $1/25$。在 $1+z=1100$ 的 CMB 时期，玻尔半径缩小至当前量级的 $1/1100$。
	
	\medskip
	今日的玻尔半径为 $a_{\scriptscriptstyle\mathrm{B},0}=5.292\times10^{-11}\,\mathrm{m}$，质子直径为 $d_p=1.6815\times10^{-15}\,\mathrm{m}$，二者今日之比约为 $3.15\times10^4$。在 $1+z=1100$ 时，该比值降至约 $28.6$，表明早期宇宙中电子围绕质子的轨道极其狭窄。这一比值是维持中性氢所需的最低尺度。在宇宙演化到达这一阶段之前，轨道半径过小，无法形成稳定的氢原子。
	
	\medskip
	这一图景在既有研究中极少被探讨。当玻尔半径仅为质子直径的数十倍时，多电子原子与复杂分子键几乎无法形成，因为电子壳层被压缩到彼此重叠的构型。此外，复合过程中从等离子体到中性氢的转变既脆弱又可逆；新形成的原子结合松散，很可能被再次电离。
	
	\medskip
	原子的轨道尺度还进一步支配着化学元素的丰度。在 CMB 前后的时代，严重压缩的原子结构只允许以氢为主的天体凝聚，几乎没有复杂的重元素。在早期宇宙环境中，简单的化学组成在很长时期内占主导。
	
	
	\medskip
	\bigskip
	\textbf{库仑参数 $\boldsymbol{k_e(t)}$ 的演化}\,\,\,
	如表 3 所分析，库仑参数 $k_e(t)$ 与宇宙光速 $c(t)$ 的立方成标度。在红移 $1+z=25$ 处，$k_e(t)$ 为今日值的 $1/25^3$；而在 CMB 时期（$1+z=1100$），它缩小至当前量级的 $1/1100^3$。早期宇宙中急剧减弱的静电力导致微观尺度上的原子成键十分脆弱，这与上文讨论的氢原子演化特征一致。
	
	
	\medskip
	\bigskip
	\textbf{普朗克量 $\boldsymbol{\hbar(t)}$、$\boldsymbol{m_{\scriptscriptstyle\mathrm{P}}(t)}$、$\boldsymbol{\ell_{\scriptscriptstyle\mathrm{P}}(t)}$、$\boldsymbol{F_{\scriptscriptstyle\mathrm{P}}(t)}$ 与 $\boldsymbol{\rho_{\scriptscriptstyle\mathrm{P}}(t)}$ 的演化}\,\,\,
	如表 3 与表 4 所分析，约化普朗克常数 $\hbar(t) \propto c(t)^{2}$，普朗克质量 $m_{\scriptscriptstyle\mathrm{P}}(t) \propto c(t)^{5/2}$，普朗克长度 $\ell_{\scriptscriptstyle\mathrm{P}}(t)\propto c(t)^{-3/2}$，普朗克力 $F_{\scriptscriptstyle\mathrm{P}}(t)\propto c(t)^{6}$，而常规普朗克密度 $\rho_{\scriptscriptstyle\mathrm{P}}(t)$ 则按 $c(t)^7$ 成比例演化。
	
	\medskip
	在今日宇宙中，普朗克质量达到 $m_{\scriptscriptstyle\mathrm{P,0}} \approx 2.2\times 10^{-8}\,\mathrm{kg} \approx 22\, \mathrm{\mu g}$，对微观尺度而言极为巨大；而普朗克长度 $\ell_{\scriptscriptstyle\mathrm{P},0} \approx 1.6\times 10^{-35}\,\mathrm{m}$ 仍极其微小，普朗克密度 $\rho_{\scriptscriptstyle\mathrm{P},0} \approx 5.2\times 10^{96}\,\mathrm{kg/m^3}$。
	
	\medskip
	由于 $\rho_{\scriptscriptstyle\mathrm{P}}(t) \propto c(t)^7$，在 $1+z=1100$ 的 CMB 时期（${z_0}$），有 $1100^7 \approx 1.95 \times 10^{21}$，从而把普朗克密度降至 $\rho_{\scriptscriptstyle\mathrm{P}}(t) \approx 2.65\times 10^{75}\,\mathrm{kg/m^3}$。
	
	\medskip
	此外，在强子化后的等离子体时期（${z_{\star}}$），光速降至今日值 $c_0$ 的 $1/7137.66$，即约 42 km/s。我们得到 $7137.66^7 \approx 9.44 \times 10^{26}$，从而把普朗克密度降至 $\rho_{\scriptscriptstyle\mathrm{P}}(t) \approx 5.51\times 10^{69}\,\mathrm{kg/m^3}$。
		
	\medskip
	上述两个普朗克密度都远大于式 (1.14) 所定义的弗里德曼宇宙密度 $\rho_{\scriptscriptstyle\mathrm{F}}(t)$。后文讨论式 (8.10) 中的 ${v_{\scriptscriptstyle\mathrm{LEM}}(t)}$ 时，我们将重新回到这些量，那里会同时用到 $\rho_{\scriptscriptstyle\mathrm{P}}(t)$ 与 $\rho_{\scriptscriptstyle\mathrm{F}}(t)$。
	
	\medskip
	若消去 $G$ 并以 $P = Gc^2$ 取代之，则式 (1.3)、式 (1.4) 与式 (1.5) 中的四个普朗克量变为：
	\[
	m_P(t) = \sqrt{\frac{\hbar(t) \cdot c(t)^3}{P}}, \quad \ell_P(t) = \sqrt{\frac{\hbar(t) \cdot P}{c(t)^5}}, \quad \rho_P(t) = \frac{c(t)^9}{\hbar(t) \cdot P^2}, \quad F_{\scriptscriptstyle\mathrm{P}}(t) = \frac{c(t)^6}{P}\,. \eqno(7.2)
	\]
	
	
	\medskip
	\bigskip
	\textbf{玻尔兹曼电离判据 $\boldsymbol{T_\mathrm{ion}(t)}$ 的演化}\,\,\,
	如表 4 所分析，氢原子的电离温度随宇宙光速标度，遵循 ${T_\mathrm{ion}}(t) \propto c(t)^2$。今日的 ${T_\mathrm{ion}}(t)$ 约为 $10^8\,\mathrm{K}$；四个典型宇宙时期下 ${T_\mathrm{ion}}(t)$ 的取值列于表 5。
	
	\begin{table}[H]
		\centering
		\captionsetup{skip=2pt, font=small}
		\caption{基准时期 $\boldsymbol{z_i}$ 的电离温度 $\boldsymbol{T_\mathrm{ion}}$}
		\footnotesize
		\renewcommand{\arraystretch}{1.50}
		\begin{tabular}{>{\centering\arraybackslash}p{2.5cm}|>{\centering\arraybackslash}p{2cm}>{\centering\arraybackslash}p{2cm}>{\centering\arraybackslash}p{2cm}>{\centering\arraybackslash}p{2cm}}
			\toprule
			\textbf{参数} & \textbf{$\boldsymbol{z_0}$} & \textbf{$\boldsymbol{z_1}$} & \textbf{$\boldsymbol{z_2}$} & \textbf{$\boldsymbol{z_3}$} \\
			\midrule
			$\boldsymbol{C(t)}$      & $1/1100$   & $1/25$    & $1/3.7$   & $1/1.43$  \\
			$\boldsymbol{T_\mathrm{ion}\,(\mathrm{K})}$  & $>82.64$ & $>1.60\times10^{5}$ & $>7.31\times10^{6}$ & $>4.89\times10^{7}$ \\
			\bottomrule
		\end{tabular}
		\label{tab:ion_temp}
	\end{table}
	
	据玻尔兹曼公式 (1.8)，若宇宙光速降至今日值的 \(1/7137.66\)，即 $c(t)=42.03\,\mathrm{km/s}$，则电离温度降为 \(1.96\,\mathrm{K}\)。换言之，在这一光速下，任何高于 \(1.96\,\mathrm{K}\) 的温度都足以把中性氢原子解离为等离子体。
	
	\medskip
	如第 9 节将进一步讨论的，把温度从 \(1.96\,\mathrm{K}\) 升至 \(2.75\,\mathrm{K}\)，会引发夸克解囚禁与强子解离，使等离子体转化为自由夸克与胶子之海。我们将在同一节中对该现象以及最初的零光速零温度（ZLZT）宇宙状态作详细分析。
	
	
	\medskip
	\bigskip
	\textbf{史瓦西半径 $\boldsymbol{r_s(t)}$ 的演化}\,\,\,
	如表 4 所分析，式 (1.11) 所定义的半径按 $r_s(t) \propto c(t)^{-4}$ 标度。在高红移时期 $1+z=25$ 与 $1+z=7$，$r_s(t)$ 分别为今日值的 $25^4 = 390{,}625$ 倍与 $7^4 = 2{,}401$ 倍。在 $1+z=3.7$ 的宇宙正午，该半径比今日大 $3.7^4 \approx 187.4$ 倍。
	
	\medskip
	早期宇宙中 $r_s(t)$ 的大幅增大，对结构的形成产生了强烈影响。引力作用范围远为扩大之后，早期恒星与致密天体能够轻易地吸积周围物质，从而迅速清除自身附近的物质，使行星几乎没有形成的机会。与此同时，更强的引力把这些恒星压得更紧，使它们更快地燃尽燃料、更快地演化。这为“为何最早的星系似乎以恒星为主、极少有行星系统的迹象”提供了一种新的理解——这一点在以往的相关讨论中尚未被提出。
	
	\medskip
	随着宇宙 $c(t)$ 继续向当前值增长，$r_s(t)$ 按 $c(t)^{-4}$ 稳步收缩。引力影响的这种减弱，意味着较晚形成的恒星远不如早期恒星那样能够俘获周围物质。到太阳系在 $1+z \approx 1.43$ 形成时，$r_s(t)$ 已降至宇宙正午时值的约 $(1.43/3.7)^4 = 22.3\%$。因此，早期宇宙是一个引力活动剧烈的场所，随着时间推移逐渐平静下来，从而容许像我们这样较为“安静”的行星系统在较晚的时期形成。
	
	\medskip
	这一图景的一个自然延伸，是把行星形成受到的抑制量化为红移的函数。可以启发式地写出
	\[
	P_{\rm planet}(z) \; \propto \; \left(\frac{r_s(z)}{r_{\rm disk}(z)}\right)^{-\beta} \propto \; (1+z)^{4\beta},\,\, \eqno(7.3)
	\]
	
	其中 $r_{\rm disk}(z)$ 是原行星盘的特征尺度，$\beta>0$ 是唯象指数。更严格的推导可从引力聚焦或 Bondi 吸积物理入手，把增强的俘获截面与盘的存续概率以及随后的行星组装联系起来。
	 
	\medskip
	然而，这样的定量处理超出本工作的范围。我们给出这一表达式，只是作为一项启发性的提示：若未来观测数据可得，增长-$c(t)$ 框架或可对行星系统形成的红移依赖性给出可检验的预言。
	
	
	\medskip
	\bigskip
	\textbf{狄拉克大数 $\boldsymbol{N_{\mathrm{Dirac,1}}(t)}$ 的演化}\,\,\,
	如表 4 所分析，式 (1.16) 中的狄拉克大数比~\cite{Dirac:1937} 按 $N_{\mathrm{Dirac,1}}(t) \propto c(t)^{5}$ 标度。今日的比值约为 $\approx 2.27 \times 10^{39}$。在 $1+z=1100$ 的 CMB 时期，该大数缩小至当前值的 $(1/1100)^5$，给出 $N_{\mathrm{Dirac,1,CMB}} \approx 1.40\times 10^{24}$。式 (1.17) 中电磁力与引力之比同宇宙年龄与质子光穿越时间之比在数值上的接近，不过是今日这一时期的偶然巧合；在早期宇宙中，这两个数并不接近。
	
	\medskip
	请注意，电磁力与引力之比在 CMB 时期出现如此剧烈的下降，是我们增长-$c(t)$ 框架的一项独特预言。较小的 $N_{\mathrm{Dirac,1}}$ 意味着，在早期宇宙的动力学中，引力效应相对而言更为重要。

	
	\medskip
	\bigskip
	\textbf{重新审视 SI 中 \texttt{秒}与 \texttt{米}的定义}\,\,\,
	一\texttt{秒}被定义为一颗静止的铯-133 原子在特定能级跃迁中辐射出的、恰好 \texttt{9,192,631,770} 个周期的持续时间~\cite{BIPM:2019}。铯原子频率是一个纯粹的微观量子不变量，由精细结构常数与原子的内禀能级结构决定；该频率并不依赖光速。一\texttt{米}则被定义为光在真空中一\texttt{秒}内所行路径长度除以 \texttt{299,792,458}。
	
	\medskip
	上述\texttt{米}的定义隐含地假设：真空中的宇宙光速永久恒定。在今日光速 $c_0$ 下，铯-133 跃迁所对应的波长为：
	\smallskip
	\[
	\lambda_{\text{cs}} = \frac{c_0}{\nu_{\text{cs}}}
	= \frac{299\,792\,458\ \mathrm{m/s}}{9\,192\,631\,770\ \mathrm{Hz}}
	\approx 0.032612\ \mathrm{m} = \texttt{32.612}\ \mathrm{mm}\,\, \eqno(7.4)
	\]
	
	\medskip
	按照上述\texttt{米}的定义，若早期宇宙的宇宙光速降至 $75\,000\ \mathrm{km/s}$，即今日值 $300\,000\ \mathrm{km/s}$ 的四分之一，则铯-133 原子的本征辐射频率保持不变，而对应波长 $\lambda_{\text{cs}}$ 缩小至现代量级的四分之一，即 $\texttt{8.153}\ \mathrm{mm}$。
	
	\medskip
	在增长-$c(t)$ 的公设之下，逻辑上必须区分两个不同概念：\texttt{现代米}与\texttt{原子米}，分别记为 $\boldsymbol{L_\mathrm{m}}$ 与 $\boldsymbol{L_\mathrm{a}}$。此处 ${L_\mathrm{m}}$ 表示当前使用的标准米，而 ${L_\mathrm{a}}$ 指铯-133 原子辐射的 ${9\,192\,631\,770} / {299\,792\,458}\simeq 30.6633$ 个波形周期内所行经的距离。
	
	\medskip
	换言之，${L_\mathrm{a}}$ 以固定的原子振荡周期为基准来定义，故周期总数永久恒定。然而，${L_\mathrm{a}}$ 所代表的实际空间延展，却随宇宙光速按增长-$c(t)$ 演化而改变。这一构想与教科书中的常规定义显著不同，请读者务必留意并充分理解其内在逻辑。
	
	
	\medskip
	\bigskip
	\textbf{重新审视光子频率、波长与光速}\,\,\,
	我们从基本物理出发。受精细结构常数与原子系统内禀量子化能隙的约束，光子的本征频率 $\nu$ 与角频率 $\omega$ 满足：
	\[
	\Delta E = h\,\nu = 2\pi\hbar\,\nu = \hbar\,\omega \,\,  \eqno(7.5)
	\]
	
	\smallskip
	$\Delta E$ 与 $\hbar$ 同步演化，且都按 $c(t)^2$ 标度。因此，光子的发射频率 $\nu$ 内禀恒定，与光速无关。
	
	\medskip
	相比之下，波长并非光子的内禀属性。它随光速与传播介质的折射率而变，遵循电磁波的基本关系：
	\[
	c = \lambda \,\nu \,\,  \eqno(7.6)
	\]
	
	我们以标准光学玻璃来直观演示波长的变化。对真空光速 300\,000 km/s，考察一个 $\nu_{n0}=1.0\times10^{15}\,\mathrm{Hz}$、$\lambda_{n0}=300\,\mathrm{nm}$ 的光子。当它进入折射率 $N_{\text{glass}}=1.5$ 的玻璃后，其传播速度降至 200\,000 km/s，波长随之缩短为 $200\,\mathrm{nm}$。当它离开玻璃、返回真空时，波长恢复为原来的 $300\,\mathrm{nm}$。
	
	\medskip
	\textit{\textbf{注：}真空光速充当普适的宇宙背景参数，而非单个光子的属性。这一共同速度由一个迄今尚未被探究的宇宙暗场所支配，第 8 节与第 9 节将把它确认为 LEM-SCDF 场。}
	
	\medskip
	在标准宇宙学红移的框架内，我们考察一个本征频率 $\nu_{n1}=1.0\times 10^{15}\,\mathrm{Hz}$ 的远古光子，它在红移 $z=3.0$（即 $1+z=4$）处发出。在宇宙膨胀的驱动下，其观测频率在今日降至 $2.5\times 10^{14}\,\mathrm{Hz}$，等于初值的四分之一，而实测波长则变为原来的四倍。
	
	\medskip
	对一个跨越整个宇宙尺度 $a(t)$ 传播的光子，标准红移理论确立的结论是：沿其轨迹的波周期总数守恒。当宇宙尺度从 $\frac{1}{4}a_0$ 膨胀至当前尺度 $a_0$ 时，光子频率下降、波长伸展以适应扩大的空间尺度，而波峰与波谷的总数保持不变。我们把这一现象称为“红移的波周期数不变性”。
	
	\medskip
	我们的增长-$c(t)$ 模型拓展并丰富了与宇宙学红移相关的物理不变量集合。在 $1+z=4.0$ 的时期，宇宙光速为 $75\,000\,\mathrm{km/s}$，是今日 $c_0$ 的四分之一。当宇宙空间尺度从 $a(t)=\frac{1}{4}\, a_0$ 演化至 $a_0$ 时，光速 $c(t)$ 同步增大。光子穿越整个宇宙尺度所需的旅行时间，在整个宇宙膨胀过程中守恒。
	
	\medskip
	波周期数与传播时间这一双重不变量，揭示了基本物理定律深刻的内在和谐与优雅之美。本质上，任何被发出并跨越整个宇宙尺度 $a(t)$ 传播的光子，都将保持固定的波周期数与恒定的旅行时间，不受宇宙膨胀的影响。这种内在的相合性在我们的增长-$c(t)$ 模型中得到了清晰的体现。
	
	\medskip
	我们进一步分析那个 $\nu_{n0}=1.0\times10^{15}\,\mathrm{Hz}$、在 $1+z=4.0$ 处发出的远古光子。以 $L_\mathrm{m}$ 为标定基准，当时宇宙光速为 $75\,000\,\mathrm{km/s}$，光子的初始波长为 $75\,\mathrm{nm}$。经过长距离传播至今，宇宙学红移把其频率降至 $2.5\times 10^{14}\,\mathrm{Hz}$。与此同时，宇宙真空光速已升至 $300\,000\, \mathrm{km/s}$，观测到的最终波长达到 $1200\,\mathrm{nm}$，即发射时波长的 16 倍。
	
	\medskip
	总计 16 倍的波长放大来自两个彼此独立的效果：其中四倍贡献于宇宙学红移，另四倍贡献于宇宙真空光速 $c(t)$ 的单调上升。
	
	\medskip
	在光速恒定的假设下，16 倍的波长移动显得违反直觉。然而，若改用 $L_\mathrm{a}$ 来解读，我们便看到：光子的频率不过是降至发射时值的四分之一。
	
	
	\medskip
	\bigskip
	\textbf{以比喻说明光子频率背后的物质性}\,\,\,
	正如一个人的先天特质由父母的基因决定，光子在发射时的频率仅由精细结构常数与其辐射源的结构决定，与外部宇宙环境无关。
	
	\medskip
	在脱离其源之后，光子的固有频率便开始受到宇宙背景的调制。宇宙膨胀与大质量天体的引力场不断改变这一本征频率，主要产生红移效应——这很像一个人进入社会之后，社会的物质条件对其施加的塑造。
	
	
	
	\medskip
	\bigskip
	至此，我们已经描绘了一系列彼此显著不同的宇宙演化物理图景。关于爱因斯坦 $G_{\mu\nu}$ 场方程以及本工作所提出的光—电—磁同步宇宙暗场（\textit{LEM-SCDF}）的讨论，将在随后的第 8 节展开。
	
	
	
	
	
	\section*{
		\begin{center}
			8. 修正后的 $\boldsymbol{G_{\mu\nu}}$、常数 $\boldsymbol{P}$ \\
			与 LEM‑SCDF 场
		\end{center}
	}
	\medskip
	\noindent
	
	我们回到第 2 节的宏观公设：存在一个普适且不随时间变化的宇宙耦合常数 $\boldsymbol{P}$，而含时光速 $c(t)$ 满足如下守恒律
	\[
	G(t)\tinysp c(t)^2 = \boldsymbol{P} \,\,.  \eqno(8.1)
	\]
	
	\smallskip
	借助这一公设，我们修正爱因斯坦的曲率场方程，并考察 $G(t)c(t)^2$ 这一幂律关系的物理合理性。此外，我们引入 LEM‑SCDF，即光、电、磁三者全局同步的宇宙暗场，作为本框架一个具有理论动机的构造。
	
	\medskip
	\textit{\textbf{注：}早在 1937 年与 1938 年，狄拉克就已设想基本常数或许也随宇宙时间演化。他进一步论证，若引力常数 $G$ 随宇宙时期按 $G(t^\Delta) \propto 1/t^\Delta$ 变化，则物理学与宇宙学中大无量纲数之间惊人的巧合便可得到解释；而这样的变化必然也牵涉光速 $c(t)$ 的改变~\cite{Dirac:1937, Dirac:1938}。}
	
	
	\medskip
	\bigskip
	\textbf{修正后的 $\boldsymbol{G_{\mu\nu}}$ 方程}\,\,\,
	把上述公设代入式 (1.9) 所定义的爱因斯坦曲率场方程~\cite{Einstein:632320}，得到：
	$$
	G_{\mu\nu} = \frac{8\pi\tinysp G(t)}{{c(t)}^4} T_{\mu\nu} = \frac{8\pi\tinysp P}{{c(t)}^6} T_{\mu\nu} \,\,.  \eqno(8.2)
	$$
	
	\smallskip
	在式 (8.2) 中，$T_{\mu\nu}$ 表示描述质能分布的能量—动量张量。虽然早期宇宙 $T_{\mu\nu}$ 的完整形式无法精确确定，但请注意：若早期星系形成时期（$1+z=25$）与 CMB 时期（$1+z=1100$）的光速分别为今日光速 $c_0$ 的 $1/25$ 与 $1/1100$，则时空曲率将极其巨大，其量级为 $25^6\approx 2.4\times 10^{8}$ 或 $1100^6\approx 1.8\times 10^{18}$。
	
	\medskip
	原初时期如此极端的时空曲率，在物质与天体之间引发了强耦合，进而要求这些天体用极长的时标才能达到弛豫平衡。因此，早期时期的时空曲率（由爱因斯坦张量 $G_{\mu\nu}$ 反映）极为巨大，并在相继的宇宙时代中呈现定性的递减趋势：CMB 时期 $z_0$、早期星系形成时期 $z_1$、宇宙正午时期 $z_2$，以及太阳系形成时期 $z_3$。
	
	\medskip
	张量 $T_{\mu\nu}$ 包含能量密度、动量通量、压强、应力与电磁场。在宏观尺度上，采用弱场近似与慢动近似并结合关系 $E=m c^2$，式 (8.2) 约化为
	$$
	G_{\mu\nu} \thinspace \approx  \thinspace \frac{8\pi\tinysp P}{{c(t)}^6}\cdot  M_{\mu\nu}\tinysp \cdot {{c(t)}^2} 
	\thinspace = \thinspace \frac{8\pi\tinysp P}{{c(t)}^4} M_{\mu\nu}\,\,  \eqno(8.3)
	$$
	
	\smallskip
	此处，$M_{\mu\nu}$ 表示物质的密度分布。
	
	\medskip
	在公设 $G(t)\tinysp c(t)^2={P}$ 之下，比较式 (8.2) 与式 (8.3) 可见，两式右端在物理上完全一致。这表明关系 $\boldsymbol{G\tinysp c^2=P}$ 以\textbf{\textit{某种方式}}本就内嵌于爱因斯坦场方程之中。
	
	\medskip
	反之，若采用诸如 $G(t)\tinysp c(t)={P}$、$G(t)\tinysp c(t)^3={P}$ 或 $G(t)\tinysp c(t)^4={P}$ 之类其他幂律形式，这种一致性便不复存在。尽管与 $G(t)$ 相关的量在纯计算的意义上是数学上容许的，它们却缺乏物理上合理的解释。因此，本工作不再进一步讨论这些替代幂律形式。
	
	\medskip
	利用式 (7.1) 所定义的麦克斯韦方程组替换式 (8.2) 中的 $c(t)^6$，我们得到修正 $G_{\mu\nu}$ 方程一个涉及 $\mu_0(t)$ 与 $\varepsilon_0(t)$ 的关联形式，如下：
	$$
	G_{\mu\nu} = 8\pi\tinysp P\cdot \mu_0(t)^3\cdot \varepsilon_0(t)^3\cdot T_{\mu\nu} \,\,   \eqno(8.4)
	$$
	
	
	\medskip
	\bigskip
	\textbf{修正 $\boldsymbol{G_{\mu\nu}}$ 方程的理论洞见}\,\,\,
	宇宙质量与能量的主要部分源自强相互作用。夸克与胶子之间的束缚能造就了可见物质的大部分质量，因此能量—动量张量 $T_{\mu\nu}$ 实质上反映的是宇宙中强相互作用效应的宏观分布。
	
	\medskip
	爱因斯坦的原始场方程构成一个纯粹单向的几何响应模型。由 $T_{\mu\nu}$ 描述的一切质能分布变化，都只通过时空曲率 $G_{\mu\nu}$ 来平衡，没有任何额外的物理自由度可以分担这一动力学负担。
	
	\medskip
	与此相对，我们的修正框架建立了一个闭环调控系统。时空曲率 $G_{\mu\nu}$ 与含时光速 $c(t)$ 协同作用，共同响应并平衡 $T_{\mu\nu}$ 的动态变化。
	
	\medskip
	由于 $c(t)$ 由真空电磁参数 $\mu_0(t)$ 与 $\varepsilon_0(t)$ 决定，整个宇宙系统便有效地把微观强相互作用的质源、大尺度的时空几何效应，以及底层的真空电磁规则连接起来。这一动力学过程中的所有物理变量，都严格受不变量宇宙常数 $P$ 的约束。
	
	\medskip
	因此，修正后的 ${G_{\mu\nu}}$ 方程不再属于纯粹的几何引力理论。它构成一个统一的宇宙动力学框架：以真空电磁基础为根基，以时空曲率与含时光速为双重调控通道，并以强相互作用主导的质能分布为核心源项。
	
	
	\medskip
	\bigskip
	\textbf{常数 $\boldsymbol{P}$ 的内在含义}\,\,\,
	${P}$ 被称为宇宙耦合常数，因为它刻画了宇宙质能分布、时空曲率、含时光速 $c(t)$，以及含时真空磁导率 $\mu_0(t)$ 与真空介电常数 $\varepsilon_0(t)$ 之间的基本耦合。
	
	\medskip
	对式 (8.4)、式 (8.2) 与式 (8.3) 作代数变换，可得：
	$$
	\boldsymbol{P} \thinspace = \thinspace \dfrac{G_{\mu\nu}}{8\pi\tinysp \mu_0(t)^3\cdot \varepsilon_0(t)^3\cdot T_{\mu\nu}}
	\thinspace = \thinspace \dfrac{{{c(t)}^6}\cdot {G_{\mu\nu}}}{8\pi\tinysp \cdot T_{\mu\nu}}
	\thinspace \approx \thinspace \dfrac{{{c(t)}^4} \cdot {G_{\mu\nu}}}{8\pi\tinysp \cdot M_{\mu\nu}} \,\,   \eqno(8.5)
	$$
	
	
	\medskip
	\bigskip
	\noindent
	$\triangleright \triangleright \triangleright$\enspace
	历史上，爱因斯坦对引力常数抱有一种矛盾的态度。尽管广义相对论把引力现象完全归因于时空曲率，使 $G$ 不再是基本独立量，爱因斯坦却仍把 $G$ 保留在场方程之中。这一选择主要是为了确保方程在弱场与慢动极限下能退化为牛顿引力，而非出于严格的理论必然性。
	
	\medskip
	1915 年之后的约十年间，爱因斯坦与当时大多数物理学家都坚持静态宇宙范式——直到 1929 年哈勃发现宇宙膨胀，这一观点才被撼动。这一思想框架使人们几乎无法设想变化的引力或光速；然而，不变量关系 \(G\tinysp c^2=P\) 却以某种方式内嵌于这个方程之中。直到 1937 年，狄拉克才开始质疑 \(G\) 是否随宇宙时间演化：$G(t^\Delta) \propto 1/t^\Delta$。\enspace$\triangleleft \triangleleft \triangleleft$	
	
	
	
	
	
	\medskip
	\bigskip
	\textbf{常数 $\boldsymbol{P}$ 在弗里德曼层面的解读}\,\,\,
	当我们把修正后的场方程带入弗里德曼框架，$P$ 的内在含义便更加清晰。这不是理论的扩展，而是在宇宙膨胀的层面上对同一条公设的重新解读。
	
	\medskip
	在标准宇宙学中，式 (1.14) 的第一个弗里德曼方程把膨胀速率 $H(t)$ 与宇宙平均密度 $\rho_{\scriptscriptstyle\mathrm{F}}(t)$ 联系起来：
	\[
	H^2(t) = \frac{8\pi G}{3}\rho_{\scriptscriptstyle\mathrm{F}}(t) + \cdots \,\,.  \eqno(8.6)
	\]
		
	\medskip
	此处 $G$ 被当作固定常数，于是宇宙膨胀的全部负担都只落在 $\rho_{\scriptscriptstyle\mathrm{F}}(t)$ 上。当宇宙稀释时，$H(t)$ 必然相应减小。把宏观公设 $G(t)=P/c(t)^2$ 代入同一方程，则得到一种结构上不同的平衡：
	\[
	H^2(t) = \frac{8\pi P}{3\,c(t)^2}\rho_{\scriptscriptstyle\mathrm{F}}(t) + \cdots \,\,.  \eqno(8.7)
	\]
	
	\medskip
	常数 $G$ 消失了。剩下的是膨胀速率 $H(t)$、宇宙平均密度 $\rho_{\scriptscriptstyle\mathrm{F}}(t)$ 与宇宙光速 $c(t)$ 之间的三元关系，而三者被不变量 $P$ 锁合在一起。
	
	\medskip
	重排式 (8.7)，这种锁合关系便显式呈现：
	\[
	\boldsymbol{P} \thinspace \sim \thinspace \frac{H^2(t)\,c(t)^2}{\rho_{\scriptscriptstyle\mathrm{F}}(t)} \,\,.  \eqno(8.8)
	\]
	
	\medskip
	以这种形式来读，$P$ 既不是密度，也不是速度，也不是曲率。它是把三者绑定在一起的不变乘积。若宇宙密度随膨胀而下降，且光速按 $c(t)$ 上升，那么膨胀速率 $H(t)$ 必须以这样一种方式调整：使右端的组合保持固定。$P$ 的量纲结构支持这一读法。由 $[G c^2]=\mathrm{L^5 M^{-1} T^{-4}}$，可以把该量分离成两个物理上可辨识的部分：
	
	\[
	\mathrm{L^5 M^{-1} T^{-4}}
	\quad  = \thinspace
	\underbrace{\mathrm{L^3 M^{-1}}}_{\text{密度倒数 $1/\rho$}}
	\quad \quad  \thinspace \times \quad 
	\underbrace{\mathrm{T^{-2} \cdot L^2 T^{-2}}}_{\text{膨胀率 $H^2$ 与光速 $c^2$ 因子}} \,\,.  \eqno(8.9)
	\]
	
	\medskip
	第一个因子是 $1/\rho$，即密度的倒数。第二个因子带有 $H^2 c^2$ 的量纲，恰是式 (8.8) 中出现的那一组合。于是不变量 $P$ 可以读作一把“密度—膨胀率—光速”之锁：它确定了宇宙密度的稀释、宇宙膨胀的速率与真空光速的增长之间必须如何相互补偿。
	
	\medskip
	这便为 $P$ 在弗里德曼层面的含义给出了清晰的物理图景。随宇宙膨胀，$\rho_{\scriptscriptstyle\mathrm{F}}(t)$ 减小。在标准框架中，仅这一减小就使 $H(t)$ 变慢。在我们的新框架中，同样的减小可以部分地被 $c(t)$ 的增大所吸收，而 $G(t)$ 随 $G(t)=P/c(t)^2$ 同步下降。这三个量并非各自独立的行动者，而是同一条不变关系的三条协同通道。
	
	\medskip
	因此，我们可以用一句话概括 $P$ 的角色：\textit{$\boldsymbol{P}$ 是宇宙密度 ${\rho_{\scriptscriptstyle\mathrm{F}}(t)}$、宇宙膨胀速率 $H(t)$ 与光速 $c(t)$ 之间的锁合常数}。它是在容许 $G(t)$ 与 $c(t)$ 演化之后依然存留的不变量，也是把修正后的爱因斯坦场方程与宇宙膨胀的弗里德曼描述联系起来的那一量。



		
	\medskip
	\bigskip
	\textbf{LEM-SCDF 场}\,\,\,
	在增长-$c(t)$ 框架内，真空电磁参数 $c$、$\varepsilon_0$ 与 $\mu_0$ 不再是普适常数。遵循理论物理中的标准公设范式，本工作引入光—电—磁同步宇宙暗场（LEM-SCDF）作为基本的宇宙背景场。
	
	\medskip
	本研究聚焦于 LEM–SCDF 暗场在宇宙演化中的角色：从强子化后的等离子体时期，历经 CMB 时代，直至今日宇宙。一切由 QCD 支配的夸克—胶子动力学与强子化过程，均不在当前讨论范围之内。
	
	\medskip
	LEM--SCDF 暗场的核心功能，是在整个宇宙空间中实现含时基本参数 $c(t)$、$\varepsilon_0(t)$ 与 $\mu_0(t)$ 的全局同步，同时保持它们严格的空间均匀性。
	
	\medskip
	我们以 $\boldsymbol{v_{\scriptscriptstyle\mathrm{LEM}}(t)}$ 表示 LEM-SCDF 暗场的传播速度，它满足 $v_{\scriptscriptstyle\mathrm{LEM}}(t) \gg c(t)$。我们定义如下不变关系
	\[
	v_{\scriptscriptstyle\mathrm{LEM}}(t)\cdot \rho_{\scriptscriptstyle\mathrm{F}}(t) = c(t)\cdot \rho_{\scriptscriptstyle\mathrm{P}}(t) \,\,\,,   \eqno(8.10)
	\]
	或等价地，
	\[
	v_{\scriptscriptstyle\mathrm{LEM}}(t)= c(t)\cdot \frac{\rho_{\scriptscriptstyle\mathrm{P}}(t)}{\rho_{\scriptscriptstyle\mathrm{F}}(t)} \,\,\,.   \eqno(8.11)
	\]
	
	\smallskip
	此处，$\rho_{\scriptscriptstyle\mathrm{F}}(t)$ 表示出现在第一与第二弗里德曼方程（式 (1.14) 与式 (1.15)）中的宇宙平均密度，是支配宇宙膨胀的核心密度项。$\rho_{\scriptscriptstyle\mathrm{P}}(t)$ 则是第 7 节讨论过的普朗克密度，它是量子引力所施加的密度上限，是时空自身的一种固有属性。
	
	\medskip
	在式 (8.10) 的右端 $c(t)\cdot\rho_{\scriptscriptstyle\mathrm{P}}(t)$ 中，$c(t)$ 是局域时空中电磁与引力相互作用的基准传播速度，而 $\rho_{\scriptscriptstyle\mathrm{P}}(t)$ 是时空在量子引力约束下所能承载的最大质量密度。因此二者之积给出单位体积内相互作用传递的内禀上限。
	
	\medskip
	在左端 $v_{\scriptscriptstyle\mathrm{LEM}}(t)\cdot\rho_{\scriptscriptstyle\mathrm{F}}(t)$ 中，弗里德曼平均密度 $\rho_{\scriptscriptstyle\mathrm{F}}(t)$ 表示膨胀宇宙中物质被稀释后的实际分布，而 $v_{\scriptscriptstyle\mathrm{LEM}}(t)$ 是 LEM--SCDF 暗场的特征传播速度。
	
	\medskip
	这一对称的乘积结构，在微观量子时空极限与宏观大尺度宇宙演化之间建立起一种对偶对应，且对任意宇宙时间 $t$ 都成立。两个乘积相等，使时空的量子尺度上限与整个宏观宇宙的协调需求之间维持永久的平衡。该定义在数学上性质良好，没有奇点，也没有不可微点，兼具形式上的优雅与清晰的物理图景。
	
	\medskip
	由于 $\rho_{\scriptscriptstyle\mathrm{P}}(t)$ 是量子引力所容许的绝对密度上限，且除宇宙膨胀的初始奇点态（那是一种逻辑上的理想化）之外，在所有后续演化阶段都有 $\rho_{\scriptscriptstyle\mathrm{P}}(t) \gg \rho_{\scriptscriptstyle\mathrm{F}}(t)$，这一不等式便直接意味着：在任意给定时刻都有 $v_{\scriptscriptstyle\mathrm{LEM}}(t) \gg c(t)$。
	
	\medskip
	$v_{\scriptscriptstyle\mathrm{LEM}}(t)$ 超出局域真空光速许多个数量级，故暗场的传播时延在所有可观测宇宙尺度上都小到可以忽略。这一特性自然产生近乎瞬时的全局同步，从而为本工作所建立的核心不变量 $G(t)\tinysp c(t)^2=\text{constant}$ 提供支撑。
	
	\medskip
	由 LEM--SCDF 暗场主导的全局同步，不需要任何物质介质，不涉及能量或物理信号的传输，也不意味着发射粒子与接收粒子之间存在因果关联。它的全局协调本质上是宇宙背景自身的一种态跃迁。因此，其超光速传播并不违反狭义相对论的局域光锥因果约束。
	
	\medskip
	\textit{\textbf{注：}速度 $\boldsymbol{v_{\scriptscriptstyle\mathrm{LEM}}(t)}$ 恰恰就是支配量子纠缠的速度——量子纠缠是一种植根于量子真空的基本非局域关联。这便化解了爱因斯坦在与玻尔的论战中以“鬼魅般的超距作用”一语著称、悬置百年的悖论。爱因斯坦与玻尔激烈争论的那个形而上学之谜，如今获得了物理上的本体：LEM-SCDF 场本身就是量子纠缠在宇宙尺度上实现全局同步的介质。}
	
	\medskip
	宇宙中的一切质量与能量共同参与塑造这些 LEM 基本参数的演化。因此，LEM-SCDF 暗场提供了一个统一的基础物理框架，把麦克斯韦电磁学、杨--米尔斯规范场论与爱因斯坦广义相对论一致地整合起来。
	
	
	\bigskip
	\textit{\textbf{注：}标准的麦克斯韦真空参数 $c$、$\varepsilon_0$ 与 $\mu_0$ 是由今日洁净真空实验测得的常数。没有任何证据表明这些固定值适用于宇宙演化的所有时代。在 LEM-SCDF 框架中，我们放弃静态的真空电磁常数 $c$、$\varepsilon_0$ 与 $\mu_0$，转而用含时的全局平均背景参数来描述宇宙的电磁基底。现代式的洁净真空对应一种理想的电磁环境，不能被当作宇宙的普遍状态。这样的理想真空条件，至少在两个关键宇宙时期并不存在：CMB 退耦之前的强子化后等离子体阶段，以及直到红移 $1+z=25$ 的再电离时期——那时整个宇宙充满电离等离子体。把今日的真空电磁标准直接套用于这些早期宇宙阶段，是不合理的。}
	
	
	
	\medskip
	\bigskip
	\textbf{统一场论难题的另一条路径}\,\,\,
	传统统一场论（UFT）研究长期受两大根本障碍的困扰：引力的量子化，以及光子散射问题。
	
	\medskip
	引力的量子化构成第一个障碍：广义相对论与标准的杨--米尔斯规范场框架在结构上存在本质的不相容，使常规的量子化方案无从施行。
	
	\medskip
	我们修正后的 $G_{\mu\nu}$ 方程放弃了传统的常数引力参数 $G$，改用随时间演化的变量 $c(t)$。通过转向 LEM 场的量子化，这一修正的曲率表述或可为绕开引力量子化的困境提供一条可行的替代路径。
	
	\medskip
	作为进一步探索的一个尝试性方向，我们指出：结合式 (8.5) 与式 (1.6)，可以在宏观引力部分与微观量子领域之间建立一种形式上的联系：
	\[
	\boldsymbol{P} \thinspace = \thinspace \dfrac{G_{\mu\nu}}{8\pi\tinysp \mu_0(t)^3\cdot \varepsilon_0(t)^3\cdot T_{\mu\nu}}
	 \thinspace = \thinspace \dfrac{c(t)^6 \cdot G_{\mu\nu}}{8\pi \cdot T_{\mu\nu}}
	 \thinspace = \thinspace \dfrac{\hbar(t)^4 \cdot \rho_{\scriptscriptstyle\mathrm{P}}(t)}{m_{\scriptscriptstyle\mathrm{P}}(t)^6} \,\,\,. \eqno(8.12)
	\]
	
	\medskip
	上式把广义相对论与量子力学的语言汇聚于一处。其左端以时空曲率与质能说话，右端则以普朗克尺度的量说话。
	
	\bigskip
	\noindent
	$\triangleright \triangleright \triangleright$\enspace
	在广义相对论完成之后的几年（1915--1920），爱因斯坦重新把时空定义为一个具有内禀属性的物理连续体。他 1920 年在莱顿大学关于新以太的演讲《以太与相对论》~\cite{Einstein:1922ether}——这部经典文本近期已有现代译本《爱因斯坦：论以太与相对性》~\cite{LeidenLecture:2025}——明确反对空无一物的空间概念。他主张：
	\begin{itemize}
		\item 时空并非空无的几何背景；几何本身就是时空的物理属性，曲率是时空自身固有的物质属性；
		\item 物质—能量张量改变时空的物理属性（曲率与密度），这些时空属性反过来又支配物质与光的轨迹；
		\item 不存在任何剥除物质属性的“\textit{真空}”——\textit{真空}是一种具有可变物理结构的连续介质；
		\item 按照广义相对论，没有以太的空间是不可设想的：在这样的空间中，不仅光的传播将不可能，而且量尺与时钟——空间与时间之所以为空间与时间的那些标准——也不可能有物理上的存在，从而无法定义物理意义上的时空间距。
	\end{itemize}
	
	在反思引力场与电磁场的关系时，爱因斯坦指出：引力场与空间本身不可分离，而“完全可以设想某一部分空间没有电磁场；因此与引力场相反，电磁场似乎只是一种次级现象”。他随即表达了这样的愿望：“{\textit{如果我们能够成功地把引力场与电磁场理解为同一种统一构型，那将是一个巨大的进步；那时，以太与物质之间的对立便会消融}}。”
		
	\medskip
	我赞同爱因斯坦的愿望，但不赞同他把电磁场视为次级现象的判断。在我们所设想的 LEM-SCDF 物理图景中，这两种场在量子层面本质上是统一的；量子化后的 LEM-SCDF 本身已把时空曲率包含在内，这提示电磁学或许是这一综合中更为基本的伙伴；而 LEM-SCDF 本身则在普朗克尺度上与一种全局量子场论相统一。不过，对其中底层物理的详细阐述超出当前讨论的范围，因为这一框架仍不完整，尚需进一步发展。\enspace$\triangleleft \triangleleft \triangleleft$
	

	\bigskip
	第二个关键问题是光子散射问题。标准 \(\Lambda\)CDM 宇宙学假设一个炎热原初宇宙，且今日 \(c\) 恒定，这会产生强烈的光子散射，阻碍早期的粒子组装与强子化。宇宙从不透明等离子体到透明介质的转变，受若干耦合散射过程支配，例如自由电子上的汤姆孙散射与中性氢上的瑞利散射，现有文献对此已有记载~\cite{Coulton:2021}。
	
	\medskip
	为解决视界疑难与平坦性疑难，Magueijo 提出了一个早期 VSL（变光速）模型~\cite{Magueijo:2003}，其中假定早期宇宙的光速远比今日的 \(c\) 为大。然而这一拟设放大了光子相互作用率与平均自由程，反而加剧了强子化与中性原子形成所面临的散射壁垒——而这恰恰是一个成功的宇宙学必须解决的问题。
			
	\medskip
	在增长-\(c(t)\) 框架内，该问题通过所设定的 \textit{ZLZT}（零光速与零温度）原初宇宙而得到处理：其中光速与温度都渐近趋于零。这一初始条件消除了标准宇宙学与传统 VSL 理论所固有的强烈光子散射。
		
	
	\medskip
	\bigskip
	接下来的第 9 节将进一步讨论 ZLZT 原初公设的物理含义。该节考察从强子化后等离子体阶段直至 CMB 时期的宇宙演化，重点在于{\textit{元粒子物质}}的物理生成，以及光速从初始近零状态逐渐增长到 CMB 时期 $c(t)=1/1100$ 的过程。
	
	
	
	
	
	\section*{
		\begin{center}
			9. ZLZT 原初态、\\
			囚禁生成物质、\\
			与 $\boldsymbol{c(t)}$ 的演化
		\end{center}
	}
	\medskip
	\noindent
	
	在光速恒定的假设下，SI 对质量的定义与光速、能量之间建立起紧密的联系。本节引入 $\boldsymbol{M_{\mathrm{aps}}}$，代表{\textit{元粒子物质}}，用以描述与光速解耦的纯物质。
	
	\medskip
	依托已知的束缚强子态静态 QCD 结果，我们以 $M_{\mathrm{aps}}$ 为尺度重构了常规的能量—质量框架，并提出一个独特的观点：{\textit{囚禁生成物质}}。这一想法与“囚禁对应于能量积累”的传统解释存在根本差别。
	
	\medskip
	本节研究宇宙物质如何形成，以及宇宙光速 $c(t)$ 从强子化后等离子体阶段到 CMB 时代的逐渐上升。本工作不探讨强子化本身由 QCD 驱动的内在动力学机制。
	
	\medskip
	\textit{\textbf{注：}在我们的模型中，物质即{\textit{元粒子物质}}（$M_{\mathrm{aps}}$），是一种独立于光速的内禀物理实体。}
	
	
	\medskip
	\bigskip
	\textbf{关于新 SI 体系质量定义的一点商榷}\,\,\,
	2018--2019 年国际单位制（SI）的修订以普朗克常数 $h$ 重新定义了千克，旨在把完整的单位体系建立在一组极简的量子化基本物理常数之上。这一修订大幅提升了计量标准的长期稳定性与测量准确度。
	
	\medskip
	然而，这一框架模糊了内禀物质与能量之间的概念边界。官方标准把质量的定义完全锚定在能量框架之内，从而把质量仅仅当作能量的一种派生表现。相比之下，传统物理教育中的静质量概念能提供清晰的物理直觉。
	
	\medskip
	在作者看来，仅仅为了追求量纲基础方面的{\textit{极简主义}}，就把质量在基本层面定义为一种{\textit{能量等价形式}}，是不恰当的。更合理的做法是：在 SI 标准的概念层面保持清晰的区分，而在高精度测量层面采用现行的量子化方案。
	
	
	\medskip
	\bigskip
	\textbf{元粒子物质}\,\,\,
	{\textit{arche}} 一词源自古希腊语，意为起源、开端与根本原理。与{\textit{粒子物质}}合在一起，$M_{\mathrm{aps}}$ 表示宇宙粒子的原初基本物质。
	
	\medskip
	在现代物理学中，电子的静止属性常以 MeV 给出，而这是静能量的单位，并非内禀质量的单位。真正的静止惯性物质量，须以该能量除以 $c^2$ 求得。
	
	\medskip
	所提出的量 $M_{\mathrm{aps}}$ 以电子静止惯性质量 $M_\mathrm{e}$ 作为其标定基准，从而把物质概念与能量转换清楚地区分开来。
	
	
	
	
	\medskip
	\bigskip
	\textbf{QCD 中的强子化囚禁能}\,\,\,
	在量子色动力学（QCD）中，质量是以恒定光速所定义的能量来描述的。质子的静能约为 $m_p c^2 \approx 938.3$ MeV，它由四个部分组成。
	\[
	E_{p} = m_p c^2 = (2m_u + m_d)c^2 + T_{\rm quark} + E_{\rm gluon} + V_{\rm potential} \,\,.   \eqno(9.1)
	\]
	
	\smallskip
	第一部分来自两个上夸克与一个下夸克的贡献。其三个价夸克（$uud$）的静能之和为
	\[
	(2m_u + m_d)c^2 \approx 8.8 \text{ MeV} \,\,.   \eqno(9.2)
	\]
	
	\smallskip
	这部分夸克静能占质子总能量不足 $1\%$。在 QCD 中，其余近 $99\%$ 都源自囚禁，即强子化能。因此，质子质量主要由被囚禁的场能贡献，而非夸克的静质量。详细的能量构成表述如下：
	\begin{itemize}
		\item $T_{\rm quark}$ 表示三个被囚禁夸克的动能。这好比三个台球在强子内部以相对论性速度剧烈运动，该项约为 $600 \text{ MeV}$。
		
		\item $E_{\rm gluon}$ 表示强子内部的胶子场能。胶子如同黏度极高的胶水，在强子内部快速运动，维持强子袋的体积，其能量约为 $340 \text{ MeV}$。
		
		\item $V_{\rm potential}$ 指由一阶胶子交换产生的微扰势能，约为 $-25 \text{ MeV}$。
	\end{itemize}
	
	\medskip
	在上述 QCD 模型中，质子的总能量 $E_{p}$ 为 938.3 MeV，自由电子的能量 $E_{e}$ 为 0.511 MeV。前者为后者的 $\boldsymbol{1836.15}$ 倍，这就是众所周知的质子与电子质量之比。
	
	
	
	\medskip
	\bigskip
	\textbf{囚禁生成物质}\,\,\,
	在 QCD 中，囚禁{\textit{捕获}}能量~\cite{Peskin:1995}；在本工作中，{\textit{囚禁生成物质}}。这是一个根本性的区别。在光速恒定的假设下，能量与物质几乎可以互换；而在变 $c$ 框架中，这两个概念必须严格区分。
	
	\medskip
	囚禁并非一个简单的储能过程。相反，囚禁过程本身就是一个{\textit{成物过程}}。宇宙物质的大部分源自强子化，因此囚禁是宇宙物质生成的主导机制。
	
	\medskip
	在 CERN、费米实验室、BNL 与 SLAC 进行的高能实验所研究的，是处于高度激发态的质子。所观测到的夸克与胶子行为，是极端激发下的能量驱动响应，未必反映强子内部真实的基态物质结构。
	
	\medskip
	质子的内部结构从未被直接观测过。因此，粒子碰撞中所探测到的能量特征，不能等同于驻留于强子内部的内禀物质本性。
	
	\medskip
	\textit{\textbf{注：}在这一模型中，强子化等同于{\textit{物质化}}。强子好比宇宙中的不动产——固定的物质资产——而光速则充当宇宙的“能量—物质”汇率。宇宙物质的大部分，是在极低 $c(t)$ 的早期强子化过程中形成的，对应着这一低汇率。}
	
	\medskip
	QCD 把宇宙质量归因于恒定光速下、在超高温中“冻结”的能量；相比之下，本工作把宇宙物质追溯到早期低 $c$ 值下由囚禁驱动的物质生成。这两种物理图景存在根本差异。
	
	\medskip
	一旦形成，被囚禁的物质 $M_{\mathrm{aps}}$ 便永久作为物质而非能量存在。任何向能量的转化都按当代汇率 $c(t)^2$ 进行，因此一定量物质的能量当量随 $c(t)$ 的增大而上升。由此，$M_{\mathrm{aps}}$ 守恒而能量不守恒，构成本模型的一条核心公设。
	
	\medskip
	对质子而言，QCD 把囚禁量化为约 938.3 MeV 的能量。作为对照，本工作把被囚禁的内容描述为约 $\boldsymbol{1836.15}$ 个电子—$M_{\mathrm{aps}}$ 单位的物质，从而把囚禁效应直接联系到物质的量级尺度。
	
	
	\medskip
	\bigskip
	\textbf{光子散射与原初温度}\,\,\,
	标准 QCD 宇宙学假设一个炎热的原初宇宙，却带有两个未解的缺陷：强烈的光子散射阻碍强子化，而宇宙膨胀无法解释早期热分布的均匀性。
	
	\medskip
	增长-$c(t)$ 的 ZLZT 态始于近乎零的温度与光速，从而抹除了强烈的原初光子散射，消除了强子化的障碍。
	
	\medskip
	强子化俘获周围能量以形成 $M_{\mathrm{aps}}$ 物质。这一耗能过程使宇宙冷却，并自然产生均匀的早期热场。
	
	
	\medskip
	\bigskip
	\textbf{ZLZT 拟设中的温度与光速阈值}\,\,\,
	由于光速 $c$ 普遍被当作常数，关于强子到夸克—胶子等离子体（QGP）转变的标准量子色动力学（QCD）研究，只考察解囚禁临界点的温度依赖性。尚无文献讨论可变光速会如何改变这一临界阈值。
	
	\medskip
	氢的电离是由 QED 支配的两体解离过程，服从玻尔标度律，从而给出临界电离温度对光速的二次依赖 $\boldsymbol{T_{\rm ion} \propto c^2}$。
	
	\medskip
	为构建自洽的变-$c(t)$ 宇宙学模型，我们采用一个简单的幂律关系，以区分 QED 与 QCD 相变。 
	
	\medskip
	在强子内部，夸克与虚胶子表现出强烈的多体集体相互作用。我们因此假设 QCD 解囚禁温度遵循光速的三次标度 $\boldsymbol{T_{\rm QGP} \propto c^3}$。这一关系只是 ZLZT 框架的一个唯象假设，并非固定 $c$ 的 QCD 的标准结果。
	
	\medskip
	在这一假设下，若光速降至今日值的 $1/7137.66$，即 $c(t) \approx 42\,\mathrm{km/s}$——一个趋近于零的极低值——则氢的电离温度按比例降至 $1.96\ \mathrm{K}$。换言之，在这一降低后的光速阈值处，任何高于 $1.96\ \mathrm{K}$ 的温度都会把中性氢原子解离为电磁等离子体。
	
	\medskip
	若把这一三次标度规则应用于 QCD 解囚禁，则在 ZLZT 早期阶段，当前 $10^{12}\,\mathrm{K}$ 的 QGP 转变温度将按 $1/7137.66^3$ 的因子降低。由此得到约 $2.75\,\mathrm{K}$ 的残余温度。 
	
	\medskip
	由这一逐步推断可知：若在 QGP 阈值时期 $\boldsymbol{z_{\star}}$ 宇宙光速降到 $42\,\mathrm{km/s}$（或 $21\,\mathrm{km/s}$、$14\,\mathrm{km/s}$）以下，而周围温度保持在 $2.75\,\mathrm{K}$ 附近，则一切强子结构完全解离为自由夸克与胶子，形成一个不含束缚强子的纯夸克—胶子宇宙。 
	
	\medskip
	这便构成我们对原初零光速零温度（ZLZT）态之光速与温度阈值的推断。当温度固定于 $2.75\,\mathrm{K}$，而光速从 $21\,\mathrm{km/s}$ 升至 $42\,\mathrm{km/s}$ 或 $63\,\mathrm{km/s}$ 时，强子化过程逐渐加速，从夸克—胶子等离子体到强子的转变愈发频繁。相应地，“囚禁生成物质”过程也变得日益活跃。
	
	
	
	\medskip
	\bigskip
	\textbf{LEM-SCDF 对 QGP 时期的不适用性}\,\,\,
	麦克斯韦方程组只描述电磁（U(1)）相互作用，其适用范围有限。相比之下，夸克—胶子等离子体（QGP）时期主要由强力（SU(3) 色相互作用）支配。这就划出了一条明确的物理边界：经典电动力学无法完整刻画 QGP 内部出现的行为。
	
	\medskip
	此外，所有既有的 QGP 理论与实验研究都采用真空光速 $c_0$ 恒定的假设，这与 LEM-SCDF 所提出的变 $c$ 前提存在尖锐冲突。 
	
	\medskip
	LEM-SCDF 框架被设定为从强子化后等离子体时期起、历经 CMB 退耦时期直至今日都成立；但我们明确把 QGP 与强子化转变排除在 LEM-SCDF 的适用范围之外。
	
	
	
	\medskip
	\bigskip
	\textbf{物质成团塑造 $\boldsymbol{c(t)}$ 的动力学}\,\,\,
	本小节提出一个尝试性的方向性猜想：在 CMB 时期之前，宇宙光速 $c(t)$ 随时间的增长源自宇宙物质的局域成团。
	
	\medskip
	强子化完成之后，原本均匀弥散的物质开始在局域空间区域内聚集，束缚强子物质逐渐变得更致密。这就是我们此处所讨论的物质成团过程。
	
	\medskip
	这种结构上的差异打破了宇宙物质早期的均匀分布，也改变了由 LEM-SCDF 场所调控的整体电磁背景。 
	
	\medskip
	随着物质不断积聚为致密的局域结构，宇宙物质分布的空间不均匀性加剧。由 LEM-SCDF 协调的背景电磁参数同步演化，进而推动宇宙 $c(t)$ 的逐渐增长。
	
	\medskip
	这就形成一个自我强化的反馈回路。为给这一思路提供直观的数学参照，我们采用统计方法，构建一个简化的二维模型，以量化物质的成团程度。
	
	\bigskip
	考虑一个面积为 $A$ 的二维正方形区域 $\Omega$。设 $\rho(x,y)$ 为平面上的面密度分布，总物质量守恒，即
	\[
	S = \iint_\Omega \rho(x,y)\,\mathrm{d}x\mathrm{d}y = \mathrm{constant} \,\,.   \eqno(9.3)
	\]
	
	\smallskip
	全局平均密度为
	\[
	\bar{\rho} = \frac{S}{A} \,\,.   \eqno(9.4)
	\]
	
	\smallskip
	对完全均匀的分布，$\rho(x,y) \equiv \bar{\rho}$ 在 $\Omega$ 内处处成立，这对应我们所假设的强子化后早期完全均匀的等离子体态。
	
	\medskip
	当物质在总量不变的前提下逐渐聚集成局域团块时，密度场变得不均匀，某些区域更致密，另一些区域更稀薄。
	
	\medskip
	为量化空间不均匀性与弥散程度的大小，我们引入物质密度场的方差：
	\[
	\sigma^2_\rho = \frac{1}{A}\iint_\Omega \big[\rho(x,y) - \bar{\rho}\big]^2 \mathrm{d}x\mathrm{d}y \,\,.   \eqno(9.5)
	\]
	
	\smallskip
	在早期均匀态下，$\sigma^2_\rho = 0$。 
	
	\medskip
	随着物质成团、分布愈发不均，$\sigma^2_\rho$ 单调上升。我们定义一个无量纲的变异系数
	$$
	C_\rho = \frac{\sigma_\rho}{\bar{\rho}}
	= \frac{1}{\bar{\rho}} \sqrt{\frac{1}{A}\iint_\Omega \big[\rho(x,y) - \bar{\rho}\big]^2 \mathrm{d}x\mathrm{d}y} \,\,.   \eqno(9.6)
	$$
	
	\smallskip
	以便在不同宇宙演化阶段之间作交叉比较。$C_\rho$ 的值越大，对应物质的空间分化越强，这与我们猜想中所述的 $c(t)$ 上升趋势一致。
	
	\bigskip
	从信息论与统计熵的角度，我们把归一化概率密度表述为
	\[
	p(x,y) = \frac{\rho(x,y)}{S}   \,\,,   \eqno(9.7)
	\]
	
	\smallskip
	它满足归一化条件 $\iint_\Omega p(x,y)\,\mathrm{d}x\mathrm{d}y = 1$。刻画物质分布有序程度的信息熵为
	\[
	H = -\iint_\Omega p(x,y)\,\ln p(x,y)\,\mathrm{d}x\mathrm{d}y   \,\,.   \eqno(9.8)
	\]
	
	\smallskip
	在总物质量与几何边界固定的条件下，完全均匀的分布给出最大熵 $H_{\max} = \ln A$。当物质组装成团簇结构时，系统变得更加有序，熵随之减小。我们据此定义熵减
	\[
	\Delta H = H_{\max} - H   \,\,,   \eqno(9.9)
	\]
	
	\smallskip
	作为物质聚集与空间不均匀性的直接度量。 
	
	\medskip
	合起来看，密度方差 $\sigma^2_\rho$ 与熵减 $\Delta H$ 构成对等离子体时期物质弥散与成团的一致双重统计描述。
	
	\bigskip
	该模型只作为一项启发性的辅助说明，用以澄清物质不均匀性与 $c(t)$ 演化趋势之间的相关性。把物质成团与 LEM-SCDF 的背景电磁参数联系起来的完整数学构造与严格物理证明，留待未来研究。
	
	\medskip
	\textit{\textbf{注：}在这一图景中，宇宙光速 \(c(t)\) 是宇宙的一种统计集体属性——正如气体温度不是任何单个分子的属性，而是它们集体运动的统计平均。每一颗恒星、每一毫克物质、每一焦耳能量，都通过其电磁构型参与 LEM-SCDF 背景的全局响应。随着物质凝聚，这些源的统计分布发生移动，\(c(t)\) 也随之演化。简言之，\(c(t)\) 之于宇宙，正如温度之于气体：它是一个宏观统计量，而非形而上的常数。}
	
	
	\medskip
	\bigskip
	\textbf{中性原子的黎明：迈向 CMB 退耦}\,\,\,
	随着宇宙物质局域凝聚的持续，光速 $c(t)$ 稳步上升。光速的这种增大直接增强了质子与电子之间的库仑相互作用。于是，由 $c(t)$ 与精细结构常数共同决定的特征轨道半径，从极小的初值逐渐扩张，直至大到足以容纳电子绕质子运行——达到第 7 节所导出的今日玻尔半径的约 $1/1100$。
	
	\medskip
	在这一增强的电磁束缚力之下，电子终于获得了稳定绕质子运行所需的物理条件。尽管在这一早期阶段，正在形成的原子结构仍易被再次电离，但随着 $c(t)$ 增长，维持稳定束缚态的概率持续上升。随着越来越多的电子成功与质子结合、原初氢原子大量累积，宇宙开始脱离以等离子体为主导的状态。 
	
	\medskip
	在从强子化后等离子体时期（$\boldsymbol{z_{\star}}$）到 CMB 退耦时期（$\boldsymbol{z_{0}}$）的宇宙演化中，光速从低于 $42\ \mathrm{km/s}$ 逐渐增至今日 $c_0$ 的 $1/1100$（$272\ \mathrm{km/s}$）。在早期宇宙光速如此之低的假设下，这一演化过程可能跨越 $4$--$5\ \mathrm{Gyr}$，远长于标准 $\Lambda$CDM 模型预言的 $360$--$370\ \mathrm{kyr}$。该跨度由 $c(t)$ 曲线作简单的幂律估计得到，覆盖从 $c_0$ 的 $1/7137.66$ 升至 $1/1100$ 的过程。 
	
	\medskip
	中性氢的这种持续形成，最终引导宇宙演化走向 CMB 退耦时期，开启一个由爱因斯坦场方程支配的全新宇宙疆域。
	
	
	
	
	
	
	
	\section*{
		\begin{center}
			10. 宇宙学红移：一个未被回答的\\
			因果缺口及其物理闭合		 
		\end{center}
	}
	\medskip
	\noindent
	
	整整一个世纪以来，现代宇宙学始终未能处理一个根本性的逻辑缺口。本节将直接面对这一缺口。
	
	\medskip
	

	\bigskip
	\textbf{我们的诘问}\,\,\,
	我们从第 7 节的一个具体结果出发。在那里，我们重新审视了光子频率、波长与光速的行为。考察一个 $\nu_{n0}=1.0\times10^{15}\,\mathrm{Hz}$ 的远古光子，它在红移 $1+z=4.0$ 处发出。经过近一百亿年在星系际空间中的传播，宇宙学红移已把它的频率在今日降至 $2.5\times10^{14}\,\mathrm{Hz}$。
	
	\medskip
	我们要向宇宙学家提出一个直接的问题：\emph{“这个光子早已离开它的宿主星系，在近乎空无一物的空间中旅行了一百亿年。它凭什么如此顺从、如此严格地遵循规则，知道自己必须随宇宙膨胀而红移？是谁通知了光子？又是如何通知的？”}
	
	\medskip
	标准答案我们都很熟悉。一种说法是：空间本身在膨胀，把光子的波长拉长，正如气球被吹大时，画在气球上的波纹会随之伸展。另一种说法是：在 FRW 度规中，标度因子 $a(t)$ 进入零测地线方程，频率自然按 $a(t)^{-1}$ 标度；计算既已给出红移，便无需再追问机制。
	
	\medskip
	然而，这两种回答都没有直接触及光子红移背后的物理机制。它们要么是几何类比，要么是数学上的重新表述，实际上遮蔽了底层的因果问题。
	
	\medskip
	让我们回到诘问本身。这条逻辑链条异常清晰：
	\begin{itemize}
		\item \textbf{前提\_1}：光子早已脱离一切物质源——它的发射星系、大质量天体以及局域引力场。
		
		\item \textbf{前提\_2}：在飞行途中，当宇宙尺度膨胀时，光子的频率严格按比例红移。
		
		\item \textbf{前提\_3}：若飞行中的光子与宇宙其余部分已无任何联系，它又如何\emph{知道}宇宙已经膨胀？又如何\emph{据此}调整自己的频率？
		
		\item \textbf{前提\_4}：若光子在整段飞行中始终与宇宙保持持续接触，以便即时接收膨胀信息并作出红移响应，那么这一接触机制的传播速度，岂非\emph{必须}远大于光速本身？
		
		\item \textbf{前提\_5}：这一接触所携带的信息，是关于“宇宙膨胀尺度”的直接指令吗？这似乎不可能。尺度膨胀这一几何量，与光子的电磁振荡之间并无内在的物理联系。
		
		\item \textbf{前提\_6}：那么这一接触是否改为向每个光子分别下达各自的频率调整指令？这同样不可能。全宇宙的光子拥有极大范围的不同频率，不存在一条能告诉每个光子究竟该偏移多少的普适指令。
		
		\item \textbf{前提\_7}：在排除其他一切可能之后，\emph{您难道不认同}只剩下唯一一个自洽的解答吗？一个全局场所传递的底层信号，只能是光自身的本征属性——即真空光速 $c(t)$ 的全局同步变化。不存在多种备选答案。
	\end{itemize}
	
	\smallskip
	在我们看来，这一现存的逻辑缺口，构成标准红移解读中一个真正的\textbf{\textit{因果赤字}}。
	
	\medskip
	这些疑虑也并不新鲜。宇宙学红移背后的物理机制此前已受到质疑：费曼（1960）~\cite{Feynman:2011} 指出，它不是多普勒频移，而是光子在旅途中自身发生的变化；戴维斯（1995）~\cite{Davies:1995} 进一步指出了一个持续的因果赤字，认为仅把红移归因于空间膨胀，是以几何描述取代了物理机制。 
	
	\medskip
	\bigskip
	\textbf{一种由历史传授的方法}\,\,\,
	物理学史告诉我们：从观测表象中推断既有理论存在因果缺口，并由此揭示出一个尚未被发现的全局物理实体，是一种强大而富有成效的思想方法——它在科学史的关键时刻曾一再奏效。
	
	\medskip
	爱因斯坦面对惯性质量与引力质量惊人的等价性，拒绝把它当作巧合；他断定引力不可能是超距作用力，而必然是时空几何本身的表现——由此发现了广义相对论。 
	
	\medskip
	奥斯特首先证明了电流能够产生磁场。法拉第借这一对称性猜想其逆效应：随时间变化的磁场应当感应出电场，并以实验验证了这一现象，从而建立了电磁感应定律。
	
	\medskip
	在这两个事例中，关键的一步都不是拟合数据，而是认识到既有框架缺少一个物理作用者——然后要求这个作用者必须存在。
	
	
	
	\medskip
	\bigskip
	\textbf{唯一因果上成立的结论}\,\,\,
	我们在此运用同样的推理。唯一因果上成立的结论是：必定存在某种机制——其传播远快于光，且覆盖整个宇宙——在光子的整段旅程中持续地通知并调制它。 
	
	\medskip
	宇宙必定充满着一个全局的、极其快速的同步场，每一个光子在整条传播路径上都与之保持耦合。该场响应宇宙平均密度的全局稀释与光速的普遍演化，并依据膨胀着的背景持续调整光子的频率。
	
	\medskip
	换言之，宇宙学红移的物理解读，要求存在一个传播速度远超光速的普适背景场。而这恰恰正是 LEM--SCDF 暗场所承担的核心功能。
	
	\medskip
	式 (8.10) 中 LEM 传播速度的定义并非唯一；我们采用普朗克密度与弗里德曼密度作为启发性的标度因子，为 LEM-SCDF 场定义一个足够大的传播速度。其他在数学上自洽的形式同样可以构造。然而，不可选项是 LEM-SCDF 场本身的物理必要性：它的传播速度必须远大于 $c(t)$。
	
	\medskip
	这是逻辑上也是物理上的必然：一个光子在离开一切局域源很久之后，仍继续遵循一种无法仅由局域相互作用解释的频率演化。在这一观测行为背后，必定存在某种远快于光、并与整个宇宙保持持续接触的物理机制。
	
	
	\medskip
	\bigskip
	\textbf{宇宙学红移就是宇宙光速膨胀}\,\,\,
	从我们提出的上述诘问，到我们所达到的结论，浮现出一个\emph{不可回避}的真相：宇宙学红移就是全局光速 $c(t)$ 的膨胀！
	
	\medskip
	\emph{对一个一百亿年前发出的远古光子而言，其漫长旅途所累积的红移，本身就是一百亿年真空光速演化的记录。}
	
	\medskip
	近一个世纪以来，我们观测手段的固有局限使我们视而不见：远古光子及其红移是可以观测的，但对于早期宇宙真实的光速，我们掌握的直接证据为\textbf{\textit{零}}。如今，在全新的 LEM-SCDF 范式之下，我们可以向世界宣告：
	
	\medskip
	\emph{宇宙空间膨胀根本上就是光速膨胀。红移仅仅是全局增长-$c(t)$ 变化可观测的印记；二者是同一个统一物理过程的两副面孔。}
	
	\medskip
	相对于宇宙演化的时标，四个世纪微不足道。自伽利略（1638）与 R\o mer（1676）首次指出光以有限速度传播以来，一代又一代的测量不断强化着我们的信念：光速是一个普适常数，由造物主设定，永不改变！
	
	\medskip
	而我们现在开始意识到：对宇宙这位巨人而言，四百年不过一次心跳。如果光速并非由造物主钦定，而是由宇宙中每一毫克尘埃、每一焦耳能量共同塑造——那其实是一个更和谐、也更美的宇宙。
	
	
	
	\medskip
	\bigskip
	\textbf{作为一个相干物理整体的宇宙}\,\,\,
	在传统物理学中，没有任何内禀机制把所有宇宙质能普遍地联系起来，于是宇宙只不过是一大堆彼此孤立的局域沙粒的集合，缺乏全局协调。然而每一粒沙都默默恪守着它那两个神圣参数——$c$ 与 $G$——纹丝不动，仿佛出自造物主的钦定。
	
	\medskip
	LEM-SCDF 场提供了一种物理机制，通过同步调制把每一个粒子与光子瞬时联系起来。它把全部耦合质能的分布编码于自身，并把大尺度的宇宙变化均匀地传递给每一个参与组分。于是，宇宙成为一个动力学上相干的物理整体，由真实而连续的场相互作用统一起来。
	
	\medskip
	以往关于宇宙的定义含混而不精确。随着 LEM-SCDF 的引入，这一概念变得严格而无歧义：每一个基本的质量与能量单元都参与塑造 LEM-SCDF 的演化；而所有此类质能的完整集合，便属于我们所定义的宇宙——无论它分布在何处，也无论它距任何参考中心有多远。
	
	\medskip
	在这一框架内，宇宙膨胀获得了严格而明确的物理含义。此前存在一个众所周知的逻辑悖论：如果宇宙在空间上已经无限，那么弗里德曼膨胀描述的究竟是什么？无限尺度的膨胀？在我们的新图景中，答案很清楚：膨胀描述的是 LEM-SCDF 耦合系统内部，分布区域的\textit{统计}扩大以及密度的相应\textit{统计}稀释。
	
	\bigskip
	在新图景中，每一个光子都经受着一种主动而持续的全局调制，而施加调制的，正是那个唯一连接宇宙万物的实体。这不是理论上的装饰性补充；它是红移现象本身所要求的因果闭合。
	
	
	
	\section*{
		\begin{center}
			11. 关于微小演化 $\boldsymbol{c(t)}$ 实验\\
			可探测性的讨论
		\end{center}
	}
	\medskip
	\noindent
	
	自 1983 年 SI 修订以来，真空光速被固定为 $c = \texttt{299,792,458} \text{\,m/s}$，使 $c$ 不再是可独立测量的物理量。$c$ 的恒定性由此成为几乎所有计量仪器之设计与运行的基本公理。
	
	\medskip
	按我们增长-$c(t)$ 模型的推导，$c(t)$ 的年相对变化率处于 $10^{-11}$ 量级。如此微小的偏差，至少超出人类现有全部测量仪器与实验技术的精度极限两到三个数量级。 
	
	\medskip
	只有顶尖的实验物理学家、创新的测量方案与最先进的装置，才有可能实现对微小演化 $c(t)$ 的定量探测。本节所作的简要讨论，并不构成任何可实际实施的实验提案。

	
	
	\bigskip
	\textbf{光速的历史测量}\,\,\,
	1987 年 NBS 报告~\cite{NBS:1987} 记载的历史干涉测量工作，给出的相对测量上限为 $3.67\times 10^{-9}$。一系列历史实验通过彼此不同的物理机制，逐步提升了 $c$ 的测量精度，简述如下：
	\begin{itemize}
		\item 菲索旋转齿轮法（1849）：地面直接测量光的飞行时间，精度约 $10^{-4}$，是首次在地面上对 $c$ 作出的定量测定。

		\item 傅科—迈克耳孙旋转镜方案（傅科 1850，迈克耳孙 1926 年八面镜改进）：通过优化光路布局抑制机械系统误差，把测量精度提升至约 $10^{-6}$。	
		\item 微波谐振腔（1946--1970 年代）：通过标定谐振频率与腔体几何，由恒等式 $c=\lambda f$ 算出 $c$；精度约 $10^{-9}$，是 1983 年 SI 重新定义之前的主要计量标准。	
		
		\item 稳频激光干涉测量（1972 年 NBS 实验）：借助原子钟与干涉技术，分别测量稳频激光的频率与真空波长，再合并结果计算 $c=\lambda\nu$；该实验实现了 \(6\times10^{-10}\) 的频率测量不确定度，较以往 $c$ 的测定提升了 100 倍，并为 1983 年第 17 届 CGPM 采用 \(c=299\,792\,458\ \mathrm{m/s}\) 的固定定义奠定了实验基础。
	\end{itemize}
	
	
	\bigskip
	\textbf{增长-$\boldsymbol{c(t)}$ 模型变化率的估计}\,\,\,
	在紧凑模型中，$C(t)$ 的导数由式 (5.14) 给出。取当前时刻 $t=0$，并代入式 (5.20) 中的参数 $k$ 与 $\beta$，得到 $\dot{C}(t)=0.03145$，对应 $3.145 \times 10^{-11}\,\mathrm{yr}^{-1}$。折算到 $c_0$，这给出 $9.43\,\mathrm{mm\!\cdot\!s^{-1}\!\cdot\!yr^{-1}}$ 的长期变化。
	
	\medskip
	对超精细模型，我们采用 \texttt{Planck18} 标准红移数据集，并用式 (6.1) 中的 $z_1$ 进行计算，取最近 $10^4$ 年的平均值。这一设置给出每 Gyr $6.89\%$ 的相对增长率，等价于 $6.89 \times 10^{-11}\ \mathrm{yr}^{-1}$。归一化到 $c_0$，对应每年 $20.7\ \mathrm{mm\!\cdot\!s^{-1}\!\cdot\!yr^{-1}}$ 的变化。
	
	\medskip
	若以日尺度评估该速率，把上述年相对变化率除以 $365$，便得到介于 $[0.86,\,1.89]\times 10^{-13}\ \mathrm{day}^{-1}$ 区间的每日相对增长率，折算为 $c$ 的每日增量 $[+26,\,+57]\, \mu\mathrm{m/s}$。
	
	\medskip
	\textit{\textbf{注：}除上述精度约束之外，另一个关键顾虑在于避免循环论证。由于计量标准定义 \texttt{Metre} 与光子波长 $\lambda$ 都建立在光速不变这一基本公理之上，必须另行设计全新的测量框架。}
	
	
	\medskip
	\bigskip
	\textbf{探索微小 $\boldsymbol{c(t)}$ 演化的 HEPS 装置条件}\,\,\,
	位于北京怀柔的高能光子源（HEPS）是一台尖端的第四代同步辐射装置，可提供高精度、高分辨率的实验条件~\cite{HEPS:2022}。据公开报道，HEPS 的建设总投入约为 7 亿美元。其主要性能优势归纳如下：
	
	\begin{itemize}
		\item HEPS 拥有 6\,GeV 的高能电子储存环，可产生高纯度的硬 X 射线束，用于物理参数的稳定标定。
		
		\item 它支持原子尺度的电荷密度成像、单纳米颗粒成像与高相干衍射，突破了传统条纹计数测量方法的极限。
		
		\item 其全无油超高真空系统可维持低于 $1\times10^{-9}\,\mathrm{Pa}$ 的本底压强。该配置消除了空气的折射干扰，为与光速相关的实验提供纯净的真空环境。
		
		\item 该装置的纳米探针可实现 9.5--12.5\,nm 的聚焦光斑尺寸，为光路与晶格几何的精确标定提供原子级空间分辨率。
		
		\item 光束线可提供皮秒级时间分辨率（$10^{-12}\,\mathrm{s}$），配合超快探测器可提升至 100\,fs，用于对微弱瞬态物理信号作高精度计时。
		
		\item HEPS 采用原子级抛光的柯克帕特里克—贝兹（KB）折转聚焦镜，其表面平整度达亚埃量级（$<0.1\,\mathrm{nm}$），把光路畸变降至最低。
	\end{itemize}
	
	\medskip
	HEPS 的实验群体是否会采纳增长-$c(t)$ 猜想，并进一步设计专门的实验方案，目前尚无定论。
	
	
	\medskip
	\bigskip
	\textbf{Mission-M 研究计划}\,\,\,
	在光速恒定假设占主导的现行物理范式之下，要为国家公共科研经费申请到专门用于搜寻 \(c(t)\) 时变性的实验，实际上不可行。因此，作者打算牵头发起一项名为 \textit{Mission-M} 的民间非政府研究计划，预算 7500 万美元，计划为期三十年。该计划围绕三条核心原则组织：以空前的精度\textit{测量}光速，若探测到任何变化则\textit{建模}其底层物理，以及无论最终结果如何都要\textit{探秘}未知。
	
	\medskip
	第一个十年用于理论分析、实验方案设计，以及新型仪器的发明与研制。其余二十年则投入持续的观测活动。总体目标是把光速的测量精度推进到 \(10^{-14}\) 至 \(10^{-15}\) 水平，相当于相对 1983 年 SI 基准提升 \(10^5\) 至 \(10^6\) 倍，并将在国内组建一支高水准的实验物理学家团队。
	
	\medskip
	可以设想两种结局。即便测量证实光速在这一提高后的精度下依然恒定，该实验仍将构成一项有价值的科学成果，堪与 1887 年具有里程碑意义的迈克耳孙--莫雷实验相比——后者排除了静态以太与以太风的存在。反之，若探测到接近 \(10^{-13}\) 的每日相对变化，则将为基础物理研究开辟广阔的新途径。

	
	
	\section*{
		\begin{center}
			12. 以增长-$\boldsymbol{c(t)}$ 框架\\
			重新审视若干悬而未决的问题
		\end{center}
	}
	\medskip
	\noindent
		
	本节将在该增长-$c(t)$ 模型下重新审视若干悬而未决的问题，包括 DESI DR2 对动态暗能量的偏好、宇宙光子预算危机、早期星系恒星质量张力、星族 III 恒星时标冲突、早期太阳锂丰度之谜、高级智慧生命出现得如此之晚，以及其他若干开放问题。
	
	\medskip
	\bigskip
	\textbf{DESI DR2 张力的另一种解释}\,\,\,
	在标准 \(\Lambda\)CDM 范式中，暗能量充当驱动宇宙加速的均匀背景流体~\cite{Riess:1998}。然而，近期的 DESI DR2 BAO 测量结合 Ia 型超新星与 CMB 数据，揭示出对纯宇宙学常数的持续偏离~\cite{Gu:2025}。这一偏离并非孤立发现。正如 Abreu 与 Turner~\cite{AbreuTurner:2026} 最近所强调的，DESI 数据暗示暗能量密度在 \(z \simeq 0.5\) 附近达到极大，这一行为难以与简单的宇宙学常数的图像相协调，并提示存在新物理。暗能量状态方程 \(w(z)\) 被发现随红移变化：在高红移（\(z\sim 0.75\)）处，数据偏向 \(w < -1\)（幽灵区）；在低红移（\(z\lesssim 0.2\)）处，数据偏向 \(w > -1\)（精质区）。这一穿越 \(w = -1\) 的行为对应“精质—幽灵 B”（quintom B）型，显著性超过 \(3\sigma\)，并给出贝叶斯证据 \(\Delta\ln E = 5.2 \pm 0.7\)，支持动态暗能量优于 $\Lambda$CDM。
	
	\medskip
	在我们增长-\(c(t)\) 框架内，这一观测图景可获得直截了当的物理解释。第 8 节的修正爱因斯坦场方程表明，宇宙膨胀与光速 \(c(t)\) 的增长直接耦合。随着物质聚集成致密团块，星系际体积变得稀薄，并通过 LEM-SCDF 场重塑全局电磁背景。按 \(G(t) \propto c(t)^{-2}\)，$c(t)$ 的稳步增长削弱了有效的长程引力耦合。关键在于，高红移处 \(w(z) < -1\)、低红移处 \(w(z) > -1\) 的穿越行为并非什么奇异物理，而是 $G(t)$ 随 $c(t)$ 增长而单调下降的自然结果：它改变了膨胀历史，而当用标准的 \(w_0\)-\(w_a\) 参数化来解读时，恰好给出 DESI DR2 中所见的幽灵穿越图样。在这一图景中，“动态暗能量”并非一种神秘流体，而是一种误认——它是一个引力耦合与光速由不变量 \(P = G(t)c(t)^2\) 锁合、共同演化的宇宙，所呈现出的宏观印记。

	
	\bigskip
	\noindent
	$\triangleright \triangleright \triangleright$\enspace
	\textit{\textbf{悖论：}暗能量是畏惧人类，从而主动回避太阳系吗？} 我仔细考察了近期 DESI DR2 的观测数据，并在标准宇宙学内部发现了一个根本性的悖论——据我们所知，这一悖论在文献中从未被严格面对过。这个悖论直指现代宇宙学的核心。直白地说：	
	\begin{itemize}
		\item 演化年龄完全相同的恒星系统（约 $t^\star = -4.67$ Gyr），仅因它们与我们观测视点的空间距离不同，便表现出截然不同的暗能量行为；		
		\item 在近处，在我们这个已有 46.7 亿年历史的太阳系内，暗能量的印记几乎不存在；
		\item 然而在遥远的宇宙环境中，结构类似、年龄相同（同样为 $t^\star \approx -4.67$ Gyr）的系统却表现出显著的暗能量效应，而且这种效应随距离增加而增强；
		\item 没有任何现有的 $\Lambda$CDM 机制能在物理上为这种不自然的宇宙空间偏置提供依据； 
		\item 相同的宇宙年龄，截然不同的物理。唯一的变量，是距地球与太阳系的距离。
	\end{itemize}
	
	于是我们再一次提出问题：暗能量真的作为宇宙的一个基本组分而存在吗？抑或它的表观行为，不过是我们自身观测位置与光行时间的镜像？
	
	\medskip
	不变量关系 $G(t)c(t)^2=P$ 或可化解这一不一致。被归因于暗能量的表观变化，可能自然而然地来自 \(c(t)\) 与 \(G(t)\) 宇宙演化的观测伪迹。我们并不声称要彻底排除宇宙学常数；我们只是指出，在增长-$c(t)$ 框架内存在一条通往宇宙加速的独立物理路径。那些挑战经典 $\Lambda$CDM 模型的 DESI DR2 观测，与这一演化图景完全相容，并且可能构成本工作所提出的 $G(t) \propto c(t)^{-2}$ 标度律的第一批观测证据。
	\enspace$\triangleleft \triangleleft \triangleleft$
	

	\medskip
	\textit{\textbf{注：}我们对暗物质框架持同样的审慎立场。早期 $c(t)$ 较小所带来的更强有效引力，可以部分解释高红移处的引力亏缺。然而，我们的机制无法化解低红移的 coma 星系团中所观测到的 400 倍质量失配~\cite{Zwicky:1933}。}


	\medskip
	\bigskip
	\textbf{重新解读 $\boldsymbol{1+z=25}$ 处的宇宙光子预算危机}\,\,\,
	标准 $\Lambda$CDM 建立在两条前提之上：热大爆炸冷却图景，以及 tSZ 热标度关系 $T(z)=T_0(1+z)$（$T_0=2.725\ \mathrm{K}$）。该关系预言 $1+z=25$ 处的星系际温度为 $68\,\mathrm{K}$，$1+z=7$ 处为 $19\,\mathrm{K}$，前者约为后者的 3.58 倍。
	
	\medskip
	按原子物理，能量高于 13.6\,eV 的光子会立即电离氢；而在较低的气体温度下，复合会急剧加速。在 $1+z=25$ 处，早期星系发射出大量与 91.2\,nm 电离阈值相匹配的远紫外光子。标准理论便以 $68\,\mathrm{K}$ 的温度论证复合很慢，因而丰富的紫外光子会在高红移处把星系际氢完全电离。
	
	\medskip
	这一预言与联合观测相冲突~\cite{Madau:2023}。JWST 在 $1+z=25$ 处探测到极强的紫外光子输出，而 $21\tinysp$cm 射电信号证明同一红移处仍有广泛存在的中性氢。这一矛盾即构成光子预算危机；标准宇宙学只有在 $1+z=7$ 附近，借助紫外产生率下降与复合加速才能化解它。 
	
	\medskip
	近期在标准范式内部化解这一危机的努力，恰好凸显了危机的严重程度。例如，Austin 等人~\cite{Austin:2025} 以及其他近期研究~\cite{Reionization:2025} 试图借助强反馈机制或 HII 成团来调和 JWST 观测与再电离。这些尝试虽提供了有价值的洞见，却往往需要精细调节的参数。相比之下，增长-\(c(t)\) 框架通过解耦 \(T_{\rm rad}\) 与 \(T_{\rm gas}\)，给出了一种更根本的解决方案——正如我们所论证的。
	
	\medskip
	$T(z)=T_0(1+z)$ 公式假设光速固定，这与我们的增长-$c(t)$ 模型冲突——后者在 $1+z=25$ 处有 $c(t)=c_0/25$，在 $1+z=7$ 处有 $c(t)=c_0/7$。我们把两个彼此无关的温度量区分开来：
	\begin{itemize}
		\item $T_{\rm rad}(z)$：由 tSZ 效应测得的远处光子的有效黑体温度；
		\item $T_{\rm gas}(z)$：支配电离与复合的弥散星系际氢的真实热力学温度。
	\end{itemize}
	
	标准宇宙学错误地把 $T_{\rm rad}(z)$ 与 $T_{\rm gas}(z)$ 等同起来；我们的框架则把它们彻底解耦。
	
	\medskip
	在我们的 ZLZT 冷原初态下，$1+z=25$ 与 $1+z=7$ 处的 $T_{\rm gas}(z)$ 都低于 $10\ \mathrm{K}$。超冷介质造成极高的复合率。每次光子碰撞之后，被电离的粒子迅速复合为中性氢，从而无需任何人为拟合参数，便自然地与全部观测数据相容。 
	
	\medskip
	另一个可能的原因是：$\Lambda$CDM 模型对 $1+z=25$ 时期恒星的质量与密度作出了错误估计，从而导致其紫外光子输出被严重高估。这一点将在下一小节分析。
	
	\medskip
	\textit{\textbf{注：}本文只聚焦于与再电离相关的狭义光子预算悖论，与全波段 EBL 问题无关。我们接受 tSZ 测量，但不接受把辐射温度与星系际气体温度等同起来这一未经证明的做法。}
	
	
	\medskip
	\bigskip
	\textbf{重新评估恒星质量与密度}\,\,\,
	在标准牛顿天体力学与传统宇宙学中，所有恒星与星系的质量标定根本上都依赖引力常数 $G_0$，而 $G_0$ 是在太阳系内局域标定的。标准宇宙学框架隐含地把 $G=G_0$ 视为在所有宇宙红移处都不变，其后关于恒星密度、半径与恒星演化的一切计算，都服从这一固定的标度规则。
	
	\medskip
	近期 JWST 观测在 $z\sim10$ 至 $z\sim25$ 的极高红移处识别出若干致密星系的候选体。在引力常数 $G_0$ 固定的假设下，这一高红移区间天体的估计平均密度接近甚至超过核饱和密度~\cite{Labbe:2023}。
	
	\medskip
	若密度超过这一饱和阈值，恒星内部便会点燃失控的核聚变，促使其迅速塌缩为黑洞。这一理论结果与空间望远镜所观测到的稳定、长寿的高红移系统相冲突。标准理论试图借助未经证实的物理机制来化解这一张力：超爱丁顿吸积，或大质量原初气体云的直接塌缩；然而这两种情形都未通过局域天文观测的验证。
	
	\medskip
	回到牛顿天体力学，行星、双星与位力化星系的轨道运动只约束组合乘积 $GM$。这一简并直接源自角动量守恒。对一个约化质量为 $\mu=\frac{Mm}{M+m}$ 的两体引力系统，轨道角动量取形式 $L=\mu\sqrt{G M a(1-e^2)}$，其中 $a$ 与 $e$ 分别表示轨道半长轴与偏心率。
	
	\medskip
	由于表达式中只出现 $GM$，轨道、角动量与光变曲线观测都无法把 $G$ 与中心质量 $M$ 分离开来。对角度动量守恒的封闭系统，任何具有相同 $GM$ 的 $(G,M)$ 组合都给出完全相同的轨道、周期与速度分布。
	
	\medskip
	我们的增长-$c(t)$ 模型假设引力常数 $G(t)$ 按 $G(t)\propto c(t)^{-2}$ 的关系随红移上升。这表明在早期高红移宇宙中，引力耦合远为更强。若保留轨道观测不变量 $G(t)M(t)$，则天体的内禀重子质量 $M(t)$ 与密度都将远低于标准天文模型的预言。 
	
	\medskip
	增长-$c(t)$ 框架的宇宙标度律打破了 $G$ 与 $M$ 之间的简并，把恒星密度分布推移到更温和、物理上自洽、且与重子物质的核饱和极限相容的区间。
	
	\medskip
	增长-$c(t)$ 模型为标准的 $\Lambda$CDM 分析框架引入了一个以往被忽视的基本标度自由度。未来工作需要把由玻尔半径导出的原子标度、原初核合成，以及双星轨道运动学结合起来作联合约束，以化解 $GM$ 简并，并建立适用于所有宇宙时期的修正质量标度关系。
	
	
	
	
	
	\medskip
	\bigskip
	\textbf{化解星族 III 恒星的时标冲突}\,\,\,
	在标准 $\Lambda$CDM 模型中，由于原初气体云缺乏金属谱线冷却，第一代无金属的原初恒星——星族 III 恒星——不可避免地以极大质量形成，介于十个到数百个太阳质量之间。标准 $\Lambda$CDM 模型预言，这些极大质量的星族 III 恒星在主序上仅停留两三百万年~\cite{Topping:2023}。
	
	\medskip
	这一经典预言与近期 JWST 在 $z\sim10$ 至 $z\sim25$ 高红移区间对天体的观测，产生了严重的时标冲突。JWST 已识别出在大爆炸之后仅数亿年便出现的、完全成形的大质量星系。若这些原初大质量恒星都在短短几百万年内燃尽核燃料并死亡，那么星系组装、超新星反馈与金属增丰就必须以不合物理的速度推进，才能产生所观测到的成熟星系。
	
	\medskip
	为缓解与 JWST 观测的这一冲突，宇宙学家提出了若干未经证实的特设物理假设，包括恒星核心内部的暗物质湮灭加热，以及强烈的顶重初始质量函数（IMF）分布。
	
	\medskip
	我们认为，标准 $\Lambda$CDM 模型中原初恒星寿命极短，根源在于一条未经检验的假设：引力常数 $G=G_0$ 在所有宇宙时期都不变。这一固定 $G$ 的先验，系统性地高估了高红移原初天体的重子质量。
	
	\medskip
	我们的增长-$c(t)$ 框架遵循 $G(t)\propto c(t)^{-2}$ 的关系，表明在早期宇宙的高红移处，引力耦合显著更强。在轨道可观测量守恒的约束下，早期增强的 $G(t)$ 必然降低原初恒星的内禀重子质量 $M(t)$。关键在于，这一质量降低效应压过了引力耦合的增强，从而大幅削弱了早期恒星的核心压缩。
	
	\medskip
	恒星物理确立了一条稳健的质量—光度关系 $L\propto M^{3\sim4}$，意味着恒星的能量消耗随质量超线性上升。经典研究已考察过恒星演化对引力常数取值的敏感性~\cite{Prather:1976}，其中表明 $G$ 的改变会直接改变核燃烧时标。内禀恒星质量的中等幅度降低——正如早期宇宙中较大的 $G(t)$ 所暗示的——会剧烈压低核燃烧速率，而核燃料总量却只是线性减少。
	
	\medskip
	因此，增长-$c(t)$ 框架或许无需借助奇异加热机制或精细调节的初始质量分布，便能化解这一严重的时标冲突。通过打破 $GM$ 简并并修正高红移恒星质量标度，早期恒星演化时标回到物理上自洽的区间，从而为在极早期宇宙中所观测到的成熟星系结构，提供一种优雅而自洽的解释。
	
	
	

	\medskip
	\bigskip
	\textbf{早期太阳锂丰度之谜}\,\,\,
	标准恒星演化理论假设光速 $c$ 恒定。基于这一前提，它们为早期太阳预言一个确定的原初锂丰度。然而，高精度光谱观测早已揭示出一个显著差异：太阳表面的锂含量仅为理论预言的三分之一。这一众所周知的不一致，被称为长期悬而未决的太阳锂之谜。
	
	\medskip
	为修补这一问题，既有理论附加了额外的假设，包括更强的内部混合、早期质量损失，或锂藏于行星之中~\cite{Korn:2019}。然而这些想法都没有直接证据，也未能提供统一而自洽的物理图景。
	
	\medskip
	在本研究所提出的增长-$c(t)$ 框架内，太阳系形成于宇宙红移 $1+z\approx1.43$ 处。按标度律 $c(z)=c_0/(1+z)$，太阳形成时期的真空光速约为今日值的 70\%。相应地，氢、锂、碳以及所有其他元素的玻尔半径都缩小至今日尺寸的 70\%，从而对原太阳的初始锂储量施加了内在限制。
	
	\medskip
	锂是一种轻元素。较低的 $c(t)$ 与较小的玻尔半径共同使锂更容易启动核燃烧。锂原子在年轻的太阳中迅速烧尽，因此大量锂在早期便已损失。这一过程造就了远低于采用固定光速之计算结果的太阳表面锂水平。
	
	\medskip
	许多年轻恒星位于太阳十光年之内。这些恒星形成于 $z=0$ 附近的红移处，当时宇宙光速已接近今日值。我们的模型给出一个可检验的预言：邻近的年轻恒星损失的锂要少得多，因此其表面锂丰度将明显高于太阳。
	
	
	
	
	
	\medskip
	\bigskip
	\textbf{地球生物演化与 $\boldsymbol{G(t)}$}\,\,\,
	寒武纪大爆发与石炭纪的碳埋藏，代表地球生物演化史上的两个标志性事件。在标准的地质学与古生物学框架中，这些事件主要通过行星尺度的内部过程来解释——例如超大陆裂解、大陆风化速率、海洋氧化还原分层、降解木质素真菌的演化时序，以及沼泽沉积物的保存条件。
	
	\medskip
	按主流观点，在寒武纪之前的近三十亿年间，地球上只有简单的生命形式，演化进程大体上被多重环境阈值所压制。直到 2500 万年（540--515~Ma）这一狭窄窗口，关键的行星条件才同时越过各自的演化阈值，触发了寒武纪辐射——大多数主要动物门类在地质学上极为短暂的一段时期内涌现。
	
	\medskip
	石炭纪拥有一个独特稳定的行星气候，持续长达 6000 万年（360--300~Ma）。温暖而富氧的环境促成了大规模的地质碳封存；随后大气氧含量的下降淘汰了脑容量有限的巨型生物，从而为哺乳动物大脑的演化打开了生态位。
	
	\medskip
	于是地球的演化记录呈现出一种引人注目的不对称：数十亿年的长期停滞，被极短地质窗口内的骤然生物转变所打断。当前许多理论把智慧生命的出现归因于随机的巧合——即所谓的{\textit{幸运地球}}假说。
	
	\medskip
	在 $G(t)$ 与 $c(t)$ 持续演化的背景之下，我们尝试性地提出：一个完整的图景或许需要把三个尺度上的嵌套约束纳入其中：{\textit{宇宙、恒星与行星}}。尤其是，含时引力系数 $G(t)$ 可能在寒武纪与石炭纪经历了关键的弛豫过程。 
	
	\medskip
	随着引力耦合趋于稳定，太阳系内的轨道扰动被降到最低，地表太阳辐射通量的涨落被大幅抑制，地球的自转与岁差周期进入一个长期的动力学规则阶段。
	
	\medskip
	行星尺度的分析固然仍不可或缺，但把宇宙背景场 $G(t)$ 与 $c(t)$ 所设定的宏观边界一并纳入，或许能提供更宽广的语境。 
	
	\medskip
	于是人们或许更倾向于认为：人类智慧恰在此刻出现于地球，并非纯粹的偶然，而可能是由一连串必要的物理条件所开启的一种偶然结果。不过，这一解读仍高度思辨，我们在此谨慎地提出。
	
	
	
	\medskip
	\bigskip
	\textbf{LLR 并不能证明 $\boldsymbol{G}$ 严格恒定}\,\,\,
	最精确的 $G$ 直接实验室测量由罗俊团队（2018）报告~\cite{Luo:2018}，采用两套独立的卡文迪什型扭秤方案，实现了 $11.6\times 10^{-6}$ 的相对不确定度。 
	
	\medskip
	相比之下，由月球激光测距（LLR）给出的 $\Delta G/G \sim 10^{-13}\,\mathrm{yr}^{-1}$ 约束并非直接测量，而是建立在“光速 $c$ 恒定”这一先验模型假设之上的推断~\cite{Mueller:2014}。我们下面剖析这一推断背后的推理。
	
	\medskip
	在扣除太阳与其他行星已知的扰动效应之后，LLR 观测在光速恒定的假设下给出一项表观的月球退行速率 ${\Delta D} = \boldsymbol{3.8}\ \mathrm{cm\,yr^{-1}}$~\cite{Murphy:2013}；标准 LLR 模型把这一退行全部归因于地月系统内部的潮汐角动量耗散。这对应往返光行时的相对增长：
	\[
	\frac{\Delta T}{T}\thinspace =\thinspace \frac{\Delta D}{D} \thinspace =  \thinspace \frac{3.8\ \mathrm{cm/yr}}{3.844\times 10^{10}\ {\mathrm{cm}}} \thinspace = \thinspace 9.89\times 10^{-11}\ \mathrm{yr}^{-1}\,\,.   \eqno(12.1)
	\]
	
	\medskip
	\textit{\textbf{注：}关键在于，$\Delta D$ 并非直接测量所得；LLR 只测量往返光行时 $2T$ 及其年变化 $2\Delta T$。退行距离 $\Delta D$ 是在假设光速 $c$ 恒定的前提下，由 $\Delta T$ 推断出来的。}
	
	\bigskip
	如第 11 节所述，在我们的超精细模型中，$\Delta c/c = 6.89\times 10^{-11}\ \mathrm{yr}^{-1}$。由于 $T = D/c$（单程光行时），取对数微分即得真实的轨道扩张比：
	
	\[
	\frac{\Delta D_{\mathrm{real}}}{D} \thinspace \approx \thinspace \frac{\Delta {T}}{T} + \frac{\Delta c}{c}\,\,,   \eqno(12.2)
	\]
	\[
	\frac{\Delta D_{\mathrm{real}}}{D} 
	\thinspace \approx  \thinspace 9.89\times 10^{-11} + 6.89\times 10^{-11}
	\thinspace = \thinspace 16.78\times 10^{-11}\ \mathrm{yr}^{-1}\,\,,   \eqno(12.3)
	\]
	
	
	对应一个物理退行速率：
	\[
	\Delta D_{\mathrm{real}} \thinspace = \thinspace 3.8\times\left(1+\frac{6.89}{9.89}\right) \thinspace \approx \thinspace \boldsymbol{6.45}\ \mathrm{cm\,yr^{-1}}\,\,.   \eqno(12.4)
	\]
	
	\medskip
	因此，在我们的超精细增长-$c(t)$ 模型下，月球以 $\boldsymbol{6.45}\ \mathrm{cm\,yr^{-1}}$ 的速率退离地球——显著大于标准 LLR 分析所推断的 $3.8\ \mathrm{cm\,yr^{-1}}$。
	
	\medskip
	此外，基于本工作所采用的核心关系 $G c^{2}=\mathrm{constant}$，一阶微分给出引力常数的时间变化关系：
	\[
	\frac{\Delta G}{G} \thinspace = \thinspace -2\ \frac{\Delta c}{c} \thinspace = \thinspace -13.78\times 10^{-11}\ \mathrm{yr}^{-1}.   \eqno(12.5)
	\]
	
	\medskip	
	在地月两体系统角动量守恒的条件下，$G R=\mathrm{constant}$，由此得到
	\[
	\frac{\Delta D_\mathrm{gravity}}{D} \thinspace = \thinspace -\frac{\Delta G}{G} \thinspace = \thinspace 13.78\times 10^{-11}\ \mathrm{yr}^{-1}    \,\,.   \eqno(12.6)
	\]
	
	\medskip
	\textit{\textbf{注：}这一部分轨道扩张完全源自引力 $G$ 的宇宙演化，与潮汐耗散无关。}
	
	\medskip	
	从总速率中减去宇宙贡献，剩余部分即对应\textit{\textbf{潮汐}}耗散所诱发的轨道扩张。于是我们有：
	
	\[
	\frac{\Delta D_\mathrm{tide}}{D}
	\thinspace = \thinspace \frac{\Delta D_\mathrm{real}}{D} - \frac{\Delta D_\mathrm{gravity}}{D}
	\thinspace = \thinspace \boldsymbol{3.00}\times 10^{-11}\ \mathrm{yr}^{-1}    \,\,.   \eqno(12.7)  
	\]
	
	\[
	\Delta D_{\mathrm{tide}} \thinspace = \thinspace 6.45\times\left(\frac{3.00}{16.78}\right) \thinspace \approx \thinspace \boldsymbol{1.15}\ \mathrm{cm\,yr^{-1}}.   \eqno(12.8)
	\]
	
	\[
	\Delta D_{\mathrm{gravity}} \thinspace = \thinspace 6.45\times\left(\frac{13.78}{16.78}\right) \thinspace \approx \thinspace \boldsymbol{5.30}\ \mathrm{cm\,yr^{-1}}.   \eqno(12.9)
	\]
	
	\medskip
	因此，与把全部观测退行都归因于潮汐耗散的标准 LLR 模型不同，我们的变 \(c\) 模型给出一种截然不同的解读。在我们的模型中，\(\Delta D_{\mathrm{tide}}\) 仅为 \(1.15\ \mathrm{cm\,yr^{-1}}\)，显著小于标准 LLR 框架所推断的 \(3.8\ \mathrm{cm\,yr^{-1}}\)。
	
	\medskip
	因此我们得出结论：LLR 模型对 \(\Delta G/G\) 的 \(10^{-13}\) 约束，不能被视为对本工作所提猜想的坚实否定。
	
	
	
	\medskip
	\bigskip
	\textbf{卡西米尔效应并非真空能灾难的证据}\,\,\,
	在常规量子场论中，真空被假定含有可观的量子零点能，而卡西米尔力被解读为其存在性的决定性实验证据。在纳米级间距下精确测得的标度关系 $F \propto D^{-4}$，被视为这类真空涨落充满全部空无一物空间的直接证据。 
	
	\medskip
	若这一真空能确实无处不在，其密度将超出观测到的宇宙暗能量密度约 $10^{120}$ 倍。这就是所谓的{\textit{真空能灾难}}，长期以来被视为标准 $\Lambda$CDM 框架中一个未解的根本性危机。
	
	\medskip
	然而，这一灾难并非直接的观测结果。Jaffe（2005）~\cite{Jaffe:2005casimir} 证明，卡西米尔力可以在不引入任何全局零点真空能的情况下导出；该效应只测量两个边界之间的能量差，而非绝对的真空能。 

	\medskip
	本工作遵循并拓展了 Jaffe 的论证。我们在 LEM-SCDF 背景场框架内重新解读卡西米尔效应，并把这一推理进一步推进到宇宙学暗能量的考量之中。
	
	\medskip
	这一严重的理论困境，完全源自一次不合理且不合物理的外推。在主流理论中，两块材料板之间测得的局域边界效应，被错误地扩展为充满整个宇宙的均匀真空属性。在此提出的 LEM-SCDF 背景场图景中，卡西米尔力源自相邻材料边界对普适背景场造成的局域压缩。板间的压缩必然在周围空间触发互补的稀薄化，从而建立背景场密度在全局上的守恒再分布。所观测到的力只反映这一局域失衡。该力只在受限的实验室几何下显现，也没有任何经验证据确认这种极端真空能均匀充满宇宙的空无一物空间。
	
	\medskip
	一旦放弃这一人为的全局外推，真空能灾难便不复存在。卡西米尔效应不再能充当暗能量的有效证据，也再无必要引入假想的暗能量组分来挽救宇宙观测。

	
	\medskip
	\bigskip
	\textbf{把局域更新误读为宇宙创生}\,\,\,
	标准宇宙学借用我们今天所见之物，去解释太初所发生之事。超新星与黑洞是成熟宇宙中的局域事件，研究者却把它们向时间上游外推，用以描述早期宇宙。 
	
	\medskip
	这种全球性早期宇宙相变与局域性晚期宇宙事件之间的逻辑错配，此前已有学者注意：Milne 的宇宙学原理~\cite{Milne:1933} 强调早期宇宙的全局均匀性，而 Barrow 与 Tipler~\cite{Barrow:1986} 则强调早期时期的极端物理条件与晚期局域环境之间的根本区别。
	
	\medskip
	这构成一种根本性的逻辑错配。所有观测到的超新星与黑洞都来自晚期宇宙，即原子、星系与恒星形成之后。相比之下，早期宇宙经历的是强子化、等离子体耦合与 CMB 退耦阶段，那时没有恒星，没有致密天体，也没有引力塌缩。
	
	\medskip
	早期宇宙的演化是全球性的：整个宇宙作为一个统一系统经历相变。先是强子物质，然后是等离子体，再是中性原子与 CMB——这一序列构筑了宇宙的基本物质基础。这才是全部重子物质的真正创生。
	
	\medskip
	黑洞与超新星则纯粹是局域的。它们只在宇宙背景完全就位之后才发生。它们重新分配质量、释放能量、循环利用恒星物质。它们是星系内部的更新过程——而不是宇宙本身的创生过程。
	
	\medskip
	我们应当在这两个不同的物理层次之间划出清晰的界线。这一区分将消除一个长期存在的混淆，并为理解早期宇宙结构形成提供一个物理上自洽的基础。
	
	
	
	
	\medskip
	\bigskip
	\textbf{关于重建弗里德曼方程}\,\,\,
	标准弗里德曼方程（式 (1.14) 与式 (1.15)）由爱因斯坦场方程导出。既然修正后的 $G_{\mu\nu}$ 现由式 (8.2) 给出，我们就需要重建弗里德曼方程，并对 \texttt{Planck18} 红移数据施加必要的修正。
	
	
	\medskip
	\bigskip
	\textbf{为何刻意回避哈勃张力}\,\,\,
	本文并不试图完整重建现代宇宙学，也不继承 $\Lambda$CDM 的全套假设。此外，解决哈勃张力并未被设为首要目标。本研究的核心目标，是为一个时变光速建立内部自洽的标度体系与猜想，并给出可经经验检验、且有可能被否证的定量预言。
	
	\medskip
	众所周知的哈勃张力，源自哈勃常数的两种不同推导方法。然而，$\Lambda$CDM 宇宙学本身嵌套着一长串相互依赖的先验假设：FLRW 度规、标度因子 $a(t)$ 的物理解释、把宇宙学红移归因于空间膨胀、均匀且各向同性的宇宙（宇宙学原理）、各类宇宙组分（物质、辐射、暗能量）的状态方程，以及引力常数与真空光速的恒定性。本文并不以解释哈勃张力为目标，也不试图在标准 FLRW 框架内拟合 $H_0$ 的两个测量值。这并非模型的缺陷，而是对研究边界的刻意划定。
	
	\medskip
	在本文的物理概念框架中，光速随宇宙时间上升，而 $G(t)$ 随 $c(t)$ 增大而减小。引力束缚系统（例如地月系统）的距离将不仅因潮汐效应而改变，还会因引力常数的逐渐减小而改变。换言之，束缚系统同样受 $c(t)$ 与 $G(t)$ 演化的影响——这与“束缚系统不膨胀”这一标准宇宙学前提形成尖锐对照。
	
	\medskip
	关于对标的度因子 $a(t)$ 与宇宙学红移的基本立场：传统观点认为，宇宙学红移等价于空间膨胀，并以标度因子 $a(t)$ 描述，而星系内部的引力束缚系统并不参与这一膨胀；只有星系际空洞的时空在膨胀。作者的立场是：承认宇宙的整体膨胀效应，但拒绝把宇宙学红移完全归因于空间尺度的拉伸。我认为，不能预先把光速的演化与标度因子 $a(t)$ 之间的关系设定为简单的线性关系。宇宙标度因子并非可直接测量的实体，而只是标准宇宙学内部的一个抽象构造；$a(t)$ 的物理含义完全取决于那一整套理论假设。
	
	\medskip
	在增长-$c(t)$ 模型框架内，由于真空光速随宇宙时间演化、引力常数同步改变，上述整套先验假设便不再天然成立。若强行在弗里德曼方程框架内拟合哈勃张力，必将出现无法跨越的量级冲突。因此，本文刻意回避对哈勃张力的讨论。我不否认宇宙的膨胀效应，但我不预设宇宙学红移唯一地源自空间尺度拉伸，也不预设光速演化与传统标度因子之间存在简单的线性关系。本文并不试图为全部宇宙学观测提供统一解释，也不承担重建完整宇宙动力学方程组的责任。
	
	\medskip
	本文所提猜想是否成立，最终将由未来为期三十年的 \textit{Mission-M} 光速测量计划来判定。若后续经验检验证实真空光速随时间缓慢增长，那么重新推导宇宙动力学方程组、重新解读宇宙学观测，将成为未来其他物理学家的研究课题。
		
	
	\medskip
	\bigskip
	\textbf{关于现代宇宙学的方法论反思}\,\,\,
	有评论者指出，现代标准宇宙学框架有退化为所谓“\textit{Sapeur 物理学}”的风险：繁复的数学形式充当华丽的外衣，而基础的物理前提却建立在未经证实的层层叠加假设之上。当 $\Lambda$CDM 范式内部浮现观测张力时，惯常采用的对策是引入额外的自由参数与特设组分，以调和观测上的不匹配。这一反复修补的过程被比作\textit{过拟合}——无数修正补丁被堆叠到一个本身就有缺陷的底层理论结构之上。
		
	\medskip
	我们应当回到海森堡所强调、并由杨振宁重申的基本原则：物理理论应当围绕可观测展开。中间性的理论构造不必可直接测量。例如，广义相对论中的时空曲率无法用任何单一独立仪器探测，但其导出的结果——光线偏折、引力红移与引力波——都是可定量核验的可观测量。这条推理路线保持着坚实的经验基础。
		
	\medskip
	这一严格标准的一个代表性基准，是量子力学的哥本哈根诠释。尽管其直觉性的物理图景在许多人看来违反常识、难以想象，该框架却在所有量子实验情形下给出了完全可操作的数学计算与精确、可检验的物理预言。其抽象的内在逻辑从未脱离可测量的光谱、自旋与干涉观测量——这正是它区别于当代宇宙学中纯思辨的多层假设之处。
		
	\medskip
	基于这一标准，必须划出一条关键界线。人们可以容忍一层不可观测的中间实体，用以在理论与测量之间搭桥。但若层层嵌套不可检验的猜想，且每一新层都缺乏独立的观测锚点，则在哲学上是站不住脚的。这样的做法抛弃了物理学的经验内核。一个理论模型不能仅仅为了迁就既有数据，就无限期地依赖未被观测到的假想成分。
		
	\medskip
	本工作也严格以这一经验纪律自我约束。尽管光速 $c(t)$ 极其细微的时变行为——其变化幅度约在 $10^{-13}$ 至 $10^{-14}$ 量级——在未来十年内还无法直接测量，所提出的框架仍对多个宇宙可观测量给出了具体、可否证的预言，包括宇宙红移的演化、随 $c(t)$ 一同演化的黑洞形成概率的时期依赖，以及依赖红移的暗能量状态方程 $w(z)$。本研究的全部理论推断，最终都可追溯到天文上可及的观测物理量。
	
	
	
	
	
	
	
	
	\section*{
		\begin{center}
			后记
		\end{center}
	}
	\smallskip
	\noindent
	
	\textbf{一首献给呼吸宇宙与智慧生命的赞歌}\,\,\,
	在最初的日子里，宇宙与我们这颗蓝色地球都是笨拙的新手。巨大的古老恒星炽烈燃烧又骤然崩塌，寿命太短，来不及孕育行星系统。我们的星球曾经连赤道都被冻结，也曾承受五次毁灭性的物种大灭绝。生命在无数次试炼中升起、倾覆、重启，一次又一次。每一次崩塌都不是终点，而是新的开始。每一次挫折，都在为生命最终的绽放积蓄无声的力量。
	
	\medskip
	经过数十亿年的演化，宇宙的各项基本物理量落入了完美的平衡。光速、引力、地球 23.5 度的自转轴倾角、大气中的氧与臭氧层，都稳定在适宜维持生命的理想状态。这才是呼吸宇宙的真正含义：它依循时间无尽的节律运行。恒星、行星与一切生灵都与宇宙一同呼吸，而智慧生命正是诞生于这一永恒的节律。
	
	\medskip
	愿智慧之光永不熄灭。地球已等了四十亿年，宇宙等得更久，只为等有意识的生命出现。愿人类文明不负这漫长的等待，把一瞬的火花化作横跨星系的恒久光辉。我们脆弱，却无比珍贵。我们并非宇宙中的匆匆过客。相反，我们是宇宙在数十亿年的沉默之后，终于用以看见并理解自己的\textit{那双眼睛}。
	
	\bigskip
	\textbf{关于本工作的缘起}\,\,\,
	我最初对光与引力的本质产生疑问，已是几十年前的事。一个突如其来的洞见在三月初的一次驾车途中降临。意识到这是物理学最深奥的谜题之一，我甚至不敢让自己去想它。到了五月初，日益增长的冲动催促我把这项研究写下来。在这项工作中，我尝试构建一个逻辑自洽的物理框架。即便我的理论最终与未来的观测证据不符，我也希望它能充当一个富有想象力的思想实验。
	


	\vskip 0.2in
	\section*{
		\begin{center}
			附录：\\
			倘若该猜想成立
		\end{center}
	}
	\medskip
	\noindent
	
	\textit{免责声明：能够回应许多长期悬而未决的谜题，只不过是一张入场券。它并不构成该框架在物理上正确的证明。本附录仅呈示概念性的思想实验推理。}
		
	\bigskip
	若其核心公设成立，这一框架将从根本上有力挑战 $\Lambda$CDM 标准宇宙学范式。其影响将遍及广义相对论、粒子物理、天体物理、计量学与科学方法论。它或许能化解二三十个长期存在的观测与理论冲突，汇总如下。
	
	\medskip
	这一框架仅建立在两条简洁的核心公设之上：
	\begin{itemize}
		\item 宏观公设：存在一个不变的宇宙耦合常数 $\boldsymbol{P}$。关系 $G(t)c(t)^2 = \boldsymbol{P}$ 在全部宇宙历史中成立。
		
		\item 微观公设：精细结构常数 $\boldsymbol{\alpha}$ 在所有宇宙时期都保持不变。
	\end{itemize}
	
	\medskip
	标准 $\Lambda$CDM 为迁就观测张力，动用了暗能量、暗物质与暴胀，并附带层层辅助假设。增长-$c(t)$ 模型并不引入多个额外自由参数，而是用单一演化变量 $c(t)$ 来回应众多开放的物理谜题。这是它在概念上的主要优势。
		
			
	\medskip
	\bigskip
	\textbf{I.\ 宇宙学：正文所讨论的十个核心谜题}
	\begin{itemize}
		\item \textbf{01\_} \textit{宇宙学红移背后的因果机制。}\,\,\,标准宇宙学给出的是几何图景：空间膨胀并拉长光子波长。然而它无法回答那个因果问题：当光子脱离发射者、穿过近乎空无一物的空间之后，究竟是什么物理作用者持续地把宇宙膨胀告知光子？在本模型中，宇宙学红移是真空光速的宇宙演化所产生的观测印记。全局同步的 LEM-SCDF 背景场，提供了这一缺失的因果环节。
		
		\item \textbf{02\_} \textit{拒绝热大爆炸初始条件。}\,\,\,标准模型假设一个极端温度与密度的初始状态。本猜想采用 ZLZT（零光速与零温度）原初态。宇宙从近乎零的光速逐步演化，而非始于一场剧烈的爆炸。CMB 退耦之前的宇宙时标，变得远长于 $\Lambda$CDM 所预言的 370 kyr。
		
		\item \textbf{03\_} \textit{CMB 退耦处的电离温度悖论。}\,\,\,标准宇宙学预言 CMB 时期附近存在极高的温度。相比之下，本框架从 ZLZT 原初假设出发，得出 CMB 时期的光速约为 272 km/s。复合时对应的电离温度仅约 82.64 K，从根本上修正了宇宙复合与再电离的物理图景。
				
		\item \textbf{04\_} \textit{史瓦西半径的演化与高红移星系中行星系统的稀少。}\,\,\,由于 $r_s(t)\propto c(t)^{-4}$，早期宇宙中的史瓦西半径变得极为巨大，引力俘获截面被大幅放大。早期恒星迅速吸积周围物质并抹去原行星盘，行星系统几乎得不到形成的机会。只有当 $c(t)$ 上升、$r_s(t)$ 收缩之后，宇宙环境才趋于稳定，容许行星广泛形成。这与 JWST 的观测自然吻合：高红移星系以恒星为主，行星候选体极少。启发式公式式 (7.1) 进一步给出依赖红移的行星形成概率预言，可供未来观测检验。
		
		\item \textbf{05\_} \textit{星族 III 恒星的时标危机。}\,\,\,JWST 在极高红移处观测到完全组装好的大质量星系。在标准 $\Lambda$CDM 中，无金属的星族 III 恒星主序寿命极短，而光速恒定的时间线无法为星系组装留出足够时间。然而在本模型中，早期 $G(t)$ 在 $GM$ 简并之下取较大值：由此推断的 $M$ 小于标准估计，于是核燃烧进行得更慢（尽管仍相对较快），有效物理时标被拉长。标准年表所计的 $\sim 1$ Myr，在 $1+z \sim 10$ 附近或可对应 $\sim 10$ Myr 或更长。原初恒星的真实重子质量下降，而核燃烧时标延长。这无需任何奇异恒星核心加热，或极端的初始质量函数假设。
				
		\item \textbf{06\_} \textit{DESI DR2 动态暗能量张力（精质—幽灵穿越行为）。}\,\,\,标准宇宙学难以解释观测到的 $w(z)$ 穿越 $w=-1$ 的现象。本模型并不引入真实的暗能量流体。随着 $c(t)$ 增大，$G(t)$ 单调下降，从而修正宇宙膨胀历史。当用标准的 $\Lambda$CDM 参数管线来分析时，这种演化便重现出表观上的精质—幽灵穿越信号。表观的动态暗能量，或许最终只是一个观测伪迹。
		
		\item \textbf{07\_} \textit{再电离期间的宇宙光子预算危机。}\,\,\,JWST 在 $z=25$ 处探测到丰富的远紫外光子，而 21 cm 测量显示同一时期有相当数量的中性氢共存。标准模型需要精细调节的反馈与成团因子才能调和。本框架严格区分辐射黑体温度 $T_\mathrm{rad}(z)$ 与气体热力学温度 $T_\mathrm{gas}(z)$。高红移处星系际气体保持极冷，从而给出极高的复合率。这便在无需人为拟合参数的情况下缓解了张力。
		
		\item \textbf{08\_} \textit{太阳锂丰度之谜。}\,\,\,太阳系形成于 $1+z\approx1.43$ 处，当时 $c(t)$ 约为今日值的 70\%。玻尔半径相应缩小。锂经历了高效的早期核燃烧。这解释了所观测到的太阳表面锂亏缺。由此可得一项可否证的预言：邻近的同龄年轻恒星应表现出明显更高的表面锂丰度。
		
		\item \textbf{09\_} \textit{月球激光测距（LLR）的约束不再能限定 $G$ 恒定。}\,\,\,众所周知的 LLR 对 $\Delta G/G$ 的上限，是在真空光速恒定这一公理下导出的。若 $c(t)$ 本身存在长期漂移，往返光行时信号便混合了两种贡献：真实的月球轨道退行，以及 $c(t)$ 的宇宙变化。标准 LLR 的限制不能充当引力常数不随时间变化的严格证据。
		
		\item \textbf{10\_} \textit{重新审视狄拉克大数假说。}\,\,\,狄拉克在 1937--1938 年猜想基本常数或许随时间演化，但未给出自洽的框架。在本框架中，$G(t)\propto c(t)^{-2}$ 且 $N_\mathrm{Dirac,1}\propto c(t)^5$。数值上的接近只是一项暂态特征，并非普适定律。 
	\end{itemize}
	
	
	\bigskip
	\textit{\textbf{注：}以下所有条目均为概念性推断，而非已被证明的结论。受篇幅所限，它们未在正文中充分展开。}
	
	
	\medskip
	\bigskip
	\textbf{II.\ 延伸的天体物理谜题}
	\begin{itemize}
		\item \textbf{11\_} \textit{高红移星系的系统性质量估计偏差。}\,\,\,天文可观测量只约束乘积 $GM$，无法把 $G$ 与重子质量 $M$ 分开。目前所有高红移星系与黑洞的质量，都采用太阳系局域值 $G_0$ 来评估。若早期时期的 $G(t)$ 更大，则宇宙天体的内禀重子质量将远低于标准估计。若干与“超大质量黑洞为何在极高红移处就已存在”相关的张力，或可得到部分缓解。
		
		\item \textbf{12\_} \textit{类星体与活动星系核性质的红移依赖。}\,\,\,早期扩大的史瓦西半径增强了引力俘获截面与气体吸积效率。这为理解高红移类星体为何异常明亮提供了新的视角，且无需奇异的黑洞种子形成图景。
		
		\item \textbf{13\_} \textit{星系旋转曲线与暗物质问题的一部分。}\,\,\,早期宇宙中增强的引力耦合 $G(t)$，意味着普通重子物质也能产生更强的引力效应。星系动力学中一部分“缺失质量”信号，或源自随时间演化的 $G(t)$。尽管如此，正文明确指出该机制无法完全解释低红移星系团的质量失配，因此暗物质不能被彻底排除。
		
		\item \textbf{14\_} \textit{恒星演化图的全面修订。}\,\,\,恒星结构方程直接依赖 $G$。而在不同宇宙红移处，$G(t)$ 并不相同。零龄主序轨迹、质量—光度关系，以及变星的脉动行为，都将变得依赖红移。造父变星是宇宙距离阶梯的一根支柱。若本猜想成立，整个距离尺度都必须重新标定，这为化解哈勃张力提供了新路径。
		
		\item \textbf{15\_} \textit{跨宇宙时间的伽马射线暴与超新星性质。}\,\,\,Ia 型超新星被当作标准烛光。若 $G(t)$ 以及随 $c(t)$ 标度的原子尺度量随红移演化，则超新星峰值光度并非普适常数。基于超新星的宇宙加速分析所得出的结论，需要重新评估。
		
		\item \textbf{16\_} \textit{重新审视大爆炸核合成（BBN）。}\,\,\,反应截面与库仑势垒依赖 $\varepsilon_0(t)$ 与 $G(t)$。标准 BBN 计算假设 $c$ 与 $G$ 固定。在本模型中，原初核合成的温度—密度条件与宇宙时序完全改变。这为观测到的原初氦与氘丰度打开了新的解释空间。
		
		\item \textbf{17\_} \textit{卡西米尔效应与真空能灾难。}\,\,\,常规量子场论把卡西米尔力当作全局零点真空能存在的证据，从而引出所谓的真空能灾难。本工作在 LEM-SCDF 框架内拓展了 Jaffe 的论证：卡西米尔力源自普适背景场的局域压缩，而这一压缩必然伴随别处的互补稀薄化，使总场密度在全局上守恒。一旦放弃这一人为的全局外推，真空能灾难便不复存在，卡西米尔效应也不再能充当暗能量的有效证据。
	\end{itemize}
	
		
	\medskip
	\bigskip
	\textbf{III.\ 基础理论物理}
	\begin{itemize}
		\item \textbf{18\_} \textit{爱因斯坦场方程的重新解读与修正。}\,\,\,原场方程中的常数 $G$ 被 $P/c(t)^2$ 取代。广义相对论不再是纯粹的几何引力理论。时空曲率动力学与真空电磁背景 LEM-SCDF 场相耦合。质能分布、时空几何与电磁真空参数构成一个闭合的动力学环路。
		
		\item \textbf{19\_} \textit{通往统一场论（UFT）的另一条路径，绕开常规引力量子化的死胡同。}\,\,\,传统研究纲领试图直接量子化时空几何，却遭遇严重障碍。本框架把动力学自由度转移到 LEM-SCDF 背景场。人们可以优先量子化这一宇宙背景场，而非量子化引力本身。这为爱因斯坦 1920 年莱顿演讲中关于物理性时空新以太的构想，提供了一种现代实现。
		
		\item \textbf{20\_} \textit{量子纠缠“鬼魅般超距作用”的物理本体。}\,\,\,LEM-SCDF 场的传播远快于局域光速。关键在于，它不输运能量或因果信号，只实现全局背景态的同步。它不违反狭义相对论的局域光锥因果，却为非局域量子关联提供了物理基底。纠缠不再仅仅是一个数学解释上的特征。
		
		\item \textbf{21\_} \textit{相对历史上变光速（VSL）模型的重大升级。}\,\,\,Magueijo 早先的 VSL 方案旨在化解视界与平坦性疑难，却加剧了光子散射的病态。相比之下，ZLZT 原初条件以其近乎零的初始光速，从一开始便绕开了那些散射问题。
		
		\item \textbf{22\_} \textit{SI 单位体系的修正基础逻辑。}\,\,\,现代\texttt{米}的定义锚定于固定的常数 $c$。在本猜想之下，必须区分现今计量学意义上的米 $L_m$ 与原子米 $L_a$——后者由固定的原子振荡周期定义，而其物理空间延展随 $c(t)$ 变化。许多计量学测量都需要重新解读。 
		
		\item \textbf{23\_} \textit{一切普朗克尺度量都成为导出的演化量。}\,\,\,普朗克质量、普朗克长度、普朗克密度与普朗克力不再是永恒的普适常数。它们服从与 $c(t)$ 的幂律标度。这一尺度显现为演化中的宇宙背景的一种属性，而非量子引力物理的固定基准。这将重塑量子引力研究的起始假设。
		
		\item \textbf{24\_} \textit{囚禁生成物质：物质与能量概念的清晰分离。}\,\,\,标准 QCD 把强子质量解读为储存的囚禁能量。本工作引入 $M_\mathrm{aps}$，即元粒子物质。物质本身守恒；$E = M c(t)^2$ 充当物质与能量之间一个随时间变化的宇宙汇率。物质与能量成为概念上彼此独立的范畴。这修正了对质能等价原理的物理与哲学理解。
		
		\item \textbf{25\_} \textit{宇宙作为一个全局相干的实体。}\,\,\,传统物理学把宇宙视为一堆彼此断开的局域沙粒的集合，而常数 $c$ 与 $G$ 则由外部强加。在这一图景中，所有质能组分共同塑造 LEM-SCDF 背景场。每个粒子与光子都保持着持续的全局耦合。宇宙成为一个动力学上相干的整体，而非一片孤立碎块的拼集。
	\end{itemize}
	
		
	\medskip
	\bigskip
	\textbf{IV.\ 行星科学与生物演化}
	\begin{itemize}
		\item \textbf{26\_} \textit{系外行星族群的红移依赖统计预言。}\,\,\,由 $r_s(z)$ 的标度关系，我们得到可检验的统计预期：在高红移处，行星出现率应被显著压制。随着宇宙时间推进，行星系统逐渐变得普遍。未来的 JWST 与 PLATO 系外行星巡天，可以对这一趋势作统计上的证实或否证。
		
		\item \textbf{27\_} \textit{原初宇宙磁场的演化。}\,\,\,由于 $\mu_0(t)$ 与 $\varepsilon_0(t)$ 随 $c(t)$ 一同演化，宇宙磁场生成的全部历史都被改写。这为求解星系磁场起源提供了新的路径。
		
		\item \textbf{28\_} \textit{宇宙背景对地球生物演化的约束。}\,\,\,随着 $G(t)$ 缓慢衰减，太阳系的轨道扰动与太阳辐射通量的变化幅度都在减小。寒武纪大爆发与石炭纪事件这类重大演化转折，并非纯粹的地质偶然。它们或许对应于宇宙尺度引力背景条件弛豫到适宜范围的时间窗口。智慧生命的出现不是纯粹的随机运气，而是宇宙物理逐步弛豫所开启的一种偶然结果。
	\end{itemize}
	
	
	\bigskip
	\textbf{概念上的吸引力并不等于物理上的真理}\,\,\,
	回应既有的观测张力只是一个必要条件，而非充分证明。解释力只是入场资格，不是判决。决定性的分量落在可否证的预言上：微小的每日相对变化，需要 HEPS 或拟议中的 Mission-M 三十年计划这类未来装置；系外行星出现率随红移的统计关系；高红移恒星与造父变星的观测性质；以及更高精度的 DESI 与未来 BAO 对 $w(z)$ 的测量。
	
	\medskip
	尽管该框架把自由参数的数量保持在很低水平，它仍引入了两条全新的本体论公设：LEM-SCDF 宇宙暗场与 $M_\mathrm{aps}$ 元粒子物质。目前仍缺乏针对这两类实体的直接独立实验证据。大量后续理论工作也尚待展开：弗里德曼方程、原初核合成、完整的恒星结构计算，以及宇宙距离标定，都需要作完整的自洽重推。本文应被视为一份纲领性的纲要；许多推导仍属启发性的。
		
	\bigskip
	\textbf{倘若该猜想成立}\,\,\,	
	若这一猜想得到确认，$\Lambda$CDM 标准模型将降格为仅适用于宇宙晚期的有效近似理论。剧烈的高温大爆炸将被光速缓慢增长的 ZLZT 图景所取代。暗能量以及部分暗物质信号，或许最终被证明是源自背景常数演化的观测伪迹。人们熟悉的引力常数 $G$ 将降格为一个演化的导出量；真正不变的宇宙量将是 $P$ 与精细结构常数 $\alpha$。狭义相对论的局域洛伦兹不变性将得以保留，但在宇宙尺度上，真空光速将表现为一个含时的全局背景。修正后的爱因斯坦场方程将通过 LEM-SCDF 场把引力与电磁学联系起来，为统一场论研究开辟新方向。宇宙距离阶梯、恒星质量、黑洞质量、超新星标准烛光分析，以及原初核合成计算，都将需要重新标定。JWST 与 DESI 数以太字节计的存档数据，将需在新范式下重新分析。甚至 SI 单位的基础定义，以及物质与能量之间的概念区分，都会受到影响。
	
	\bigskip
	简而言之，本文并非对既有宇宙学的一次小幅补丁。它用 $G(t)c(t)^2=P$ 这单一关系，去回应横跨宇宙学、天体物理与基础理论的二三十个谜题。然而在物理学中，优雅的逻辑自洽只不过是一个起点。最终的判决，属于望远镜与未来的高精度实验。

	
	\vskip 0.5in
	\paragraph{资助：}本工作未获得任何资助。
	
	
	\vskip 0.5in
	%% \bibliography{refs}
	
	\begin{thebibliography}{10}
		
		\bibitem{NBS:1987}
		{\em The International System of Units (SI)}.
		\newblock NBS Special Publication 330. 1987.
		\newblock Appendices I and II contain the meter definition and practical
		realization.
		
		\bibitem{Astropy:2022}
		Astropy Collaboration, A.~M. Price-Whelan, et~al.
		\newblock {The Astropy Project: Sustaining and Growing a Community-oriented
			Open-source Project and the Latest Major Release (v5.0)}.
		\newblock {\em The Astrophysical Journal}, 935(2):167, 2022.
		
		\bibitem{Planck2018:cosmo}
		Planck Collaboration, N.~Aghanim, et~al.
		\newblock {Planck 2018 results. VI. Cosmological parameters}.
		\newblock {\em Astronomy \& Astrophysics}, 641:A6, 2020.
		\newblock The cosmological parameters used in the astropy.cosmology Planck18
		implementation.
		
		\bibitem{Tipler:2008}
		Paul~A. Tipler and Gene Mosca.
		\newblock {\em {Physics for Scientists and Engineers}}.
		\newblock W. H. Freeman and Company, New York, 6th edition, 2008.
		
		\bibitem{Weinberg:2008}
		Steven Weinberg.
		\newblock {\em {Cosmology}}.
		\newblock Oxford University Press, Oxford, 2008.
		
		\bibitem{Dirac:1937}
		P.~A.~M. Dirac.
		\newblock {The Cosmological Constants}.
		\newblock {\em Nature}, 139:323, 1937.
		
		\bibitem{BIPM:2019}
		Bureau~International des Poids~et Mesures.
		\newblock {The International System of Units (SI)}, 2019.
		
		\bibitem{Dirac:1938}
		P.~A.~M. Dirac.
		\newblock {A New Basis for Cosmology}.
		\newblock {\em Proceedings of the Royal Society of London. Series A,
			Mathematical and Physical Sciences}, 165(921):199--208, 1938.
		
		\bibitem{Einstein:632320}
		A~Einstein.
		\newblock {Die Feldgleichungen der Gravitation, 25 Nov}, 1915.
		
		\bibitem{Einstein:1922ether}
		Albert Einstein.
		\newblock {\em Sidelights on Relativity}.
		\newblock Methuen \& Co. Ltd., London, 1922.
		\newblock A lecture delivered on May 5, 1920, at the University of Leyden.
		
		\bibitem{LeidenLecture:2025}
		Albert Einstein.
		\newblock {\em {Is Space Empty? Einstein Resurrects the Aether}}.
		\newblock Independently Published, 2025.
		\newblock A new translation of the 1920 Leiden Lecture 'Ether and the Theory of
		Relativity'.
		
		\bibitem{Coulton:2021}
		William~R. Coulton, Benjamin Beringue, and P.~Daniel Meerburg.
		\newblock {Primordial information content of Rayleigh anisotropies}.
		\newblock {\em Physical Review D}, 103(4):043501, 2021.
		
		\bibitem{Magueijo:2003}
		João Magueijo.
		\newblock {New varying speed of light theories}.
		\newblock {\em Reports on Progress in Physics}, 66(11):2025--2068, 2003.
		
		\bibitem{Peskin:1995}
		Michael~E. Peskin and Daniel~V. Schroeder.
		\newblock {\em {An Introduction to Quantum Field Theory}}.
		\newblock Westview Press, Boulder, 1995.
		
		\bibitem{Feynman:2011}
		Richard~P. Feynman, Robert~B. Leighton, and Matthew Sands.
		\newblock {\em {The Feynman Lectures on Physics}}.
		\newblock Basic Books, New York, 2011.
		\newblock Vol. I, Chapter 42-5.
		
		\bibitem{Davies:1995}
		Paul C.~W. Davies.
		\newblock {\em {About Time: Einstein's Unfinished Revolution}}.
		\newblock Simon \& Schuster, New York, 1995.
		
		\bibitem{HEPS:2022}
		Y.~Jiao et~al.
		\newblock {The High Energy Photon Source (HEPS): A new 6-GeV synchrotron
			radiation facility in Beijing}.
		\newblock {\em Journal of Physics: Conference Series}, 2380:012012, 2022.
		
		\bibitem{Riess:1998}
		Adam~G. Riess et~al.
		\newblock {Observational Evidence from Supernovae for an Accelerating Universe
			and a Cosmological Constant}.
		\newblock {\em The Astronomical Journal}, 116:1009--1038, 1998.
		
		\bibitem{Gu:2025}
		G.~Gu, X.~Wang, Y.~Wang, et~al.
		\newblock {Dynamical dark energy in light of the DESI DR2 baryonic acoustic
			oscillations measurements}.
		\newblock {\em Nature Astronomy}, 9:1879--1889, 2025.
		
		\bibitem{AbreuTurner:2026}
		Matilde~L. Abreu and Michael~S. Turner.
		\newblock {DESI Dark Secrets}.
		\newblock {\em arXiv e-prints}, 2026.
		\newblock Submitted to a journal.
		
		\bibitem{Zwicky:1933}
		F.~Zwicky.
		\newblock {Die Rotverschiebung von extragalaktischen Nebeln}.
		\newblock {\em Helvetica Physica Acta}, 6:110--127, 1933.
		
		\bibitem{Madau:2023}
		Piero Madau.
		\newblock {The Cosmic Reionization History and the Photon Budget Crisis}.
		\newblock {\em Research Notes of the AAS}, 7(5):98, 2023.
		
		\bibitem{Austin:2025}
		D.~Austin, T.~Harvey, C.~Conselice, et~al.
		\newblock {Resolving the ionizing photon budget crisis with JWST/NIRCam HII
			clumping constraints at z=6}.
		\newblock {\em arXiv e-prints}, 2025.
		
		\bibitem{Reionization:2025}
		No~authors specified.
		\newblock {Towards Reconciling Reionization with JWST: The Role of Bright
			Galaxies and Strong Feedback}.
		\newblock {\em arXiv e-prints}, 2025.
		\newblock Submitted to ApJ.
		
		\bibitem{Labbe:2023}
		I.~Labbé, P.~van Dokkum, E.~Nelson, et~al.
		\newblock {A population of red candidate massive galaxies ~600 Myr after the
			Big Bang}.
		\newblock {\em Nature}, 616:266--269, 2023.
		
		\bibitem{Topping:2023}
		M.~W. Topping, D.~P. Stark, R.~Endsley, et~al.
		\newblock {The Evolving Intergalactic Medium and the Implications for the First
			Galaxies}.
		\newblock {\em The Astrophysical Journal}, 954(2):139, 2023.
		
		\bibitem{Prather:1976}
		Michael~J. Prather.
		\newblock {\em {Evolution of Low Mass Stars and the Effects of a Variable
				Gravitational Constant}}.
		\newblock PhD thesis, Yale University, 1976.
		
		\bibitem{Korn:2019}
		A.~J. Korn and J.~Meléndez.
		\newblock {Lithium in the Universe}.
		\newblock {\em Astronomy \& Astrophysics Review}, 27:4, 2019.
		
		\bibitem{Luo:2018}
		Qing Li, Chao Xue, Jian-Ping Liu, Jun-Fei Wu, Shan-Qing Yang, Cheng-Gang Shao,
		Li-Di Quan, Wen-Hai Tan, Liang-Cheng Tu, Qi~Liu, Hao Xu, Lin-Xia Liu,
		Qing-Lan Wang, Zhong-Kun Hu, Ze-Bing Zhou, Peng-Shun Luo, Shu-Chao Wu, Vadim
		Milyukov, and Jun Luo.
		\newblock Measurements of the gravitational constant using two independent
		methods.
		\newblock {\em Nature}, 560:582--588, 2018.
		
		\bibitem{Mueller:2014}
		Jürgen Müller, L.~Biskupek, F.~Hofmann, and E.~Mai.
		\newblock Testing relativity by lunar laser ranging.
		\newblock In {\em Frontiers in Relativistic Celestial Mechanics}, volume~2. de
		Gruyter, 2014.
		
		\bibitem{Murphy:2013}
		T.~W. Murphy.
		\newblock Lunar laser ranging: the millimeter challenge.
		\newblock {\em Reports on Progress in Physics}, 76(7):076901, 2013.
		
		\bibitem{Jaffe:2005casimir}
		R.~L. Jaffe.
		\newblock Casimir effect and the quantum vacuum.
		\newblock {\em Physical Review D}, 72(2):021301, 2005.
		
		\bibitem{Milne:1933}
		E.~A. Milne.
		\newblock {A Note on the Cosmological Principle}.
		\newblock {\em Proceedings of the Royal Society of London. Series A},
		142:40--56, 1933.
		
		\bibitem{Barrow:1986}
		John~D. Barrow and Frank~J. Tipler.
		\newblock {\em {The Anthropic Cosmological Principle}}.
		\newblock Oxford University Press, Oxford, 1986.
		
	\end{thebibliography}
	\newpage
	% 以下为论文末尾的版权页（作者通讯信息 + 版本历史 + 引用模板 + 文档指纹），
	% 属固定版式的一整块内容，本版仍按原行距 1.0 排，以免整块被拆散到两页。
	% 若希望这一页也随正文放宽 20%，删掉下一行即可。
	\linespread{1}\selectfont
		
	\vskip 0.6in
	\noindent
	XianLiang Ma\\[10pt]
	首席量化架构师\\[2pt]
	Beijing DeepQuant.AI Technologies Co., Ltd\\[2pt]
	中国 北京市朝阳区 100029\\[10pt]
	电子邮箱：maxianliang@gmail.com\\[10pt]
	2026 年 6 月 13 日
	
	\vskip 0.1in

	
	\begin{footnotesize}
		% 版权信息块（版本历史 / 引用模板 / 文档指纹）属固定版式，保持原行距 1.0，
		% 否则本块会被拆到两页、“存档于”与其下的三条链接分离。
		% 若希望此块也随正文放宽，删掉下一行即可。
		\linespread{1}\selectfont
		\vspace{0.1em}
		\noindent{\color{lightgray}\rule{\linewidth}{0.3pt}}
		\vspace{0.1em}
		
		\noindent
		版本历史
		
		\vspace{0.1em}
		\texttt{\makebox[3.0em][l]{v2.5}2026-10-20\hspace{1.2em}第 12 节新增“为何刻意回避哈勃张力”。}
				
		\vspace{0.1em}
		\texttt{\makebox[3.0em][l]{v2.4}2026-10-16\hspace{1.2em}第 8 节新增“常数 $P$ 的弗里德曼层面解读”。}
		
		\vspace{0.1em}
		\texttt{\makebox[3.0em][l]{v2.3}2026-09-14\hspace{1.2em}第 12 节新增卡西米尔效应与真空能讨论。}
		
		\vspace{0.1em}
		\texttt{\makebox[3.0em][l]{v2.2}2026-09-07\hspace{1.2em}第 5 节新增图 2 及相关讨论。}
		
		\vspace{0.1em}
		\texttt{\makebox[3.0em][l]{v2.1}2026-09-01\hspace{1.2em}新增附录“倘若该猜想成立”。}
		
		\vspace{0.1em}
		\texttt{\makebox[3.0em][l]{v2.0}2026-08-24\hspace{1.2em}第 11 节新增 Mission-M 研究计划。}
		
		\vspace{0.1em}
		\texttt{\makebox[3.0em][l]{v1.9}2026-08-18\hspace{1.2em}第 12 节新增“月球激光测距并不能证明 $G$ 严格不变”。}

		\vspace{0.1em}
		\texttt{\makebox[3.0em][l]{v1.8}2026-08-08\hspace{1.2em}全文增补若干参考文献与注释说明。}
		
		\vspace{0.1em}
		\texttt{\makebox[3.0em][l]{v1.7}2026-08-05\hspace{1.2em}第 8 节新增爱因斯坦莱顿演讲中关于真空的物质性的论述。}
		
		\vspace{0.1em}
		\texttt{\makebox[3.0em][l]{v1.6}2026-07-30\hspace{1.2em}第 12 节新增对现代宇宙学的方法论反思。}
		
		\vspace{0.1em}
		\texttt{\makebox[3.0em][l]{v1.5}2026-07-29\hspace{1.2em}第 7 节新增史瓦西半径讨论。}
				
		\vspace{0.1em}
		\texttt{\makebox[3.0em][l]{v1.4}2026-07-27\hspace{1.2em}标题改为更直接的表述，新增 DESI DR2 讨论。}
		
		\vspace{0.1em}
		\texttt{\makebox[3.0em][l]{v1.3}2026-07-24\hspace{1.2em}摘要修订与优化。}
		
		\vspace{0.1em}
		\texttt{\makebox[3.0em][l]{v1.2}2026-06-26\hspace{1.2em}新增关于重建弗里德曼方程的说明。}
		
		\vspace{0.1em}
		\texttt{\makebox[3.0em][l]{v1.1}2026-06-24\hspace{1.2em}首个公开版本。}
		
		\vspace{0.1em}
		\texttt{\makebox[3.0em][l]{v1.0}2026-06-13\hspace{1.2em}初稿。}
				
		\vskip 0.2in
		\noindent
		引用模板
				
		\smallskip
		\hspace{-0.8em} \begin{minipage}{\linewidth}
			\begin{verbatim}
				@misc{XianLiang:2026_growing_ct,
					author       = "XianLiang.Ma",
					title        = "{XianLiang Conjecture: Cosmic Redshift is Just 
					                 the Light Speed Expansion! G(t)c(t)^2 = P}",
					year         = "2026",
					month        = "6",
					url          = "https://github.com/ma-xian-liang/2026-growing-ct",
					howpublished = "{GitHub repository}"
				}
			\end{verbatim}
		\end{minipage}
		
		
		\vskip 0.1in
		\noindent
		存档于
		
		\smallskip
		\url{https://orcid.org/0009-0003-6406-8514}
		
		\smallskip
		\url{https://doi.org/10.5281/zenodo.23279668} 
		
		\smallskip
		\url{https://github.com/ma-xian-liang/2026-growing-ct}
		
		\vskip 0.2in
		\noindent
		文档指纹（SHA-256，IPR.TSA.CN 可信时间戳）
		
		\smallskip
		\texttt{463F919377637F84C125D355DD74471CCF8A301FA061D6E1DFAB88445BF44DE2}		
	\end{footnotesize}
		
\end{document}
